% Multiplication d'un vecteur par un nombre réel — Mathématiques 3ème — Cours signé OnBuch+
% Programme officiel MINESEC. Compiler avec : tectonic N07-multiplication-vecteur-par-nombre-reel.tex
\newcommand{\DOCMATIERE}{Mathématiques}
\newcommand{\DOCNIVEAU}{3ème}
\newcommand{\DOCMODULE}{Mathématiques – classe de 3ème}
\newcommand{\DOCLECON}{N07}
\newcommand{\DOCTITRE}{Multiplication d'un vecteur par un nombre réel}
% OnBuch+ — préambule « cours signés » ENRICHI (compilé avec tectonic / XeLaTeX)
% Système visuel « Pop » d'OnBuch+ : fond crème (jamais blanc), bordures encre
% épaisses, ombres « dures » décalées, titres Archivo Black en capitales.
% Version MATHÉMATIQUES : graphes (pgfplots/tikz), boîtes pédagogiques
% (définition, propriété, méthode, attention, le savais-tu, exercice résolu,
% exercice, TP, bilan), page de garde et signature OnBuch+ ancrée en bas.
%
% Macros attendues dans main.tex AVANT \input{preamble} :
%   \DOCMATIERE  \DOCNIVEAU  \DOCMODULE  \DOCLECON (numéro)  \DOCTITRE
\documentclass[11pt]{article}

\usepackage{fontspec}
\usepackage[french]{babel}
\usepackage[a4paper,margin=2cm,top=3.4cm,bottom=2.8cm,headheight=1.4cm,footskip=1.3cm]{geometry}
\usepackage{amsmath,amssymb,amsfonts}
\usepackage{mathtools} % \xmapsto, \coloneqq... souvent utilisés par les modèles
\usepackage{cancel} % simplifications barrées
% \abs{z} (module, valeur absolue) et \norm{u} (norme), utilisés par les modèles
\providecommand{\abs}{}\let\abs\relax\DeclarePairedDelimiter{\abs}{\lvert}{\rvert}
\providecommand{\norm}{}\let\norm\relax\DeclarePairedDelimiter{\norm}{\lVert}{\rVert}
\usepackage[table]{xcolor} % [table] : fournit aussi \rowcolors (alternance de lignes)
\usepackage{pagecolor}
\usepackage{tikz}
\usetikzlibrary{arrows.meta,positioning,calc,shapes.geometric,shapes.misc,shapes.arrows,decorations.pathmorphing,decorations.pathreplacing,decorations.markings,patterns,shadows,mindmap,trees,fit,backgrounds,matrix,intersections,angles,quotes}
\usepackage{pgfplots}
\usepackage{pgf-pie} % diagrammes circulaires \pie (statistiques)
\pgfplotsset{compat=1.18}
\usepgfplotslibrary{fillbetween,groupplots} % aires entre courbes (calcul intégral)
\usepackage{enumitem}
\usepackage{fancyhdr}
% [most] charge toutes les bibliothèques tcolorbox (listings, minted, raster,
% documentation...) et gonfle inutilement la mémoire TeX sur un document long
% (~300 boîtes) : on ne charge que ce qui est réellement utilisé.
\usepackage[skins,breakable]{tcolorbox}
\usepackage{graphicx}
\usepackage{colortbl}
\usepackage{booktabs}
\usepackage{tabularx}
\usepackage{diagbox} % case d'en-tête diagonale des tableaux à double entrée
% Colonnes extensibles centrée / gauche / droite (C, L, R), souvent utilisées par les modèles
\newcolumntype{C}{>{\centering\arraybackslash}X}
\newcolumntype{L}{>{\raggedright\arraybackslash}X}
\newcolumntype{R}{>{\raggedleft\arraybackslash}X}
\usepackage{array}
\usepackage{multicol}
\usepackage{siunitx}
% parse-numbers=false : accepte aussi \SI{2,5 \times 10^{-3}}{...}
\sisetup{locale=FR,output-decimal-marker={,},group-minimum-digits=4,parse-numbers=false}
\DeclareSIUnit\ohm{\ensuremath{\symup{\Omega}}} % U+2126 absent de la police
\usepackage{qrcode}
% \degree : commande fréquemment utilisée par les modèles pour un angle ou une
% température en dehors de \SI{}{} (non fournie par les paquets chargés).
\providecommand{\degree}{^{\circ}}
% \qed (fin de démonstration), utilisé par les modèles sans amsthm
\providecommand{\qed}{\hfill$\square$}
% \barre : double barre des valeurs interdites dans un tableau de variations (tkz-tab)
\providecommand{\barre}{\ensuremath{\|}}
% environnement proof (amsthm non chargé) utilisé par les modèles
\ifdefined\proof\else\newenvironment{proof}[1][Démonstration]{\par\noindent\textit{#1.}\ }{\qed\par}\fi
\usepackage{lastpage}
\usepackage{hyperref}
\hypersetup{hidelinks}

% ============================================================
% Palette « Pop » OnBuch+ — mêmes hex que la classe P (plus_ui.dart)
% ============================================================
\definecolor{poporange}{HTML}{F59321}
\definecolor{popdark}{HTML}{E07A0C}
\definecolor{popcream}{HTML}{FFF6E8}
\definecolor{popink}{HTML}{1C1714}
\definecolor{poppurple}{HTML}{7A5AE0}
\definecolor{popgreen}{HTML}{2E9E6A}
\definecolor{poppink}{HTML}{E0457B}
\definecolor{popblue}{HTML}{2F7DE1}
\definecolor{popgold}{HTML}{C98A12}
\definecolor{popmuted}{HTML}{8A7F76}
\definecolor{popline}{HTML}{EADFD2}
\definecolor{popgray}{HTML}{8A7F76}
\colorlet{popgrey}{popgray}
% Alias tolérants (le modèle nomme parfois les couleurs différemment)
\colorlet{popwhite}{white}
\colorlet{popblack}{popink}
\colorlet{popred}{poppink}
\colorlet{popyellow}{popgold}
% Teintes claires pour fonds de tableaux / zones
\colorlet{popblueL}{popblue!10!white}
\colorlet{poporangeL}{poporange!14!white}
\colorlet{popgreenL}{popgreen!12!white}
\colorlet{poppurpleL}{poppurple!12!white}
\colorlet{poppinkL}{poppink!12!white}
\colorlet{popgoldL}{popgold!14!white}

\pagecolor{popcream}

% ============================================================
% Typographie
% ============================================================
\defaultfontfeatures{Ligatures=TeX}
\setmainfont{PlusJakartaSans-Medium.ttf}[Path=../../../fonts/,BoldFont=PlusJakartaSans-Bold.ttf]
% Maths en sans-serif (Fira Math) avec chiffres et lettres droites de Plus
% Jakarta, pour que les formules se fondent dans le texte.
\usepackage[math-style=ISO,bold-style=ISO]{unicode-math}
\setmathfont{FiraMath-Regular.otf}[Path=../../../fonts/]
\setmathfont{PlusJakartaSans-Medium.ttf}[Path=../../../fonts/,range={up/num,up/{Latin,latin},\mathup}]
\usepackage{icomma} % virgule décimale sans espace en mode math
% Caractères mathématiques Unicode absents des polices de texte (Plus Jakarta,
% Archivo Black) : rendus par leur équivalent mathématique, en texte comme en
% formule — sinon ils s'affichent en carré vide (ex. « Arithmétique dans ℤ »).
\usepackage{newunicodechar}
\newunicodechar{ℝ}{\ensuremath{\mathbb{R}}}
\newunicodechar{ℤ}{\ensuremath{\mathbb{Z}}}
\newunicodechar{ℂ}{\ensuremath{\mathbb{C}}}
\newunicodechar{ℕ}{\ensuremath{\mathbb{N}}}
\newunicodechar{ℚ}{\ensuremath{\mathbb{Q}}}
\newunicodechar{𝔻}{\ensuremath{\mathbb{D}}}
\newunicodechar{α}{\ensuremath{\alpha}}
\newunicodechar{β}{\ensuremath{\beta}}
\newunicodechar{γ}{\ensuremath{\gamma}}
\newunicodechar{δ}{\ensuremath{\delta}}
\newunicodechar{ε}{\ensuremath{\varepsilon}}
\newunicodechar{θ}{\ensuremath{\theta}}
\newunicodechar{λ}{\ensuremath{\lambda}}
\newunicodechar{μ}{\ensuremath{\mu}}
\newunicodechar{σ}{\ensuremath{\sigma}}
\newunicodechar{τ}{\ensuremath{\tau}}
\newunicodechar{φ}{\ensuremath{\varphi}}
\newunicodechar{ψ}{\ensuremath{\psi}}
\newunicodechar{ω}{\ensuremath{\omega}}
\newunicodechar{Σ}{\ensuremath{\Sigma}}
% majuscules produites par \MakeUppercase dans les titres (α-aminés -> Α)
\newunicodechar{Α}{\ensuremath{\alpha}}
\newunicodechar{Β}{\ensuremath{\beta}}
\newunicodechar{Γ}{\ensuremath{\Gamma}}
\newunicodechar{Δ}{\ensuremath{\Delta}}
\newunicodechar{Ε}{\ensuremath{\varepsilon}}
\newunicodechar{Θ}{\ensuremath{\Theta}}
\newunicodechar{Λ}{\ensuremath{\Lambda}}
\newunicodechar{Μ}{\ensuremath{\mu}}
\newunicodechar{Τ}{\ensuremath{\tau}}
\newunicodechar{Φ}{\ensuremath{\Phi}}
\newunicodechar{Ψ}{\ensuremath{\Psi}}
\newunicodechar{Ω}{\ensuremath{\Omega}}
\newunicodechar{↦}{\ensuremath{\mapsto}}
\newunicodechar{∈}{\ensuremath{\in}}
\newunicodechar{∉}{\ensuremath{\notin}}
\newunicodechar{∧}{\ensuremath{\wedge}}
\newunicodechar{⇔}{\ensuremath{\Leftrightarrow}}
\newunicodechar{⟺}{\ensuremath{\Longleftrightarrow}}
\newunicodechar{⇒}{\ensuremath{\Rightarrow}}
\newunicodechar{⟹}{\ensuremath{\Longrightarrow}}
\newunicodechar{∘}{\ensuremath{\circ}}
\newunicodechar{∪}{\ensuremath{\cup}}
\newunicodechar{∩}{\ensuremath{\cap}}
\newunicodechar{≡}{\ensuremath{\equiv}}
\newunicodechar{⊂}{\ensuremath{\subset}}
\newunicodechar{⊥}{\ensuremath{\perp}}
\newunicodechar{∃}{\ensuremath{\exists}}
\newunicodechar{∀}{\ensuremath{\forall}}
\newunicodechar{∛}{\ensuremath{\sqrt[3]{\,}}}
\newunicodechar{∜}{\ensuremath{\sqrt[4]{\,}}}
\newunicodechar{⃗}{\ensuremath{{}^{\to}}}
\newunicodechar{ⁿ}{\ifmmode^{n}\else\textsuperscript{\ensuremath{n}}\fi}
\newunicodechar{ˣ}{\ifmmode^{x}\else\textsuperscript{\ensuremath{x}}\fi}
\newunicodechar{ᵇ}{\ifmmode^{b}\else\textsuperscript{\ensuremath{b}}\fi}
\newunicodechar{ʳ}{\ifmmode^{r}\else\textsuperscript{\ensuremath{r}}\fi}
\newunicodechar{ᵘ}{\ifmmode^{u}\else\textsuperscript{\ensuremath{u}}\fi}
\newunicodechar{ʸ}{\ifmmode^{y}\else\textsuperscript{\ensuremath{y}}\fi}
\newunicodechar{ᵃ}{\ifmmode^{a}\else\textsuperscript{\ensuremath{a}}\fi}
\newunicodechar{ᵉ}{\ifmmode^{e}\else\textsuperscript{\ensuremath{e}}\fi}
\newunicodechar{ᵏ}{\ifmmode^{k}\else\textsuperscript{\ensuremath{k}}\fi}
\newunicodechar{ᵖ}{\ifmmode^{p}\else\textsuperscript{\ensuremath{p}}\fi}
\newunicodechar{ᴺ}{\ifmmode^{N}\else\textsuperscript{\ensuremath{N}}\fi}
\newunicodechar{ᶜ}{\ifmmode^{c}\else\textsuperscript{\ensuremath{c}}\fi}
\newunicodechar{ʰ}{\ifmmode^{h}\else\textsuperscript{\ensuremath{h}}\fi}
\newunicodechar{ᵠ}{\ifmmode^{q}\else\textsuperscript{\ensuremath{q}}\fi}
\newunicodechar{ₙ}{\ifmmode_{n}\else\textsubscript{\ensuremath{n}}\fi}
\newunicodechar{ₐ}{\ifmmode_{a}\else\textsubscript{\ensuremath{a}}\fi}
\newunicodechar{ᵢ}{\ifmmode_{i}\else\textsubscript{\ensuremath{i}}\fi}
\newunicodechar{ᵣ}{\ifmmode_{r}\else\textsubscript{\ensuremath{r}}\fi}
\newunicodechar{ₖ}{\ifmmode_{k}\else\textsubscript{\ensuremath{k}}\fi}
\newunicodechar{ᵦ}{\ifmmode_{\beta}\else\textsubscript{\ensuremath{\beta}}\fi}
\newunicodechar{ₓ}{\ifmmode_{x}\else\textsubscript{\ensuremath{x}}\fi}
\newfontfamily\popdisplay{ArchivoBlack-Regular.ttf}[Path=../../../fonts/]
\newfontfamily\poplogo{SpaceGrotesk-Bold.ttf}[Path=../../../fonts/]
\newfontfamily\popsemi{PlusJakartaSans-SemiBold.ttf}[Path=../../../fonts/]
\newfontfamily\popxbold{PlusJakartaSans-ExtraBold.ttf}[Path=../../../fonts/]
\newfontfamily\popxboldls{PlusJakartaSans-ExtraBold.ttf}[Path=../../../fonts/,LetterSpace=3.0]

\setlist[enumerate]{leftmargin=1.7em,itemsep=5pt,topsep=4pt}
\setlist[itemize]{leftmargin=1.7em,itemsep=4pt,topsep=4pt,label={\color{poporange}\textbullet}}
\setlength{\parindent}{0pt}
\setlength{\parskip}{5pt}
\renewcommand{\arraystretch}{1.25}

% pgfplots : style Pop par défaut
\pgfplotsset{
  every axis/.append style={axis line style={popink,line width=1pt},tick style={popink},
    grid style={popline,line width=0.5pt},label style={font=\small},tick label style={font=\footnotesize}},
  popcurve/.style={poporange,line width=1.8pt,no markers,smooth},
}
% Même style disponible dans un \draw TikZ simple (hors pgfplots)
\tikzset{popcurve/.style={poporange,line width=1.8pt}}

% ============================================================
% Pill / squiggle
% ============================================================
\newcommand{\poppill}[3]{%
  \tikz[baseline=-0.55ex]\node[draw=popink,line width=1.2pt,rounded corners=2.6pt,%
    fill=#2,inner xsep=6.5pt,inner ysep=3pt]{%
    \popxboldls\fontsize{7}{7}\selectfont\color{#3}\MakeUppercase{#1}};%
}
\newcommand{\popsquiggle}[1]{%
  \tikz[baseline=-0.4ex]{
    \draw[#1,line width=2.6pt,line cap=round]
      plot [smooth] coordinates {(0,0.09) (0.34,0.01) (0.68,0.21) (1.02,0.01) (1.36,0.10)};
  }%
}
% Mot-clé surligné (marqueur orange sous le texte)
\newcommand{\cle}[1]{{\popxbold\color{popdark}#1}}

% ============================================================
% En-tête / pied
% ============================================================
\pagestyle{fancy}
\fancyhf{}
\renewcommand{\headrulewidth}{0pt}
\renewcommand{\footrulewidth}{0pt}
\lhead{\raisebox{-0.15ex}{\poplogo\fontsize{13}{13}\selectfont\color{popink}OnBuch\color{poporange}+}}
\rhead{\poppill{\DOCMATIERE\ \textbullet\ \DOCNIVEAU\ \textbullet\ Leçon \DOCLECON}{white}{popink}}
\lfoot{\popxbold\fontsize{7.6}{7.6}\selectfont\color{popmuted}\MakeUppercase{Cours signé OnBuch+}\ \textbullet\ {\popsemi\DOCTITRE}}
\rfoot{\tikz[baseline=-0.6ex]\node[rounded corners=8.5pt,fill=popink,minimum height=17pt,minimum width=17pt,inner xsep=5pt,inner sep=0]{\popxbold\fontsize{8}{8}\selectfont\color{popcream}\thepage\,/\,\pageref*{LastPage}};}

% ============================================================
% Titres
% ============================================================
\newcommand{\chaptitre}[2]{%
  \begin{tcolorbox}[enhanced,colback=poporange,colframe=popink,boxrule=2.6pt,arc=11pt,
      drop shadow=popink,left=15pt,right=15pt,top=12pt,bottom=11pt,before skip=3pt,after skip=18pt]
    \poppill{#2}{popink}{popcream}\par\vspace{9pt}
    {\raggedright\hyphenpenalty=10000\exhyphenpenalty=10000\popdisplay\fontsize{22}{24}\selectfont\color{popcream}\MakeUppercase{#1}\par}
  \end{tcolorbox}
}
% Section numérotée : pastille orange avec le numéro + titre Archivo
\newcounter{popsec}
\newcommand{\coursec}[1]{%
  \stepcounter{popsec}\setcounter{popsub}{0}%
  \par\vspace{16pt}\noindent
  \begin{minipage}{\linewidth}
  \tikz[baseline=-0.7ex]\node[draw=popink,line width=1.6pt,fill=poporange,rounded corners=4pt,minimum size=22pt,inner sep=0]{\popdisplay\fontsize{12}{12}\selectfont\color{popcream}\thepopsec};%
  \hspace{8pt}{\popdisplay\fontsize{15}{17}\selectfont\color{popink}\MakeUppercase{#1}}\par
  \vspace{4pt}\noindent{\color{poporange}\rule{36pt}{3.6pt}}
  \end{minipage}\par\nopagebreak\vspace{8pt}
}
\newcounter{popsub}
\newcommand{\courssub}[1]{%
  \stepcounter{popsub}%
  \par\vspace{9pt}\noindent{\popxbold\fontsize{11.5}{13}\selectfont\color{popink}{\color{poporange}\thepopsec.\thepopsub}\hspace{6pt}#1}\par\nopagebreak\vspace{4pt}
}

% ============================================================
% Boîtes pédagogiques — même recette : fond blanc, bordure épaisse colorée,
% ombre dure encre, badge plein. Titre optionnel [#1].
% ============================================================
\tcbset{popbase/.style={breakable,enhanced,colback=white,boxrule=2.3pt,arc=9pt,
  shadow={3pt}{-3pt}{0pt}{popink},left=14pt,right=14pt,top=11pt,bottom=11pt,before skip=11pt,after skip=15pt,
  fontupper=\popsemi\fontsize{10.6}{14.5}\selectfont\color{popink},
  fontlower=\popsemi\fontsize{10.6}{14.5}\selectfont\color{popink}}}
% Coche dessinée (le glyphe ✓ n'existe pas dans les polices du document)
\newcommand{\popcheck}[1]{\tikz[baseline=-0.5ex]\draw[#1,line width=1.6pt,line cap=round,line join=round](-0.13,0.0)--(-0.04,-0.09)--(0.13,0.1);}
\newcommand{\popboxhead}[3]{\poppill{#1}{#2}{white}\ifx\relax#3\relax\else\hspace{6pt}{\popxbold\fontsize{10.6}{12}\selectfont\color{popink}#3}\fi\par\vspace{7pt}}

\newtcolorbox{aretenir}[1][]{popbase,colframe=popgreen,before upper={\popboxhead{À retenir}{popgreen}{#1}}}
\newtcolorbox{exemplebox}[1][]{popbase,colframe=poporange,before upper={\popboxhead{Exemple}{poporange}{#1}}}
\newtcolorbox{definition}[1][]{popbase,colframe=popblue,before upper={\popboxhead{Définition}{popblue}{#1}}}
\newtcolorbox{propriete}[1][]{popbase,colframe=poppurple,before upper={\popboxhead{Propriété}{poppurple}{#1}}}
\newtcolorbox{methode}[1][]{popbase,colframe=popgold,before upper={\popboxhead{Méthode}{popgold}{#1}}}
\newtcolorbox{attention}[1][]{popbase,colframe=poppink,before upper={\popboxhead{Attention}{poppink}{#1}}}
\newtcolorbox{experience}[1][]{popbase,colframe=popblue,colback=popblueL,before upper={\popboxhead{Expérience}{popblue}{#1}}}
% « Le savais-tu ? » : carte violette pleine (contexte, culture, Cameroun)
\newtcolorbox{savaistu}[1][]{popbase,colback=poppurple,colframe=popink,
  fontupper=\popsemi\fontsize{10.4}{14.2}\selectfont\color{white},
  before upper={\poppill{Le savais-tu\,?}{white}{poppurple}\ifx\relax#1\relax\else\hspace{6pt}{\popxbold\color{white}#1}\fi\par\vspace{7pt}}}
% Exercice résolu : énoncé en haut, \tcblower puis la solution sur fond crème
\newcounter{popexo}\newcounter{popexr}
\newtcolorbox{exoresolu}[1][]{popbase,colframe=poporange,colbacklower=popcream,
  segmentation style={popink,line width=1.2pt,dashed},
  before upper={\stepcounter{popexr}\popboxhead{Exercice résolu \thepopexr}{poporange}{#1}},
  before lower={\poppill{Solution}{popink}{popcream}\par\vspace{6pt}}}
% Exercice (non résolu en place) — la correction est dans la partie Corrigés
\newtcolorbox{exercice}[1][]{popbase,colframe=popink,
  before upper={\stepcounter{popexo}\popboxhead{Exercice \thepopexo}{popink}{#1}}}
% Corrigé
\newtcolorbox{corrige}[1][]{popbase,colframe=popgreen,colback=popgreenL,
  before upper={\popboxhead{Corrigé}{popgreen}{#1}}}
% Objectifs / prérequis / bilan
\newtcolorbox{objectifs}{popbase,colframe=popink,colback=poporangeL,
  before upper={\popboxhead{Objectifs de la leçon}{popink}{}}}
\newtcolorbox{prerequis}{popbase,colframe=popink,colback=popblueL,
  before upper={\popboxhead{Prérequis}{popblue}{}}}
\newtcolorbox{bilan}{popbase,colframe=popink,colback=white,boxrule=2.8pt,
  before upper={\popboxhead{Fiche bilan}{poporange}{L'essentiel à connaître pour l'examen}}}
% Figure encadrée avec légende
\newtcolorbox{popfigure}[1][]{enhanced,breakable=false,colback=white,colframe=popline,boxrule=1.4pt,arc=8pt,
  left=10pt,right=10pt,top=10pt,bottom=8pt,before skip=10pt,after skip=12pt,halign=center,
  lower separated=false,halign lower=center,
  fontlower=\popsemi\fontsize{9}{11.5}\selectfont\color{popmuted},
  #1}
\newcommand{\legende}[1]{\par\vspace{6pt}{\popsemi\fontsize{9}{11.5}\selectfont\color{popmuted}#1}}

% ============================================================
% Signature OnBuch+
% ============================================================
\newcommand{\poplogoinline}[1]{{\poplogo\fontsize{#1}{#1}\selectfont\color{popink}OnBuch\color{poporange}+}}
\newcommand{\signatureonbuch}{%
  \begin{tcolorbox}[enhanced,colback=popink,colframe=popink,boxrule=2.2pt,arc=10pt,
      drop shadow=poporange,left=14pt,right=14pt,top=10pt,bottom=10pt,before skip=0pt,after skip=0pt]
    \tikz[baseline=-0.7ex]\node[circle,fill=popgreen,draw=popcream,line width=1.3pt,minimum size=22pt,inner sep=0]{\popcheck{white}};%
    \hspace{9pt}\begin{minipage}[c]{0.62\linewidth}
      {\popxbold\fontsize{12}{13}\selectfont\color{popcream}Signé\ \textbullet\ {\poplogo OnBuch\color{poporange}+}}\par\vspace{2pt}
      {\raggedright\popsemi\fontsize{8}{10.5}\selectfont\color{popline}\MakeUppercase{Équipe pédagogique OnBuch+}\\\MakeUppercase{Conforme au programme officiel MINESEC}\par}
    \end{minipage}\hfill
    {\poplogo\fontsize{22}{22}\selectfont\color{popcream}OnBuch\color{poporange}+}
  \end{tcolorbox}%
}

% ============================================================
% Page de garde
% #1 titre · #2 matière · #3 niveau · #4 module · #5 n° leçon · #6 durée ·
% #7 description / objectifs (texte ou itemize)
% ============================================================
\newcommand{\pagegarde}[7]{%
  \thispagestyle{empty}
  \newgeometry{a4paper,margin=1.7cm,top=1.6cm,bottom=1.4cm}%
  \begin{tikzpicture}[remember picture,overlay]
    \fill[poporange!18] ([xshift=-3.2cm,yshift=-3.6cm]current page.north east) circle (4.6cm);
    \fill[poppurple!14] ([xshift=2.4cm,yshift=7.6cm]current page.south west) circle (3.8cm);
  \end{tikzpicture}%
  {\poplogo\fontsize{30}{30}\selectfont\color{popink}OnBuch\color{poporange}+}\hfill
  \raisebox{0.6ex}{\poppill{Cours signé \textbullet\ Premium}{popink}{popcream}}\par
  \vspace{1pt}\popsquiggle{poppurple}\par
  \vspace{8pt}
  \begin{tcolorbox}[enhanced,colback=poporange,colframe=popink,boxrule=2.8pt,arc=13pt,
      drop shadow=popink,left=18pt,right=18pt,top=12pt,bottom=13pt,after skip=10pt]
    \poppill{#2 \textbullet\ #3}{popink}{popcream}\hspace{5pt}\poppill{#4}{white}{popink}\par\vspace{8pt}
    {\popdisplay\fontsize{14}{14}\selectfont\color{popink}LEÇON #5}\par\vspace{4pt}
    {\raggedright\hyphenpenalty=10000\exhyphenpenalty=10000\popdisplay\fontsize{25}{27}\selectfont\color{popcream}\MakeUppercase{#1}\par}
    \vspace{8pt}
    {\popxbold\fontsize{9.5}{9.5}\selectfont\color{popink}\MakeUppercase{Durée conseillée : #6}}
  \end{tcolorbox}
  \begin{tcolorbox}[enhanced,colback=white,colframe=popink,boxrule=2.2pt,arc=10pt,
      drop shadow=popink,left=16pt,right=16pt,top=10pt,bottom=10pt,after skip=8pt]
    \poppill{Dans cette leçon}{poppurple}{white}\par\vspace{5pt}
    \popsemi\fontsize{9.3}{12.6}\selectfont\color{popink}#7
  \end{tcolorbox}
  \noindent\begin{minipage}[t]{0.487\linewidth}
    \begin{tcolorbox}[enhanced,colback=popgreen,colframe=popink,boxrule=2.2pt,arc=10pt,
        drop shadow=popink,left=11pt,right=11pt,top=10pt,bottom=10pt,height=3.1cm,valign=center]
      \tcbox[colback=white,colframe=popink,boxrule=1.3pt,arc=5pt,boxsep=0pt,nobeforeafter,
        left=3pt,right=3pt,top=3pt,bottom=3pt,valign=center]{\qrcode[height=1.9cm]{https://play.google.com/store/apps/details?id=com.luvvix.onbuch&hl=fr}}%
      \hspace{8pt}\begin{minipage}[c]{0.5\linewidth}
        {\popdisplay\fontsize{11}{12.5}\selectfont\color{white}\MakeUppercase{Télécharge OnBuch}}\par\vspace{4pt}
        {\raggedright\popsemi\fontsize{8.4}{11}\selectfont\color{white}Gratuit \textbullet\ Google Play\\Tuteur IA, annales, concours, résultats\par}
      \end{minipage}
    \end{tcolorbox}
  \end{minipage}\hfill
  \begin{minipage}[t]{0.487\linewidth}
    \begin{tcolorbox}[enhanced,colback=poppurple,colframe=popink,boxrule=2.2pt,arc=10pt,
        drop shadow=popink,left=11pt,right=11pt,top=10pt,bottom=10pt,height=3.1cm,valign=center]
      \tikz[baseline=-0.7ex]\node[circle,fill=white,draw=popink,line width=1.3pt,minimum size=1.6cm,inner sep=0]{\popdisplay\fontsize{24}{24}\selectfont\color{poppurple}+};%
      \hspace{8pt}\begin{minipage}[c]{0.6\linewidth}
        {\popdisplay\fontsize{11}{12.5}\selectfont\color{white}\MakeUppercase{Passe à OnBuch+}}\par\vspace{4pt}
        {\raggedright\popsemi\fontsize{8.4}{11}\selectfont\color{white}Tous les cours signés, Tuteur IA illimité, fiches PDF — onglet OnBuch+ de l'app\par}
      \end{minipage}
    \end{tcolorbox}
  \end{minipage}
  \vspace{8pt}
  \popcontenu
  \vfill
  \signatureonbuch
  \restoregeometry
  \newpage
}

% Rangée « Ce cours contient » (page de garde)
\newcommand{\poptile}[3]{% #1 couleur #2 gros texte #3 légende
  \begin{tcolorbox}[enhanced,colback=#1,colframe=popink,boxrule=2pt,arc=9pt,shadow={2.5pt}{-2.5pt}{0pt}{popink},
      left=6pt,right=6pt,top=9pt,bottom=9pt,halign=center,valign=center,height=1.75cm,width=\linewidth,nobeforeafter]
    {\popdisplay\fontsize{12.5}{13}\selectfont\color{white}\MakeUppercase{#2}}\par\vspace{3pt}
    {\centering\popsemi\fontsize{7.8}{9.5}\selectfont\color{white}#3\par}
  \end{tcolorbox}}
\newcommand{\popcontenu}{%
  \par\noindent{\popxboldls\fontsize{7.5}{7.5}\selectfont\color{popmuted}CE COURS SIGNÉ CONTIENT}\par\vspace{7pt}
  \noindent\begin{minipage}[t]{0.235\linewidth}\poptile{popblue}{Cours}{complet, illustré, pas à pas}\end{minipage}\hfill
  \begin{minipage}[t]{0.235\linewidth}\poptile{poporange}{Méthodes}{et exercices résolus}\end{minipage}\hfill
  \begin{minipage}[t]{0.235\linewidth}\poptile{poppink}{Exercices}{type examen + corrigés}\end{minipage}\hfill
  \begin{minipage}[t]{0.235\linewidth}\poptile{popgreen}{Fiche bilan}{l'essentiel à retenir}\end{minipage}\par}

% Sommaire
\newtcolorbox{sommaire}{enhanced,colback=white,colframe=popink,boxrule=2.2pt,arc=10pt,
  drop shadow=popink,left=18pt,right=18pt,top=13pt,bottom=13pt,before skip=6pt,after skip=20pt,
  before upper={\poppill{Sommaire}{poppurple}{white}\par\vspace{9pt}\popsemi\fontsize{10.6}{14.5}\selectfont\color{popink}}}

% Signature de fin de cours — poussée tout en bas de la dernière page
\newcommand{\findecours}{%
  \par\vfill\nopagebreak
  \begin{center}\popsquiggle{poporange}\end{center}\vspace{4pt}
  \signatureonbuch
}
\AtBeginDocument{\def\llbracket{\mathopen{[\mkern-3.5mu[}}\def\rrbracket{\mathclose{]\mkern-3.5mu]}}} % intervalles d'entiers (glyphe ⟦ absent de la police)
% Glyphes absents de Fira Math : redéfinis à partir de glyphes disponibles
\AtBeginDocument{%
  \def\setminus{\mathbin{\backslash}}\let\smallsetminus\setminus
  \def\vdots{\mathord{\vbox{\baselineskip3.2pt\lineskiplimit0pt\kern4pt\hbox{.}\hbox{.}\hbox{.}}}}%
  \def\ddots{\mathinner{\mkern1mu\raise6.5pt\hbox{.}\mkern2mu\raise3.6pt\hbox{.}\mkern2mu\raise0.7pt\hbox{.}\mkern1mu}}%
  \def\triangle{\mathord{\scalebox{0.9}{$\symup{\Delta}$}}}%
  \def\nabla{\mathord{\raisebox{\depth}{\scalebox{1}[-1]{$\symup{\Delta}$}}}}%
  \def\checkmark{\popcheck{popdark}}%
  \def\star{\mathbin{\ast}}\def\bigstar{\mathord{\ast}}%
  \def\diamond{\mathbin{\scalebox{0.6}{$\Diamond$}}}%
  \def\bigcup{\mathop{\mathchoice{\raisebox{-0.3em}{\scalebox{1.7}{$\cup$}}}{\scalebox{1.25}{$\cup$}}{\cup}{\cup}}\displaylimits}%
  \def\bigcap{\mathop{\mathchoice{\raisebox{-0.3em}{\scalebox{1.7}{$\cap$}}}{\scalebox{1.25}{$\cap$}}{\cap}{\cap}}\displaylimits}%
  \def\lesseqqgtr{\mathrel{\lessgtr}}\def\bigoplus{\mathop{\oplus}\displaylimits}%
}
\providecommand{\ssi}{si et seulement si} % abréviation fréquente
\AtBeginDocument{\providecommand{\wideparen}{\overparen}} % arc de cercle

\begin{document}

\pagegarde{Multiplication d'un vecteur par un nombre réel}{Mathématiques}{3ème}{Mathématiques – classe de 3ème}{N07}{2 séances (fiche de progression harmonisée)}{Cette leçon introduit le produit d'un vecteur par un nombre réel et la notion fondamentale de vecteurs colinéaires. Elle fournit des outils puissants pour démontrer le parallélisme de droites, l'alignement de points et caractériser les vecteurs directeurs d'une droite.
\vspace{4pt}
\begin{multicols}{2}\small
\begin{itemize}
\item À la fin de cette leçon, je sais construire le vecteur k·u pour un réel k et un vecteur u donnés
\item À la fin de cette leçon, je sais définir le produit d'un vecteur par un nombre réel et interpréter géométriement le signe de k
\item À la fin de cette leçon, je sais appliquer les règles de calcul sur les vecteurs (distributivité, associativité)
\item À la fin de cette leçon, je sais reconnaître et démontrer que deux vecteurs sont colinéaires
\item À la fin de cette leçon, je sais utiliser la colinéarité pour démontrer que des droites sont parallèles
\item À la fin de cette leçon, je sais utiliser la colinéarité pour démontrer que trois points sont alignés
\end{itemize}\end{multicols}}

\begin{sommaire}
\begin{multicols}{2}
\begin{enumerate}
\item Produit d'un vecteur par un nombre réel
\item Règles de calcul sur les produits de vecteurs
\item Vecteurs colinéaires
\item Colinéarité, parallélisme et alignement
\item Vecteurs directeurs d'une droite
\item Applications et situations concrètes
\item Activité d'intégration
\item Méthodes et exercices résolus
\item Exercices
\item Corrigés des exercices
\item Fiche bilan
\end{enumerate}
\end{multicols}
\end{sommaire}

\begin{objectifs}
\begin{itemize}
\item À la fin de cette leçon, je sais construire le vecteur k·u pour un réel k et un vecteur u donnés
\item À la fin de cette leçon, je sais définir le produit d'un vecteur par un nombre réel et interpréter géométriement le signe de k
\item À la fin de cette leçon, je sais appliquer les règles de calcul sur les vecteurs (distributivité, associativité)
\item À la fin de cette leçon, je sais reconnaître et démontrer que deux vecteurs sont colinéaires
\item À la fin de cette leçon, je sais utiliser la colinéarité pour démontrer que des droites sont parallèles
\item À la fin de cette leçon, je sais utiliser la colinéarité pour démontrer que trois points sont alignés
\item À la fin de cette leçon, je sais déterminer et utiliser un vecteur directeur d'une droite
\item À la fin de cette leçon, je sais rédiger et justifier une démonstration complète en géométrie vectorielle
\item À la fin de cette leçon, je sais appliquer ces notions à des situations concrètes de la vie courante (déplacements, forces, agrandissements)
\end{itemize}
\end{objectifs}

\begin{prerequis}
\begin{itemize}
\item Notion de vecteur : direction, sens, norme (classe de 4ème)
\item Égalité de deux vecteurs et vecteurs opposés
\item Somme de deux vecteurs : relation de Chasles et règle du parallélogramme
\item Coordonnées d'un vecteur dans un repère
\item Notions de base : points alignés, droites parallèles, milieu d'un segment
\end{itemize}
\end{prerequis}

\newpage

\coursec{Produit d'un vecteur par un nombre réel}

Moussa conduit sa moto-taxi de Yaoundé vers Soa. Il roule \textbf{deux fois plus vite} qu'un cycliste qui va dans la même direction. Comment décrire mathématiquement le déplacement de Moussa à partir de celui du cycliste ? Et si un véhicule fait \textbf{demi-tour}, comment traduire ce changement de sens avec les vecteurs ?

Ces questions de la vie quotidienne cachent une opération fondamentale : \cle{multiplier un vecteur par un nombre réel}. Grâce à elle, nous pourrons aussi, dans les sections suivantes, reconnaître des droites parallèles et des points alignés \emph{sans avoir à les tracer entièrement} : un outil précieux pour le BEPC.

\courssub{Activité d'introduction : additionner un vecteur avec lui-même}

\begin{experience}[Que vaut $\vec{u} + \vec{u}$ ?]
Prends une feuille, une règle et un compas.
\begin{enumerate}
\item Trace un vecteur $\vec{u} = \overrightarrow{AB}$ de longueur $2$~cm, dirigé vers la droite.
\item À partir du point $B$, reporte le même vecteur : place le point $C$ tel que $\overrightarrow{BC} = \vec{u}$. Par la relation de Chasles, $\overrightarrow{AC} = \overrightarrow{AB} + \overrightarrow{BC} = \vec{u} + \vec{u}$.
\item Observe le vecteur $\overrightarrow{AC}$ et réponds :
\begin{itemize}
\item Quelle est sa \textbf{direction} par rapport à celle de $\vec{u}$ ?
\item Quel est son \textbf{sens} par rapport à celui de $\vec{u}$ ?
\item Quelle est sa \textbf{longueur} par rapport à celle de $\vec{u}$ ?
\end{itemize}
\end{enumerate}
\textbf{Constat :} le vecteur $\vec{u} + \vec{u}$ a la \emph{même direction} que $\vec{u}$, le \emph{même sens} que $\vec{u}$, et une longueur \emph{double}. Il est naturel de le noter $2\vec{u}$ :
\[ \vec{u} + \vec{u} = 2\vec{u}. \]
De même, $\vec{u} + \vec{u} + \vec{u} = 3\vec{u}$ serait trois fois plus long, toujours dans la même direction et le même sens.
\end{experience}

Cette activité montre qu'additionner un vecteur avec lui-même revient à \emph{étirer} le vecteur sans changer ni sa direction ni son sens. Mais que se passe-t-il si l'on multiplie par un nombre \emph{négatif}, ou par une \emph{fraction} comme $\dfrac{1}{2}$ ? C'est toute la richesse de la notion que nous allons définir.

\courssub{Définition du produit d'un vecteur par un nombre réel}

Revenons à Moussa. Le déplacement du cycliste est un vecteur $\vec{v}$ : il a une direction (la route Yaoundé--Soa), un sens (vers Soa) et une longueur (la distance parcourue). Moussa va deux fois plus vite dans le même sens : son déplacement est $2\vec{v}$. Un véhicule qui fait demi-tour parcourt un déplacement \emph{opposé} : c'est $-1$ fois le vecteur initial. Généralisons.

\begin{definition}[Produit d'un vecteur par un nombre réel]
Soit $k$ un nombre réel \textbf{non nul} et $\vec{u}$ un vecteur \textbf{non nul}. Le produit de $\vec{u}$ par $k$, noté $k\vec{u}$, est le vecteur caractérisé par :
\begin{itemize}
\item sa \cle{direction} : $k\vec{u}$ a la \textbf{même direction} que $\vec{u}$ ;
\item son \cle{sens} : $k\vec{u}$ a le \textbf{même sens} que $\vec{u}$ si $k > 0$, et le \textbf{sens opposé} à celui de $\vec{u}$ si $k < 0$ ;
\item sa \cle{norme} (longueur) : $\|k\vec{u}\| = |k| \times \|\vec{u}\|$.
\end{itemize}
\end{definition}

\begin{aretenir}[La règle en trois questions]
Pour construire ou décrire $k\vec{u}$, pose-toi trois questions :
\begin{enumerate}
\item \textbf{Direction ?} Toujours celle de $\vec{u}$ (la droite support ne change pas).
\item \textbf{Sens ?} Celui de $\vec{u}$ si $k > 0$ ; l'opposé si $k < 0$.
\item \textbf{Longueur ?} On multiplie la longueur de $\vec{u}$ par $|k|$ (la valeur absolue, donc toujours un nombre positif).
\end{enumerate}
\end{aretenir}

\begin{exemplebox}[Exemple chiffré : construire $-2\vec{u}$]
On donne un vecteur $\vec{u}$ de norme $\|\vec{u}\| = 3$~cm. Décrivons le vecteur $-2\vec{u}$ :
\begin{itemize}
\item \textbf{Direction} : la même que celle de $\vec{u}$ ;
\item \textbf{Sens} : $k = -2 < 0$, donc le sens est \emph{opposé} à celui de $\vec{u}$ ;
\item \textbf{Norme} : $\|-2\vec{u}\| = |-2| \times \|\vec{u}\| = 2 \times 3 = 6$~cm.
\end{itemize}
Le vecteur $-2\vec{u}$ est donc deux fois plus long que $\vec{u}$, sur la même ligne de direction, mais pointant dans l'autre sens.
\end{exemplebox}

\begin{attention}[Piège classique : le sens ne se conserve pas toujours !]
Beaucoup d'élèves croient que $k\vec{u}$ a \emph{toujours} le même sens que $\vec{u}$. C'est \textbf{faux} dès que $k$ est négatif. Par exemple, $-3\vec{u}$ pointe dans le sens \emph{opposé} à $\vec{u}$. Retiens : le \textbf{signe de $k$ commande le sens}, la \textbf{valeur absolue de $k$ commande la longueur}. Autre erreur fréquente : écrire $\|k\vec{u}\| = k \times \|\vec{u}\|$ sans valeur absolue. Si $k = -2$, cela donnerait une longueur négative, ce qui est impossible ! La bonne formule est $\|k\vec{u}\| = |k| \times \|\vec{u}\|$.
\end{attention}

\courssub{Cas particuliers à connaître par c\oe ur}

Que se passe-t-il quand $k = 0$, quand $\vec{u} = \vec{0}$, ou quand $k = 1$ ou $k = -1$ ? La définition précédente excluait $k = 0$ et $\vec{u} = \vec{0}$ (un vecteur nul n'a pas de direction). On complète donc par des conventions naturelles.

\begin{propriete}[Cas particuliers du produit]
Pour tout vecteur $\vec{u}$ et tout réel $k$ :
\begin{itemize}
\item $0 \cdot \vec{u} = \vec{0}$ (multiplier par zéro « écrase » le vecteur) ;
\item $k \cdot \vec{0} = \vec{0}$ (le vecteur nul reste nul) ;
\item $1 \cdot \vec{u} = \vec{u}$ ;
\item $(-1) \cdot \vec{u} = -\vec{u}$ : c'est le \cle{vecteur opposé} de $\vec{u}$, de même direction, de même norme, mais de sens contraire.
\end{itemize}
En particulier : $k\vec{u} = \vec{0}$ si et seulement si $k = 0$ ou $\vec{u} = \vec{0}$.
\end{propriete}

Ces résultats sont intuitifs : si Moussa roule « zéro fois » la vitesse du cycliste, il ne bouge pas (vecteur nul). Et prendre $-1$ fois un déplacement, c'est exactement faire demi-tour et revenir en sens inverse sur la même distance.

\courssub{Interprétation géométrique : agrandir, réduire, retourner}

Multiplier un vecteur par un réel $k$, c'est lui appliquer une \emph{transformation} que l'on peut lire sur la valeur de $k$ :

\begin{center}
\begin{tabularx}{\linewidth}{|l|X|}\hline
\rowcolor{popblueL} \textbf{Valeur de $k$} & \textbf{Effet sur le vecteur $\vec{u}$} \\ \hline
$k > 1$ & Agrandissement : le vecteur est étiré, même sens. \\ \hline
$0 < k < 1$ & Réduction : le vecteur est raccourci, même sens. \\ \hline
$-1 < k < 0$ & Réduction \textbf{et} changement de sens. \\ \hline
$k < -1$ & Agrandissement \textbf{et} changement de sens. \\ \hline
$k = 1$ ou $k = -1$ & Longueur inchangée ; $-1$ inverse seulement le sens. \\ \hline
\end{tabularx}
\end{center}

\begin{popfigure}
\begin{center}
\begin{tikzpicture}[scale=0.9,>=Stealth]
% Ligne de direction
\draw[popline,dashed] (-5.6,0) -- (7.6,0);
\fill (0,0) circle (1.5pt) node[below left=1pt] {\small $O$};
% -2u
\draw[->,very thick,poppink] (0,0.9) -- (-4,0.9) node[midway,above] {\small $-2\vec{u}$};
% -u
\draw[->,very thick,poporange] (0,0.55) -- (-2,0.55) node[midway,above] {\small $-\vec{u}$};
% u
\draw[->,very thick,popblue] (0,0.2) -- (2,0.2) node[midway,above] {\small $\vec{u}$};
% 1/2 u
\draw[->,very thick,popgreen] (0,-0.35) -- (1,-0.35) node[midway,below] {\small $\dfrac{1}{2}\vec{u}$};
% 2u
\draw[->,very thick,poppurple] (0,-0.8) -- (4,-0.8) node[midway,below] {\small $2\vec{u}$};
% graduations
\foreach \x in {-4,-2,1,2,4} \draw[popmuted] (\x,-0.08) -- (\x,0.08);
\end{tikzpicture}
\end{center}
\legende{Les vecteurs $2\vec{u}$, $\frac{1}{2}\vec{u}$, $-\vec{u}$ et $-2\vec{u}$ construits à partir du même vecteur $\vec{u}$ : tous ont la même direction (la droite en pointillés) ; le signe de $k$ fixe le sens, $|k|$ fixe la longueur.}
\end{popfigure}

Observe bien la figure : les vecteurs $2\vec{u}$ et $\dfrac{1}{2}\vec{u}$ partent dans le même sens que $\vec{u}$ (car $2 > 0$ et $\dfrac{1}{2} > 0$), tandis que $-\vec{u}$ et $-2\vec{u}$ partent dans le sens opposé. Tous restent sur la même droite de direction.

\courssub{Construction géométrique à la règle et au compas}

\begin{methode}[Construire $k\vec{u}$ pas à pas]
On donne un vecteur $\vec{u} = \overrightarrow{AB}$ et un point $O$. Pour construire le point $M$ tel que $\overrightarrow{OM} = k\vec{u}$ :
\begin{enumerate}
\item \textbf{Direction} : trace la droite passant par $O$ et parallèle à la droite $(AB)$ (ou la droite $(AB)$ elle-même si $O$ est dessus).
\item \textbf{Sens} : regarde le signe de $k$. Si $k > 0$, pars dans le sens de $\vec{u}$ ; si $k < 0$, pars dans le sens opposé.
\item \textbf{Longueur} : avec le compas, mesure la longueur $AB$, puis reporte $|k|$ fois cette longueur à partir de $O$ dans le sens choisi. L'extrémité obtenue est le point $M$.
\end{enumerate}
\end{methode}

\begin{exemplebox}[Trois constructions modèles]
Soit $\vec{u}$ un vecteur de longueur $2$~cm.
\begin{itemize}
\item $\mathbf{k = 2}$ : on reporte $2$ fois la longueur de $\vec{u}$ dans le sens de $\vec{u}$. On obtient un vecteur de $4$~cm, même sens.
\item $\mathbf{k = \dfrac{1}{2}}$ : on plie la longueur en deux (médiatrice au compas) et on reporte la moitié, soit $1$~cm, dans le sens de $\vec{u}$.
\item $\mathbf{k = -3}$ : on reporte $3$ fois la longueur de $\vec{u}$, soit $6$~cm, mais dans le sens \emph{opposé} à $\vec{u}$.
\end{itemize}
\end{exemplebox}

\courssub{Lien avec la vie courante}

La multiplication d'un vecteur par un réel est partout autour de toi :

\begin{itemize}
\item \textbf{Doubler une vitesse.} Sur la route Yaoundé--Soa, si le vecteur vitesse du cycliste est $\vec{v}$, celui de la moto de Moussa est $2\vec{v}$ : même direction, même sens, double intensité.
\item \textbf{Reculer, faire demi-tour.} Une marche arrière ou un demi-tour se traduisent par un coefficient négatif : le déplacement retour est $-\vec{d}$ si l'aller est $\vec{d}$.
\item \textbf{Diluer une solution.} En chimie, au laboratoire du collège, diluer deux fois une solution, c'est multiplier la « quantité de principe actif » par $\dfrac{1}{2}$ : le vecteur concentration devient $\dfrac{1}{2}\vec{c}$.
\end{itemize}

\begin{popfigure}
\begin{center}
\begin{tikzpicture}[scale=0.85,>=Stealth]
% Route
\draw[very thick,popmuted] (0,0) -- (11,0);
\draw[popmuted,dashed] (0.5,0) -- (1.5,0);
\foreach \x in {1,2.5,4,5.5,7,8.5,10} \draw[white,thick] (\x,0) -- (\x+0.5,0);
% Cycliste
\node[popdark] (cyc) at (2,0.55) {\small Cycliste};
\draw[popdark] (2,0.1) circle (0.18);
\draw[popdark] (3,0.1) circle (0.18);
\draw[popdark] (2,0.1) -- (2.5,0.45) -- (3,0.1);
\draw[->,very thick,popgreen] (2,1.0) -- (4,1.0) node[midway,above] {\small $\vec{v}$};
% Moto
\node[popdark] at (6.5,0.55) {\small Moto de Moussa};
\draw[popdark] (6,0.12) circle (0.2);
\draw[popdark] (7.4,0.12) circle (0.2);
\draw[popdark,thick] (6,0.12) -- (6.6,0.5) -- (7.4,0.12);
\draw[->,very thick,poporange] (6,1.0) -- (10,1.0) node[midway,above] {\small $2\vec{v}$};
% Ville
\node[below] at (0.5,0) {\small Yaoundé};
\node[below] at (10.5,0) {\small Soa};
\end{tikzpicture}
\end{center}
\legende{Sur la route Yaoundé--Soa, la moto roule deux fois plus vite que le cycliste dans le même sens : son vecteur vitesse est $2\vec{v}$, deux fois plus long que $\vec{v}$, même direction et même sens.}
\end{popfigure}

\begin{savaistu}[Vecteurs en physique : forces, vitesses, accélérations]
En physique, la \textbf{force}, la \textbf{vitesse} et l'\textbf{accélération} sont des grandeurs \emph{vectorielles} : on les représente par des vecteurs. Dire qu'une force \emph{double}, c'est écrire simplement $\vec{F'} = 2\vec{F}$. Dire qu'un mobile fait demi-tour à même vitesse, c'est écrire $\vec{v'} = -\vec{v}$. Les ingénieurs qui conçoivent les ponts de Douala ou les lignes électriques de la SONATREL utilisent en permanence ces calculs de vecteurs pour équilibrer les forces. Les mathématiques que tu apprends aujourd'hui sont leurs outils de tous les jours !
\end{savaistu}

\courssub{Exercice résolu}

\begin{exoresolu}[Longueurs et sens de vecteurs multiples]
On donne un vecteur $\vec{u}$ tel que $\|\vec{u}\| = 2{,}5$~cm.
\begin{enumerate}
\item Calculer les normes des vecteurs $3\vec{u}$, $\dfrac{1}{5}\vec{u}$ et $-4\vec{u}$.
\item Pour chacun de ces vecteurs, préciser s'il a le même sens que $\vec{u}$ ou le sens opposé.
\item Construire sur une figure le vecteur $-4\vec{u}$.
\end{enumerate}
\tcblower
\textbf{Solution détaillée.}
\begin{enumerate}
\item On applique la formule $\|k\vec{u}\| = |k| \times \|\vec{u}\|$ :
\begin{align*}
\|3\vec{u}\| &= |3| \times 2{,}5 = 7{,}5~\text{cm} ;\\
\left\|\tfrac{1}{5}\vec{u}\right\| &= \left|\tfrac{1}{5}\right| \times 2{,}5 = 0{,}5~\text{cm} ;\\
\|-4\vec{u}\| &= |-4| \times 2{,}5 = 10~\text{cm}.
\end{align*}
\item Sens : $3 > 0$ donc $3\vec{u}$ a le \textbf{même sens} que $\vec{u}$ ; $\dfrac{1}{5} > 0$ donc $\dfrac{1}{5}\vec{u}$ a le \textbf{même sens} que $\vec{u}$ ; $-4 < 0$ donc $-4\vec{u}$ a le \textbf{sens opposé} à celui de $\vec{u}$.
\item Construction de $-4\vec{u}$ : on trace la droite portant $\vec{u}$ ; on choisit le sens \emph{opposé} à $\vec{u}$ (car $-4 < 0$) ; on reporte au compas $4$ fois la longueur de $\vec{u}$, soit $10$~cm. Le vecteur obtenu est $-4\vec{u}$.
\end{enumerate}
\textbf{Vérification :} une longueur est toujours positive ; $10$~cm $= 4 \times 2{,}5$~cm, le compte est bon.
\end{exoresolu}

\begin{exercice}[À toi de jouer]
Soit $\vec{u}$ un vecteur de norme $1{,}5$~cm.
\begin{enumerate}
\item Calcule $\|5\vec{u}\|$, $\left\|\dfrac{2}{3}\vec{u}\right\|$ et $\|-3\vec{u}\|$.
\item Le vecteur $-0{,}5\vec{u}$ a-t-il le même sens que $\vec{u}$ ? Justifie.
\item Trace un vecteur $\vec{u}$ horizontal vers la droite, puis construis $2\vec{u}$ et $-3\vec{u}$ à la règle et au compas.
\end{enumerate}
\end{exercice}

\begin{corrige}[Exercice : À toi de jouer]
\begin{enumerate}
\item $\|5\vec{u}\| = 5 \times 1{,}5 = 7{,}5$~cm ; $\left\|\dfrac{2}{3}\vec{u}\right\| = \dfrac{2}{3} \times 1{,}5 = 1$~cm ; $\|-3\vec{u}\| = 3 \times 1{,}5 = 4{,}5$~cm.
\item Non : $-0{,}5 < 0$, donc $-0{,}5\vec{u}$ a le \textbf{sens opposé} à celui de $\vec{u}$ (et sa longueur est $0{,}75$~cm, la moitié de celle de $\vec{u}$).
\item $2\vec{u}$ : même direction, même sens, longueur $3$~cm. $-3\vec{u}$ : même direction, sens opposé, longueur $4{,}5$~cm.
\end{enumerate}
\end{corrige}

\begin{aretenir}[Bilan de la section]
\begin{itemize}
\item Multiplier un vecteur $\vec{u}$ non nul par un réel $k$ non nul donne un vecteur $k\vec{u}$ de \textbf{même direction} que $\vec{u}$, de \textbf{même sens} si $k > 0$ et de \textbf{sens opposé} si $k < 0$, avec $\|k\vec{u}\| = |k| \times \|\vec{u}\|$.
\item Cas particuliers : $0 \cdot \vec{u} = \vec{0}$, $k \cdot \vec{0} = \vec{0}$, $1 \cdot \vec{u} = \vec{u}$, $(-1) \cdot \vec{u} = -\vec{u}$.
\item Géométriquement : $|k| > 1$ agrandit, $|k| < 1$ réduit, $k < 0$ retourne.
\item Cette opération sera la clé des sections suivantes : deux vecteurs dont l'un est le produit de l'autre par un réel seront dits \emph{colinéaires}, ce qui permettra de prouver le parallélisme et l'alignement.
\end{itemize}
\end{aretenir}

\coursec{Règles de calcul sur les produits de vecteurs}

Dans la section précédente, nous avons défini le produit d'un vecteur $\vec{u}$ par un nombre réel $k$ : c'est le vecteur $k\vec{u}$ qui a même direction que $\vec{u}$, un sens identique si $k>0$ et opposé si $k<0$, et une norme multipliée par $|k|$. Nous savons maintenant « multiplier » un vecteur par un nombre. Il nous reste à vérifier que cette opération se comporte bien avec les autres opérations déjà connues : l'addition des nombres réels, la multiplication des nombres réels entre eux, et l'addition des vecteurs. C'est l'objet de cette section : établir les \textbf{règles de calcul} qui permettront de transformer, simplifier et résoudre des égalités vectorielles exactement comme on le fait en calcul littéral.

\courssub{Les trois propriétés fondamentales}

Toutes les manipulations algébriques sur les vecteurs reposent sur trois propriétés qui prolongent celles des nombres. Elles sont admises au niveau de la classe de 3ème (elles seront démontrées rigoureusement en classe de Seconde à l'aide des coordonnées), mais nous en donnons ici une justification géométrique intuitive.

\begin{propriete}[Règles de calcul sur le produit d'un vecteur par un réel]
Soient $\vec{u}$ et $\vec{v}$ deux vecteurs quelconques, et $k$, $k'$ deux nombres réels quelconques. On a les trois propriétés suivantes :
\begin{enumerate}
    \item \textbf{Associativité (produit de nombres) :} \quad $k(k'\vec{u}) = (kk')\vec{u}$.
    \item \textbf{Distributivité par rapport à l'addition des réels :} \quad $(k + k')\vec{u} = k\vec{u} + k'\vec{u}$.
    \item \textbf{Distributivité par rapport à l'addition des vecteurs :} \quad $k(\vec{u} + \vec{v}) = k\vec{u} + k\vec{v}$.
\end{enumerate}
\end{propriete}

\begin{experience}[Justification géométrique de l'associativité : $k(k'\vec{u}) = (kk')\vec{u}$]
Prenons un vecteur $\vec{u} = \overrightarrow{AB}$.
\begin{enumerate}
    \item Construisons le vecteur $\vec{v} = k'\vec{u}$. Par définition, $\vec{v}$ a même direction que $\vec{u}$, sa norme est $|k'| \times AB$.
    \item Multiplions $\vec{v}$ par $k$ : le vecteur $k\vec{v} = k(k'\vec{u})$ a pour norme $|k| \times \|\vec{v}\| = |k| \times |k'| \times AB = |kk'| \times AB$.
    \item Construisons maintenant directement le vecteur $(kk')\vec{u}$. Sa norme est $|kk'| \times AB$ et sa direction est celle de $\vec{u}$ (sens selon le signe de $kk'$).
\end{enumerate}
Les deux vecteurs obtenus ont même origine (si on les place au même point), même direction, même sens et même norme. Ils sont donc égaux. On peut donc écrire sans ambiguïté $kk'\vec{u}$ sans parenthèses.
\end{experience}

\begin{experience}[Justification géométrique de la distributivité sur les réels : $(k+k')\vec{u} = k\vec{u} + k'\vec{u}$]
Prenons $\vec{u} = \overrightarrow{AB}$. Posons $\vec{v} = k\vec{u}$ et $\vec{w} = k'\vec{u}$. Ces deux vecteurs ont la même direction que $\vec{u}$ (donc ils sont \cle{colinéaires}).
\begin{itemize}
    \item La somme $\vec{v} + \vec{w}$ se construit par la règle de Chasles : on place $\vec{w}$ à la suite de $\vec{v}$. Comme ils sont sur la même droite, le résultant est un vecteur sur cette droite, de norme $\|\vec{v}\| + \|\vec{w}\|$ (si même sens) ou $|\|\vec{v}\| - \|\vec{w}\||$ (si sens opposés).
    \item Or $\|\vec{v}\| = |k| \cdot AB$ et $\|\vec{w}\| = |k'| \cdot AB$. La norme du résultant est bien $|k+k'| \cdot AB$ (en tenant compte des signes pour le sens).
    \item C'est exactement la définition de $(k+k')\vec{u}$.
\end{itemize}
\end{experience}

La troisième propriété, la distributivité par rapport à l'addition des vecteurs, est la plus riche géométriquement. Elle fait l'objet d'une démonstration guidée ci-dessous.

\courssub{Démonstration guidée : $k(\vec{u} + \vec{v}) = k\vec{u} + k\vec{v}$ pour $k > 0$}

Cette propriété est le cœur de la géométrie vectorielle : elle lie l'homothétie (multiplication par $k$) et la translation (addition). Nous la démontrons ici pour $k > 0$ à l'aide d'un agrandissement de figure (homothétie intuitive).

\begin{experience}[Démonstration guidée : distributivité par rapport à l'addition des vecteurs ($k>0$)]
\emph{But :} Montrer que pour $k > 0$, $k(\vec{u} + \vec{v}) = k\vec{u} + k\vec{v}$.

\begin{enumerate}
    \item \textbf{Construction de la somme $\vec{u} + \vec{v}$.} \\
    Soient $\vec{u} = \overrightarrow{AB}$ et $\vec{v} = \overrightarrow{BC}$. Par la règle de Chasles, $\vec{u} + \vec{v} = \overrightarrow{AC}$. Les points $A, B, C$ forment un triangle (ou sont alignés si $\vec{u}$ et $\vec{v}$ sont colinéaires).
    
    \item \textbf{Construction de $k(\vec{u} + \vec{v})$.} \\
    Puisque $k>0$, le vecteur $k(\vec{u} + \vec{v}) = k\overrightarrow{AC}$ est le vecteur $\overrightarrow{AC'}$ où $C'$ est l'image de $C$ par l'homothétie de centre $A$ et de rapport $k$. On a $A, C, C'$ alignés dans cet ordre et $AC' = k \cdot AC$.
    
    \item \textbf{Construction de $k\vec{u} + k\vec{v}$.} \\
    \begin{itemize}
        \item $k\vec{u} = k\overrightarrow{AB} = \overrightarrow{AB'}$ où $B'$ est l'image de $B$ par la même homothétie (centre $A$, rapport $k$). Donc $B'$ appartient à $(AB)$ et $AB' = k \cdot AB$.
        \item $k\vec{v} = k\overrightarrow{BC}$. Par définition de l'homothétie, l'image du segment $[BC]$ est le segment $[B'C']$ qui est \textbf{parallèle} à $(BC)$ et de longueur $k \cdot BC$. Le vecteur $\overrightarrow{B'C'}$ est donc bien $k\vec{v}$.
    \end{itemize}
    La somme $k\vec{u} + k\vec{v} = \overrightarrow{AB'} + \overrightarrow{B'C'} = \overrightarrow{AC'}$ (règle de Chasles dans le triangle $AB'C'$).
    
    \item \textbf{Conclusion.} \\
    Les deux constructions aboutissent au même vecteur $\overrightarrow{AC'}$. Donc $k(\vec{u} + \vec{v}) = k\vec{u} + k\vec{v}$.
\end{enumerate}

\medskip
\noindent \textbf{Interprétation géométrique :} L'homothétie de centre $A$ et de rapport $k$ transforme le triangle $ABC$ en triangle $AB'C'$. La relation de Chasles ($\overrightarrow{AB} + \overrightarrow{BC} = \overrightarrow{AC}$) est conservée par l'homothétie : $\overrightarrow{AB'} + \overrightarrow{B'C'} = \overrightarrow{AC'}$.
\end{experience}

\begin{popfigure}
\begin{tikzpicture}[scale=1.2, >=Stealth]
% Coordinates
\coordinate (A) at (0,0);
\coordinate (B) at (2,0.5);
\coordinate (C) at (3.5,2);
% Homothety k=2
\coordinate (Bp) at (4,1);
\coordinate (Cp) at (7,4);

% Original triangle ABC
\draw[popblue, thick] (A) -- (B) -- (C) -- cycle;
% Vectors u, v, u+v
\draw[popblue, very thick, ->] (A) -- (B) node[midway, below left] {$\vec{u}$};
\draw[poporange, very thick, ->] (B) -- (C) node[midway, right] {$\vec{v}$};
\draw[popgreen, very thick, ->] (A) -- (C) node[midway, above left] {$\vec{u}+\vec{v}$};
% Points labels
\node[popblue] at (A) [below left] {$A$};
\node[popblue] at (B) [below right] {$B$};
\node[popblue] at (C) [above right] {$C$};

% Homothetic triangle AB'C' (k=2)
\draw[popblue!50, thick, dashed] (A) -- (Bp) -- (Cp) -- cycle;
\draw[popblue!50, very thick, ->] (A) -- (Bp) node[midway, below left] {$2\vec{u}$};
\draw[poporange!50, very thick, ->] (Bp) -- (Cp) node[midway, right] {$2\vec{v}$};
\draw[popgreen!50, very thick, ->] (A) -- (Cp) node[midway, above left] {$2(\vec{u}+\vec{v})$};
% Points labels
\node[popblue!50] at (Bp) [below right] {$B'$};
\node[popblue!50] at (Cp) [above right] {$C'$};

% Homothety rays
\draw[popmuted, thin, dashed] (A) -- (Cp);
\draw[popmuted, thin, dashed] (A) -- (Bp);

% k label
\node[popdark, fill=popcream, rounded corners, inner sep=2pt] at (5.5, 0.5) {Homothétie $h(A, 2)$};
\end{tikzpicture}
\legende{Illustration de $2(\vec{u}+\vec{v}) = 2\vec{u} + 2\vec{v}$ : l'homothétie de centre $A$ et de rapport $2$ transforme le triangle $ABC$ en $AB'C'$ et conserve la relation de Chasles.}
\end{popfigure}

\begin{attention}[Cas $k < 0$]
La démonstration ci-dessus utilise une homothétie de rapport $k>0$ (agrandissement/réduction sans changement de sens par rapport au centre). Si $k < 0$, l'homothétie devient une homothétie de rapport négatif (rotation de $180^\circ$ + homothétie de rapport $|k|$). La relation $k(\vec{u}+\vec{v}) = k\vec{u}+k\vec{v}$ reste vraie, mais la figure géométrique est moins intuitive (les points $B'$ et $C'$ se trouvent de l'autre côté de $A$). En 3ème, on admet la propriété pour tout réel $k$ ; la preuve rigoureuse pour tout $k$ se fait en Seconde à l'aide des coordonnées.
\end{attention}

\courssub{Calcul vectoriel : simplification d'expressions}

Grâce aux trois propriétés de la \textbf{Propriété 1}, on peut manipuler les expressions vectorielles exactement comme des expressions littérales en $x$ et $y$, à une différence près : \textbf{on ne mélange jamais nombres et vecteurs}.

\begin{methode}[Réduire une expression vectorielle]
Pour simplifier une expression vectorielle $E$ contenant des sommes, différences et produits par des nombres :
\begin{enumerate}
    \item \textbf{Distribuer} les produits sur les sommes : développer $k(\vec{u} \pm \vec{v})$ en $k\vec{u} \pm k\vec{v}$.
    \item \textbf{Regrouper} les termes qui portent sur le \textbf{même vecteur} (même lettre, même origine si contexte géométrique).
    \item \textbf{Faire les calculs numériques} sur les coefficients (les nombres devant les vecteurs).
    \item \textbf{Écrire le résultat} sous la forme $a\vec{u} + b\vec{v} + \dots$ (\cle{combinaison linéaire}).
\end{enumerate}
\end{methode}

\begin{exemplebox}[Simplification d'une expression vectorielle]
Réduire l'expression : $\quad A = 2\vec{u} - 3(\vec{u} - \vec{v}) + 5\vec{v}$.

\begin{align*}
A &= 2\vec{u} - 3(\vec{u} - \vec{v}) + 5\vec{v} \\
  &= 2\vec{u} + \bigl[ (-3)\times\vec{u} + (-3)\times(-\vec{v}) \bigr] + 5\vec{v} \quad \text{(Distributivité de $-3$ sur la différence)} \\
  &= 2\vec{u} - 3\vec{u} + 3\vec{v} + 5\vec{v} \\
  &= (2 - 3)\vec{u} + (3 + 5)\vec{v} \quad \text{(Regroupement et factorisation)} \\
  &= (-1)\vec{u} + 8\vec{v} \\
  &= -\vec{u} + 8\vec{v}
\end{align*}
\emph{Résultat :} $\boxed{A = -\vec{u} + 8\vec{v}}$.
\end{exemplebox}

\begin{attention}[Pièges mortels à éviter absolument]
\begin{itemize}
    \item \textbf{Ne jamais additionner un nombre et un vecteur.} Écrire $3 + \vec{u}$, $5 - 2\vec{v}$ ou $k + \vec{w}$ \textbf{n'a aucun sens mathématique}. C'est comme ajouter des mètres et des kilogrammes.
    \item \textbf{Ne pas « simplifier » un vecteur par un nombre.} On ne peut pas passer de $2\vec{u} = 4$ à $\vec{u} = 2$. L'égalité entre un vecteur et un nombre est fausse (sauf si le vecteur est $\vec{0}$ et le nombre $0$).
    \item \textbf{Respecter l'ordre dans les produits.} On écrit $k\vec{u}$ (nombre $\times$ vecteur) et non $\vec{u}k$, même si c'est commutatif. La notation standard place le scalaire à gauche.
    \item \textbf{Ne pas confondre $-\vec{u}$ et $\vec{-u}$.} On note $-\vec{u}$ l'opposé de $\vec{u}$ (produit par $-1$). La notation $\vec{-u}$ n'existe pas.
\end{itemize}
\end{attention}

\courssub{Lien avec les coordonnées (rappel et admission)}

Dans un repère orthonormé $(O; \vec{i}, \vec{j})$, tout vecteur $\vec{u}$ se repère par ses coordonnées $(x ; y)$ : $\vec{u} = x\vec{i} + y\vec{j}$.
La multiplication par un réel $k$ s'effectue \textbf{coordonnée par coordonnée} :

\begin{propriete}[Coordonnées du produit $k\vec{u}$]
Si $\vec{u}$ a pour coordonnées $(x ; y)$ dans un repère, alors le vecteur $k\vec{u}$ a pour coordonnées $(kx ; ky)$.
\[
\vec{u} \begin{pmatrix} x \\ y \end{pmatrix} \quad \Longrightarrow \quad k\vec{u} \begin{pmatrix} kx \\ ky \end{pmatrix}
\]
\end{propriete}

\begin{aretenir}[Preuve admise en 3ème (rappel de Seconde)]
C'est une conséquence immédiate de la distributivité et de la définition des coordonnées :
\[
k\vec{u} = k(x\vec{i} + y\vec{j}) = (kx)\vec{i} + (ky)\vec{j}.
\]
Cette propriété permet de faire des calculs vectoriels très rapidement sans figure, en ne manipulant que des couples de nombres.
\end{aretenir}

\begin{exoresolu}[Calcul vectoriel avec coordonnées]
Dans un repère $(O; \vec{i}, \vec{j})$, on donne $\vec{u}\begin{pmatrix} 2 \\ -3 \end{pmatrix}$ et $\vec{v}\begin{pmatrix} -1 \\ 4 \end{pmatrix}$.
Calculer les coordonnées de $\vec{w} = 3\vec{u} - 2\vec{v}$.

\begin{align*}
\vec{w} &= 3\vec{u} - 2\vec{v} \\
&= 3\begin{pmatrix} 2 \\ -3 \end{pmatrix} - 2\begin{pmatrix} -1 \\ 4 \end{pmatrix} \\
&= \begin{pmatrix} 3\times 2 \\ 3\times (-3) \end{pmatrix} - \begin{pmatrix} 2\times (-1) \\ 2\times 4 \end{pmatrix} \quad \text{(Produit par un réel : coordonnée par coordonnée)} \\
&= \begin{pmatrix} 6 \\ -9 \end{pmatrix} - \begin{pmatrix} -2 \\ 8 \end{pmatrix} \\
&= \begin{pmatrix} 6 - (-2) \\ -9 - 8 \end{pmatrix} \quad \text{(Soustraction coordonnée par coordonnée)} \\
&= \begin{pmatrix} 8 \\ -17 \end{pmatrix}
\end{align*}
\emph{Résultat :} $\vec{w}\begin{pmatrix} 8 \\ -17 \end{pmatrix}$.
\end{exoresolu}

\begin{savaistu}[Application concrète : Redimensionnement d'image (Informatique graphique)]
Quand vous zoomez sur une photo sur votre téléphone (MTN/Orange Money, WhatsApp, etc.) ou que vous redimensionnez une image dans un logiciel comme Paint ou Photoshop, l'ordinateur applique une homothétie à chaque pixel.
Si un pixel a un vecteur position $\vec{p} = (x ; y)$ par rapport au centre de l'écran, et que vous zoomez par un facteur $k=2$ (200\%), le nouveau pixel se trouve à la position $2\vec{p} = (2x ; 2y)$. Les formules $k(\vec{u}+\vec{v}) = k\vec{u}+k\vec{v}$ garantissent que les formes géométriques (cercles, carrés, lettres) gardent leurs proportions : un cercle reste un cercle, un carré reste un carré. C'est la base de la \textbf{graphique vectorielle} (format SVG) utilisée pour les logos (comme celui de la CAMTEL ou de la FECAFOOT) qui ne pixellisent jamais, quel que soit le zoom.
\end{savaistu}

\begin{exercice}[Entraînement : Simplification et coordonnées]
\begin{enumerate}
    \item Réduire les expressions suivantes :
    \begin{enumerate}
        \item $A = 4\vec{a} - 2(\vec{a} + 3\vec{b}) + \vec{b}$
        \item $B = -5(\vec{u} - 2\vec{v}) + 3\vec{u} - \vec{v}$
        \item $C = \frac{1}{2}(2\vec{x} - 4\vec{y}) + 3\vec{y}$
    \end{enumerate}
    \item Dans un repère, $\vec{u}(3 ; -2)$ et $\vec{v}(-1 ; 5)$. Calculer les coordonnées de :
    \begin{enumerate}
        \item $2\vec{u} + 3\vec{v}$
        \item $- \vec{u} + 4\vec{v}$
        \item $\frac{1}{2}\vec{u} - \frac{3}{2}\vec{v}$
    \end{enumerate}
    \item Soit $ABCD$ un parallélogramme. Exprimer en fonction de $\overrightarrow{AB}$ et $\overrightarrow{AD}$ :
    \begin{enumerate}
        \item $\overrightarrow{AC}$
        \item $\overrightarrow{DB}$
        \item $2\overrightarrow{AB} - \overrightarrow{AD}$
    \end{enumerate}
\end{enumerate}
\end{exercice}

\begin{corrige}[Exercice n°1]
\begin{enumerate}
    \item \textbf{Simplifications :}
    \begin{enumerate}
        \item $A = 4\vec{a} - 2\vec{a} - 6\vec{b} + \vec{b} = 2\vec{a} - 5\vec{b}$.
        \item $B = -5\vec{u} + 10\vec{v} + 3\vec{u} - \vec{v} = -2\vec{u} + 9\vec{v}$.
        \item $C = \vec{x} - 2\vec{y} + 3\vec{y} = \vec{x} + \vec{y}$.
    \end{enumerate}
    \item \textbf{Coordonnées :}
    \begin{enumerate}
        \item $2\begin{pmatrix} 3 \\ -2 \end{pmatrix} + 3\begin{pmatrix} -1 \\ 5 \end{pmatrix} = \begin{pmatrix} 6 \\ -4 \end{pmatrix} + \begin{pmatrix} -3 \\ 15 \end{pmatrix} = \begin{pmatrix} 3 \\ 11 \end{pmatrix}$.
        \item $-\begin{pmatrix} 3 \\ -2 \end{pmatrix} + 4\begin{pmatrix} -1 \\ 5 \end{pmatrix} = \begin{pmatrix} -3 \\ 2 \end{pmatrix} + \begin{pmatrix} -4 \\ 20 \end{pmatrix} = \begin{pmatrix} -7 \\ 22 \end{pmatrix}$.
        \item $\frac{1}{2}\begin{pmatrix} 3 \\ -2 \end{pmatrix} - \frac{3}{2}\begin{pmatrix} -1 \\ 5 \end{pmatrix} = \begin{pmatrix} 1,5 \\ -1 \end{pmatrix} - \begin{pmatrix} -1,5 \\ 7,5 \end{pmatrix} = \begin{pmatrix} 3 \\ -8,5 \end{pmatrix}$.
    \end{enumerate}
    \item \textbf{Parallélogramme $ABCD$ :}
    \begin{enumerate}
        \item $\overrightarrow{AC} = \overrightarrow{AB} + \overrightarrow{BC} = \overrightarrow{AB} + \overrightarrow{AD}$.
        \item $\overrightarrow{DB} = \overrightarrow{DA} + \overrightarrow{AB} = -\overrightarrow{AD} + \overrightarrow{AB} = \overrightarrow{AB} - \overrightarrow{AD}$.
        \item $2\overrightarrow{AB} - \overrightarrow{AD}$ (déjà simplifié).
    \end{enumerate}
\end{enumerate}
\end{corrige}

\courssub{Bilan de la section}

\begin{aretenir}[Ce qu'il faut retenir absolument]
\begin{itemize}
    \item Les trois règles d'or du calcul vectoriel (pour tous réels $k, k'$ et vecteurs $\vec{u}, \vec{v}$) :
    \begin{align*}
        k(k'\vec{u}) &= (kk')\vec{u} & \text{(\cle{Associativité})} \\
        (k+k')\vec{u} &= k\vec{u} + k'\vec{u} & \text{(\cle{Distributivité} sur $\mathbb{R}$)} \\
        k(\vec{u}+\vec{v}) &= k\vec{u} + k\vec{v} & \text{(\cle{Distributivité} sur les vecteurs)}
    \end{align*}
    \item Pour simplifier une expression : \textbf{Développer $\rightarrow$ Regrouper les mêmes vecteurs $\rightarrow$ Calculer les coefficients}.
    \item \textbf{Interdit :} Additionner un nombre et un vecteur ($3 + \vec{u}$ est une absurdité).
    \item En coordonnées : si $\vec{u}(x;y)$ alors $k\vec{u}(kx;ky)$.
\end{itemize}
\end{aretenir}

Ces règles de calcul sont les « outils de base » qui nous permettront, dans la section suivante, de caractériser les \textbf{vecteurs colinéaires} et de comprendre le lien fondamental entre colinéarité, parallélisme des droites et alignement de points.

\coursec{Vecteurs colinéaires}

Dans la section précédente, nous avons appris à multiplier un vecteur par un nombre réel et à manipuler les règles de calcul qui en découlent : $k(\vec{u}+\vec{v}) = k\vec{u}+k\vec{v}$, $(k+k')\vec{u} = k\vec{u}+k'\vec{u}$, etc. Nous avons aussi observé un fait géométrique capital : lorsqu'on multiplie un vecteur $\vec{u}$ par un réel $k$, le vecteur $k\vec{u}$ obtenu garde \emph{la même direction} que $\vec{u}$ ; seuls sa longueur et éventuellement son sens changent. Cette observation va devenir une notion fondamentale : la \cle{colinéarité}. C'est elle qui permettra, dans les sections suivantes, de démontrer que des droites sont parallèles ou que des points sont alignés, sans aucune figure à l'échelle : uniquement par le calcul.

\courssub{Intuition : des vecteurs qui « vivent » sur des lignes parallèles}

Imagine la route nationale N3 entre Yaoundé et Ebolowa : plusieurs motos-taxis roulent sur cette route, certains vers Ebolowa, d'autres en sens inverse vers Yaoundé, à des vitesses différentes. Si l'on représente le déplacement de chaque moto par un vecteur, tous ces vecteurs ont \textbf{la même direction} (celle de la route), mais pas forcément le même sens, ni la même longueur. On dira qu'ils sont \cle{colinéaires}.

À l'inverse, un vecteur représentant le déplacement d'un piéton qui traverse la route perpendiculairement n'a pas la même direction : il n'est colinéaire à aucun des vecteurs précédents.

\begin{popfigure}
\begin{tikzpicture}[scale=0.9, >=Stealth]
  % droites parallèles (supports)
  \draw[popline, dashed] (-0.5,0) -- (6.5,0);
  \draw[popline, dashed] (-0.5,1.2) -- (6.5,1.2);
  \draw[popline, dashed] (-0.5,2.4) -- (6.5,2.4);
  % vecteurs colinéaires
  \draw[->, very thick, popblue] (0,0) -- (2,0) node[midway, below] {$\vec{u}$};
  \draw[->, very thick, poporange] (0.5,1.2) -- (3.5,1.2) node[midway, above] {$\vec{v} = 1{,}5\,\vec{u}$};
  \draw[->, very thick, popgreen] (4.5,2.4) -- (3.5,2.4) node[midway, above] {$\vec{w} = -0{,}5\,\vec{u}$};
  % vecteur non colinéaire
  \draw[->, very thick, poppink] (5,0) -- (6.2,2.2) node[midway, right] {$\vec{t}$};
  \node[popmuted, font=\small] at (6.9,0) {$(d_1)$};
  \node[popmuted, font=\small] at (6.9,1.2) {$(d_2)$};
  \node[popmuted, font=\small] at (6.9,2.4) {$(d_3)$};
\end{tikzpicture}
\legende{Les vecteurs $\vec{u}$, $\vec{v}$ et $\vec{w}$ sont portés par des droites parallèles : ils sont colinéaires. Le vecteur $\vec{t}$ n'a pas la même direction : il n'est colinéaire à aucun d'eux.}
\end{popfigure}

L'idée est donc simple : \textbf{deux vecteurs sont colinéaires lorsqu'ils ont la même direction}, c'est-à-dire lorsque les droites qui les portent (leurs \emph{supports}) sont parallèles ou confondues. Mais comment le traduire par le calcul, sans dessin ? Grâce à la multiplication par un réel : si $\vec{v}$ a la même direction que $\vec{u}$, alors $\vec{v}$ est un « agrandissement » ou une « réduction » de $\vec{u}$, éventuellement retourné. Autrement dit, $\vec{v} = k\vec{u}$ pour un certain réel $k$.

\courssub{Définition de la colinéarité}

\begin{definition}[Vecteurs colinéaires]
Soient $\vec{u}$ et $\vec{v}$ deux vecteurs \textbf{non nuls}. On dit que $\vec{u}$ et $\vec{v}$ sont \cle{colinéaires} s'il existe un nombre réel $k$ tel que :
\[ \vec{v} = k\,\vec{u}. \]
\end{definition}

Le réel $k$ est parfois appelé \emph{coefficient de colinéarité}. Son signe et sa valeur absolue nous renseignent immédiatement :

\begin{aretenir}[Lire le coefficient de colinéarité]
Si $\vec{v} = k\vec{u}$ avec $\vec{u} \neq \vec{0}$, alors :
\begin{itemize}
  \item si $k > 0$ : $\vec{u}$ et $\vec{v}$ ont le \textbf{même sens} ;
  \item si $k < 0$ : $\vec{u}$ et $\vec{v}$ sont de \textbf{sens contraires} ;
  \item dans tous les cas, les normes sont liées par $\|\vec{v}\| = |k| \times \|\vec{u}\|$.
\end{itemize}
\end{aretenir}

\begin{popfigure}
\begin{tikzpicture}[scale=1.1, >=Stealth]
  % Cas 1 : même sens
  \draw[->, very thick, popblue] (0,0) -- (2,0) node[midway, below] {$\vec{u}$};
  \draw[->, very thick, poporange] (0,0.6) -- (3,0.6) node[midway, above] {$\vec{v} = 1{,}5\,\vec{u}$};
  \node[popdark, font=\small] at (1.5,-1) {Cas $k = 1{,}5 > 0$ : même sens, $\|\vec{v}\| = 1{,}5\,\|\vec{u}\|$};
  % Cas 2 : sens contraires
  \draw[->, very thick, popblue] (6,0) -- (8,0) node[midway, below] {$\vec{u}$};
  \draw[->, very thick, popgreen] (8,0.6) -- (7,0.6) node[midway, above] {$\vec{w} = -0{,}5\,\vec{u}$};
  \node[popdark, font=\small] at (7,-1) {Cas $k = -0{,}5 < 0$ : sens contraires, $\|\vec{w}\| = 0{,}5\,\|\vec{u}\|$};
\end{tikzpicture}
\legende{À gauche, $\vec{v} = 1{,}5\,\vec{u}$ : même direction, même sens, une fois et demie plus long. À droite, $\vec{w} = -0{,}5\,\vec{u}$ : même direction, sens contraire, deux fois plus court.}
\end{popfigure}

\begin{propriete}[Convention sur le vecteur nul]
Le vecteur nul $\vec{0}$ est colinéaire à \textbf{tout} vecteur. En effet, pour tout vecteur $\vec{u}$, on peut écrire $\vec{0} = 0 \times \vec{u}$.
\end{propriete}

Cette convention est très pratique : elle évite d'avoir à traiter des cas particuliers dans les théorèmes. Retiens simplement : le vecteur nul n'a pas de direction, donc on décide qu'il « va dans toutes les directions », c'est-à-dire qu'il est colinéaire à n'importe quel vecteur.

\begin{attention}[Colinéaires ne veut pas dire égaux !]
L'erreur la plus fréquente est de confondre « colinéaires » et « égaux ». Deux vecteurs \textbf{égaux} ont même direction, même sens \emph{et} même norme : c'est le cas particulier $k = 1$. Deux vecteurs \textbf{colinéaires} ont seulement la même direction : ils peuvent avoir des normes différentes et même des sens contraires. Par exemple, $\vec{v} = -3\vec{u}$ est colinéaire à $\vec{u}$, mais il est trois fois plus long et de sens opposé : il n'est certainement pas égal à $\vec{u}$. L'égalité implique la colinéarité, mais la réciproque est fausse.
\end{attention}

\courssub{Méthode : montrer que deux vecteurs sont colinéaires}

En géométrie, les vecteurs ne sont pas toujours donnés par des coordonnées : ils sont souvent définis par des points d'une figure. La stratégie consiste alors à utiliser la \textbf{relation de Chasles} et les règles de calcul de la section précédente pour exprimer l'un des vecteurs en fonction de l'autre.

\begin{methode}[Montrer que deux vecteurs sont colinéaires]
Pour montrer que $\overrightarrow{AB}$ et $\overrightarrow{CD}$ sont colinéaires :
\begin{enumerate}
  \item À l'aide de la relation de Chasles, décomposer $\overrightarrow{CD}$ en passant par des points intermédiaires bien choisis (souvent $A$ et $B$).
  \item Utiliser les hypothèses de l'énoncé (milieu, relations vectorielles données) pour simplifier.
  \item Arriver à une égalité de la forme $\overrightarrow{CD} = k\,\overrightarrow{AB}$.
  \item Conclure : $\overrightarrow{AB}$ et $\overrightarrow{CD}$ sont colinéaires (et préciser le coefficient $k$).
\end{enumerate}
\end{methode}

\begin{exoresolu}[Colinéarité par le calcul vectoriel]
\textbf{Énoncé.} Soit $ABC$ un triangle. On définit les points $M$ et $N$ par $\overrightarrow{AM} = 2\overrightarrow{AB}$ et $\overrightarrow{AN} = 2\overrightarrow{AC}$. Montrer que les vecteurs $\overrightarrow{BC}$ et $\overrightarrow{MN}$ sont colinéaires.
\tcblower
\textbf{Solution détaillée.} On cherche à exprimer $\overrightarrow{MN}$ en fonction de $\overrightarrow{BC}$. Par la relation de Chasles :
\begin{align*}
\overrightarrow{MN} &= \overrightarrow{MA} + \overrightarrow{AN} \\
&= -\overrightarrow{AM} + \overrightarrow{AN} \\
&= -2\overrightarrow{AB} + 2\overrightarrow{AC} \\
&= 2\left(\overrightarrow{AC} - \overrightarrow{AB}\right) \\
&= 2\left(\overrightarrow{BA} + \overrightarrow{AC}\right) \quad \text{(car } -\overrightarrow{AB} = \overrightarrow{BA}\text{)} \\
&= 2\,\overrightarrow{BC} \quad \text{(relation de Chasles).}
\end{align*}
On a obtenu $\overrightarrow{MN} = 2\,\overrightarrow{BC}$ : les vecteurs $\overrightarrow{MN}$ et $\overrightarrow{BC}$ sont donc \textbf{colinéaires}, de même sens (car $k = 2 > 0$), et $MN = 2\,BC$.
\end{exoresolu}

\courssub{Critère pratique avec les coordonnées}

Lorsque les vecteurs sont donnés par leurs coordonnées dans un repère, il existe un test de colinéarité rapide, que nous admettons à ce niveau et que tu démontreras en classe de seconde.

\begin{propriete}[Critère de colinéarité (admis)]
Dans un repère, soient $\vec{u}(x\,;y)$ et $\vec{v}(x'\,;y')$. Les vecteurs $\vec{u}$ et $\vec{v}$ sont colinéaires si et seulement si :
\[ xy' - x'y = 0. \]
\end{propriete}

Ce critère se retient facilement comme un « produit en croix » : on multiplie les coordonnées en diagonale et on fait la différence. Si le résultat est nul, les vecteurs sont colinéaires ; sinon, ils ne le sont pas.

\begin{exemplebox}[Vérifier une colinéarité avec les coordonnées]
Soient $\vec{u}(2\,;-3)$ et $\vec{v}(-4\,;6)$. Sont-ils colinéaires ?
\begin{itemize}
  \item \textbf{Par le critère :} $xy' - x'y = 2 \times 6 - (-4) \times (-3) = 12 - 12 = 0$. Le critère est vérifié : $\vec{u}$ et $\vec{v}$ sont colinéaires.
  \item \textbf{Par le coefficient :} on remarque que $-4 = -2 \times 2$ et $6 = -2 \times (-3)$, donc $\vec{v} = -2\,\vec{u}$. Le coefficient est $k = -2 < 0$ : les vecteurs sont de sens contraires et $\|\vec{v}\| = 2\|\vec{u}\|$.
\end{itemize}
\textbf{Contre-exemple :} soient $\vec{u}(2\,;-3)$ et $\vec{t}(4\,;5)$. On calcule $2 \times 5 - 4 \times (-3) = 10 + 12 = 22 \neq 0$ : les vecteurs $\vec{u}$ et $\vec{t}$ ne sont \textbf{pas} colinéaires. On peut aussi le voir directement : l'abscisse a été multipliée par $2$ ($2 \to 4$) mais pas l'ordonnée ($-3 \to 5$ et non $-6$) ; aucun réel $k$ ne convient pour les deux coordonnées à la fois.
\end{exemplebox}

\begin{attention}[Le même coefficient pour les deux coordonnées]
Pour conclure que $\vec{v} = k\vec{u}$, le \textbf{même} réel $k$ doit convenir pour l'abscisse \emph{et} pour l'ordonnée. Si l'abscisse est multipliée par $3$ et l'ordonnée par $2$, les vecteurs ne sont pas colinéaires ! En cas de doute, utilise toujours le critère $xy' - x'y = 0$, qui est infaillible.
\end{attention}

\begin{exoresolu}[Trouver une coordonnée pour obtenir la colinéarité]
\textbf{Énoncé.} Dans un repère, on donne $\vec{u}(3\,;-2)$ et $\vec{v}(m\,;4)$ où $m$ est un nombre réel. Déterminer $m$ pour que $\vec{u}$ et $\vec{v}$ soient colinéaires.
\tcblower
\textbf{Solution détaillée.} D'après le critère de colinéarité, $\vec{u}$ et $\vec{v}$ sont colinéaires si et seulement si :
\begin{align*}
3 \times 4 - m \times (-2) &= 0 \\
12 + 2m &= 0 \\
2m &= -12 \\
m &= -6.
\end{align*}
\textbf{Vérification :} $\vec{v}(-6\,;4)$ et l'on a bien $\vec{v} = -2\,\vec{u}$ car $-6 = -2 \times 3$ et $4 = -2 \times (-2)$. Les vecteurs sont colinéaires, de sens contraires.
\end{exoresolu}

\courssub{Pour aller plus loin (hors exigibilité)}

\begin{savaistu}[D'où vient le mot « colinéaire » ?]
Le mot \emph{colinéaire} vient du latin \emph{linea}, qui signifie « ligne » (le préfixe \emph{co-} signifie « ensemble », comme dans \emph{copain} ou \emph{coopérer}). Des vecteurs colinéaires sont littéralement des vecteurs qui « vivent ensemble sur des lignes » : des lignes parallèles ou confondues. En classe de seconde et au lycée, tu rencontreras un outil puissant appelé le \cle{déterminant} de deux vecteurs : pour $\vec{u}(x\,;y)$ et $\vec{v}(x'\,;y')$, on note $\det(\vec{u},\vec{v}) = xy' - x'y$. Le critère de colinéarité s'énonce alors très élégamment : « deux vecteurs sont colinéaires si et seulement si leur déterminant est nul ». Le déterminant a aussi une magnifique interprétation géométrique : sa valeur absolue est l'aire du parallélogramme construit sur les deux vecteurs. Dire que le déterminant est nul, c'est dire que ce parallélogramme est « aplati »... c'est-à-dire que les deux vecteurs ont la même direction ! Tu retrouves ainsi l'intuition de départ.
\end{savaistu}

\begin{exercice}[À toi de jouer]
\begin{enumerate}
  \item Dans chaque cas, dire si les vecteurs sont colinéaires et, si oui, donner le coefficient $k$ tel que $\vec{v} = k\vec{u}$ :
  \begin{enumerate}
    \item $\vec{u}(1\,;4)$ et $\vec{v}(3\,;12)$ ;
    \item $\vec{u}(-2\,;5)$ et $\vec{v}(6\,;-15)$ ;
    \item $\vec{u}(2\,;-1)$ et $\vec{v}(-4\,;3)$.
  \end{enumerate}
  \item Soit $ABC$ un triangle et $I$ le milieu de $[AB]$. On définit $J$ par $\overrightarrow{AJ} = \dfrac{1}{2}\overrightarrow{AC}$. Montrer que $\overrightarrow{IJ}$ et $\overrightarrow{BC}$ sont colinéaires.
  \item Déterminer le réel $a$ pour que $\vec{u}(a\,;6)$ et $\vec{v}(4\,;-3)$ soient colinéaires.
\end{enumerate}
\end{exercice}

\begin{corrige}[Exercice]
\begin{enumerate}
  \item \begin{enumerate}
    \item $1 \times 12 - 3 \times 4 = 12 - 12 = 0$ : colinéaires, $\vec{v} = 3\vec{u}$ ($k = 3$, même sens).
    \item $(-2) \times (-15) - 6 \times 5 = 30 - 30 = 0$ : colinéaires, $\vec{v} = -3\vec{u}$ ($k = -3$, sens contraires).
    \item $2 \times 3 - (-4) \times (-1) = 6 - 4 = 2 \neq 0$ : non colinéaires.
  \end{enumerate}
  \item $\overrightarrow{IJ} = \overrightarrow{IA} + \overrightarrow{AJ} = -\dfrac{1}{2}\overrightarrow{AB} + \dfrac{1}{2}\overrightarrow{AC} = \dfrac{1}{2}\left(\overrightarrow{AC} - \overrightarrow{AB}\right) = \dfrac{1}{2}\overrightarrow{BC}$. Donc $\overrightarrow{IJ}$ et $\overrightarrow{BC}$ sont colinéaires, de coefficient $k = \dfrac{1}{2}$.
  \item Il faut $a \times (-3) - 4 \times 6 = 0$, soit $-3a - 24 = 0$, d'où $a = -8$. Vérification : $\vec{u}(-8\,;6) = -2\,\vec{v}$, donc $\vec{u}$ et $\vec{v}$ sont bien colinéaires.
\end{enumerate}
\end{corrige}

La colinéarité est maintenant bien définie et tu sais la vérifier aussi bien par le calcul vectoriel (relation de Chasles) que par les coordonnées (critère $xy' - x'y = 0$). Dans la section suivante, nous exploiterons cette notion pour démontrer le parallélisme de droites et l'alignement de points : c'est là que la colinéarité révèle toute sa puissance.

\coursec{Colinéarité, parallélisme et alignement}

Dans la section précédente, nous avons découvert la notion de \cle{vecteurs colinéaires} : deux vecteurs sont colinéaires lorsque l'un est le produit de l'autre par un nombre réel, c'est-à-dire lorsqu'ils ont la \emph{même direction}. Cette notion, qui peut sembler abstraite, est en réalité un outil extraordinairement puissant : elle permet de \textbf{démontrer} des propriétés géométriques — parallélisme de droites, alignement de points — par de simples \emph{calculs} sur les vecteurs. C'est toute la magie de la géométrie vectorielle : remplacer des raisonnements de figure par des calculs rigoureux. C'est exactement ce que nous allons apprendre dans cette section, et c'est une compétence très attendue au BEPC.

\courssub{Colinéarité et parallélisme des droites}

\textbf{Intuition.} Imagine deux motos qui roulent sur l'axe Yaoundé--Douala, sur deux voies parallèles. Leurs déplacements ont la même direction : si la seconde moto roule une fois et demie plus loin que la première, son déplacement est $1{,}5$ fois celui de la première. Leurs vecteurs déplacement sont colinéaires, et leurs trajectoires sont parallèles. Ce constat est en fait un théorème fondamental.

\begin{propriete}[Théorème 1 : colinéarité et parallélisme]
Soient $A$, $B$, $C$, $D$ quatre points tels que $A \neq B$ et $C \neq D$.
\[ (AB) \parallel (CD) \quad \Longleftrightarrow \quad \overrightarrow{AB} \text{ et } \overrightarrow{CD} \text{ sont colinéaires.} \]
Autrement dit : les droites $(AB)$ et $(CD)$ sont parallèles si et seulement s'il existe un réel $k$ tel que $\overrightarrow{CD} = k\,\overrightarrow{AB}$.
\end{propriete}

\begin{experience}[Démonstration guidée du théorème 1 (sens direct)]
Nous démontrons le sens direct : \emph{si les droites sont parallèles, alors les vecteurs sont colinéaires}.

\textbf{Hypothèse :} les droites $(AB)$ et $(CD)$ sont parallèles, avec $A \neq B$ et $C \neq D$.

\textbf{Étape 1 :} Par hypothèse, $(AB)$ et $(CD)$ sont parallèles, donc elles ont la \emph{même direction}. Or la direction d'une droite $(AB)$ est précisément la direction du vecteur $\overrightarrow{AB}$. Les vecteurs $\overrightarrow{AB}$ et $\overrightarrow{CD}$ ont donc la même direction.

\textbf{Étape 2 :} Construisons le point $E$ tel que $\overrightarrow{AE} = \overrightarrow{CD}$ (c'est-à-dire le point $E$ tel que $ACDE$ soit un parallélogramme). Le vecteur $\overrightarrow{AE}$ a même direction que $\overrightarrow{CD}$, donc même direction que $\overrightarrow{AB}$. Comme $\overrightarrow{AE}$ et $\overrightarrow{AB}$ partent du même point $A$ et ont la même direction, le point $E$ appartient à la droite $(AB)$.

\textbf{Étape 3 :} Puisque $E \in (AB)$, les vecteurs $\overrightarrow{AE}$ et $\overrightarrow{AB}$ ont même direction : il existe donc un réel $k$ tel que $\overrightarrow{AE} = k\,\overrightarrow{AB}$ (le réel $k$ est le rapport des longueurs, affecté d'un signe selon le sens : positif si les vecteurs ont le même sens, négatif sinon).

\textbf{Conclusion :} $\overrightarrow{CD} = \overrightarrow{AE} = k\,\overrightarrow{AB}$ : les vecteurs $\overrightarrow{CD}$ et $\overrightarrow{AB}$ sont colinéaires. $\blacksquare$

La réciproque (si les vecteurs sont colinéaires, alors les droites sont parallèles) se justifie de même : si $\overrightarrow{CD} = k\,\overrightarrow{AB}$, les deux vecteurs ont même direction, donc leurs supports $(AB)$ et $(CD)$ sont des droites de même direction, c'est-à-dire parallèles.
\end{experience}

\begin{popfigure}
\begin{center}
\begin{tikzpicture}[scale=1.1,>=Stealth]
\draw[popline] (-0.5,0) -- (5.5,0);
\draw[popline] (-0.5,1.8) -- (5.5,1.8);
\coordinate (A) at (0.3,0); \coordinate (B) at (2.3,0);
\coordinate (C) at (1.2,1.8); \coordinate (D) at (4.2,1.8);
\draw[->,very thick,popblue] (A) -- (B) node[midway,below] {$\overrightarrow{AB}$};
\draw[->,very thick,poporange] (C) -- (D) node[midway,above] {$\overrightarrow{CD} = 1{,}5\,\overrightarrow{AB}$};
\fill[popblue] (A) circle (1.5pt) node[below left] {$A$};
\fill[popblue] (B) circle (1.5pt) node[below right] {$B$};
\fill[poporange] (C) circle (1.5pt) node[above left] {$C$};
\fill[poporange] (D) circle (1.5pt) node[above right] {$D$};
\node[popmuted] at (5.3,0.9) {$(AB) \parallel (CD)$};
\end{tikzpicture}
\end{center}
\legende{Les vecteurs $\overrightarrow{AB}$ et $\overrightarrow{CD}$ sont colinéaires ($\overrightarrow{CD} = 1{,}5\,\overrightarrow{AB}$) : les droites $(AB)$ et $(CD)$ sont parallèles.}
\end{popfigure}

\begin{methode}[Démontrer que deux droites sont parallèles]
Pour prouver que $(AB) \parallel (CD)$ :
\begin{enumerate}
\item Exprimer $\overrightarrow{AB}$ et $\overrightarrow{CD}$ en fonction d'un même vecteur (ou des mêmes vecteurs) de référence, en utilisant la relation de Chasles et les hypothèses de l'énoncé.
\item Mettre en évidence un réel $k$ tel que $\overrightarrow{CD} = k\,\overrightarrow{AB}$.
\item Conclure : les vecteurs sont colinéaires, donc les droites $(AB)$ et $(CD)$ sont parallèles (théorème 1).
\end{enumerate}
\end{methode}

\courssub{Colinéarité et alignement des points}

\textbf{Intuition.} Trois poteaux électriques $A$, $B$, $C$ plantés le long d'une route droite sont \emph{alignés}. Comment le vérifier avec des vecteurs ? Si en partant de $A$, le déplacement vers $C$ est simplement un agrandissement du déplacement vers $B$ (par exemple $\overrightarrow{AC} = 2{,}5\,\overrightarrow{AB}$), alors $B$ et $C$ sont sur la même ligne issue de $A$ : les trois points sont alignés.

\begin{propriete}[Théorème 2 : colinéarité et alignement]
Soient $A$, $B$, $C$ trois points tels que $A \neq B$.
\[ A,\ B,\ C \text{ sont alignés} \quad \Longleftrightarrow \quad \overrightarrow{AB} \text{ et } \overrightarrow{AC} \text{ sont colinéaires.} \]
Autrement dit : les points $A$, $B$, $C$ sont alignés si et seulement s'il existe un réel $k$ tel que $\overrightarrow{AC} = k\,\overrightarrow{AB}$.
\end{propriete}

\begin{experience}[Justification du théorème 2]
\textbf{Sens direct.} Si $A$, $B$, $C$ sont alignés, alors $C$ appartient à la droite $(AB)$. Les vecteurs $\overrightarrow{AB}$ et $\overrightarrow{AC}$ ont donc la même direction (celle de cette droite) : ils sont colinéaires.

\textbf{Réciproque.} Supposons $\overrightarrow{AB}$ et $\overrightarrow{AC}$ colinéaires : il existe un réel $k$ tel que $\overrightarrow{AC} = k\,\overrightarrow{AB}$.
\begin{itemize}
\item Les vecteurs $\overrightarrow{AB}$ et $\overrightarrow{AC}$ ont donc la \emph{même direction} : leurs supports, les droites $(AB)$ et $(AC)$, sont parallèles.
\item Or ces deux droites ont un \cle{point commun} : le point $A$.
\item Deux droites parallèles qui ont un point commun sont \textbf{confondues}. Donc $(AC) = (AB)$, ce qui signifie que $C \in (AB)$ : les points $A$, $B$, $C$ sont alignés. $\blacksquare$
\end{itemize}
C'est précisément l'existence du point commun $A$ qui transforme le parallélisme en alignement : c'est le point clé du raisonnement.
\end{experience}

\begin{popfigure}
\begin{center}
\begin{tikzpicture}[scale=1.2,>=Stealth]
\draw[popline] (-0.5,0) -- (6,0);
\coordinate (A) at (0,0); \coordinate (B) at (2,0); \coordinate (C) at (5,0);
\draw[->,very thick,popblue] (A) -- (B) node[midway,above] {$\overrightarrow{AB}$};
\draw[->,very thick,poporange] (0,-0.15) -- (5,-0.15) node[midway,below] {$\overrightarrow{AC} = 2{,}5\,\overrightarrow{AB}$};
\fill[popdark] (A) circle (1.5pt) node[above left=1pt] {$A$};
\fill[popdark] (B) circle (1.5pt) node[above=1pt] {$B$};
\fill[popdark] (C) circle (1.5pt) node[above right=1pt] {$C$};
\end{tikzpicture}
\end{center}
\legende{$\overrightarrow{AC} = 2{,}5\,\overrightarrow{AB}$ : les vecteurs sont colinéaires avec un point commun $A$, donc $A$, $B$, $C$ sont alignés.}
\end{popfigure}

\begin{methode}[Démontrer que trois points sont alignés]
Pour prouver que $A$, $B$, $C$ sont alignés :
\begin{enumerate}
\item Choisir deux vecteurs \emph{de même origine} : $\overrightarrow{AB}$ et $\overrightarrow{AC}$.
\item Décomposer chacun d'eux avec la relation de Chasles, en fonction d'un même vecteur de référence (souvent donné par l'énoncé).
\item Montrer qu'il existe un réel $k$ tel que $\overrightarrow{AC} = k\,\overrightarrow{AB}$.
\item Conclure en deux temps : les vecteurs sont colinéaires, \emph{et} ils ont le point $A$ en commun ; donc $A$, $B$, $C$ sont alignés.
\end{enumerate}
\end{methode}

\begin{attention}[Le piège du point commun]
Pour démontrer un \textbf{alignement}, il ne suffit pas que deux vecteurs soient colinéaires ! La colinéarité de $\overrightarrow{AB}$ et $\overrightarrow{CD}$ ne donne que le \emph{parallélisme} des droites $(AB)$ et $(CD)$, qui peuvent être distinctes. Pour conclure que des points sont alignés, il faut choisir deux vecteurs ayant un \cle{point commun} : par exemple $\overrightarrow{AB}$ et $\overrightarrow{AC}$ (origine commune $A$), et non $\overrightarrow{AB}$ et $\overrightarrow{CD}$. Rédige toujours la conclusion complète : « $\overrightarrow{AB}$ et $\overrightarrow{AC}$ sont colinéaires et ont le point $A$ en commun, donc $A$, $B$, $C$ sont alignés. »
\end{attention}

\courssub{Applications : parallélogrammes et trapèzes}

La colinéarité permet de reconnaître des quadrilatères particuliers par le calcul vectoriel :
\begin{itemize}
\item $ABCD$ est un \cle{parallélogramme} si et seulement si $\overrightarrow{AB} = \overrightarrow{DC}$ (c'est-à-dire $\overrightarrow{DC} = 1 \times \overrightarrow{AB}$ : colinéarité avec $k = 1$, même longueur et même sens).
\item $ABCD$ est un \cle{trapèze} de bases $[AB]$ et $[DC]$ lorsque $\overrightarrow{DC} = k\,\overrightarrow{AB}$ avec $k \neq 0$ et $k \neq 1$ : les bases sont parallèles mais de longueurs différentes.
\end{itemize}

\begin{popfigure}
\begin{center}
\begin{tikzpicture}[scale=1.1,>=Stealth]
\coordinate (A) at (0,0); \coordinate (B) at (4,0);
\coordinate (D) at (1,2); \coordinate (C) at (3,2);
\draw[thick,popdark] (A) -- (B) -- (C) -- (D) -- cycle;
\draw[->,very thick,popblue] (A) -- (B) node[midway,below] {$\overrightarrow{AB}$};
\draw[->,very thick,poporange] (D) -- (C) node[midway,above] {$\overrightarrow{DC} = \tfrac{1}{2}\overrightarrow{AB}$};
\fill[popdark] (A) circle (1.5pt) node[below left] {$A$};
\fill[popdark] (B) circle (1.5pt) node[below right] {$B$};
\fill[popdark] (C) circle (1.5pt) node[above right] {$C$};
\fill[popdark] (D) circle (1.5pt) node[above left] {$D$};
\end{tikzpicture}
\end{center}
\legende{$\overrightarrow{DC} = \frac{1}{2}\overrightarrow{AB}$ : les droites $(AB)$ et $(DC)$ sont parallèles, donc $ABCD$ est un trapèze de bases $[AB]$ et $[DC]$.}
\end{popfigure}

\begin{exoresolu}[Le théorème des milieux démontré par les vecteurs]
\textbf{Énoncé.} Soit $ABC$ un triangle. On note $M$ le milieu de $[AB]$ et $N$ le milieu de $[AC]$.
\begin{enumerate}
\item Exprimer $\overrightarrow{MN}$ en fonction de $\overrightarrow{AB}$ et $\overrightarrow{AC}$.
\item En déduire que $\overrightarrow{MN} = \dfrac{1}{2}\overrightarrow{BC}$.
\item Conclure sur les droites $(MN)$ et $(BC)$.
\end{enumerate}
\tcblower
\textbf{Solution détaillée.}
\begin{enumerate}
\item D'après la relation de Chasles, en passant par le point $A$ :
\[ \overrightarrow{MN} = \overrightarrow{MA} + \overrightarrow{AN}. \]
\item Comme $M$ est le milieu de $[AB]$, on a $\overrightarrow{MA} = -\dfrac{1}{2}\overrightarrow{AB}$. Comme $N$ est le milieu de $[AC]$, on a $\overrightarrow{AN} = \dfrac{1}{2}\overrightarrow{AC}$. Donc :
\begin{align*}
\overrightarrow{MN} &= -\frac{1}{2}\overrightarrow{AB} + \frac{1}{2}\overrightarrow{AC} = \frac{1}{2}\left(\overrightarrow{AC} - \overrightarrow{AB}\right) = \frac{1}{2}\left(\overrightarrow{BA} + \overrightarrow{AC}\right) = \frac{1}{2}\overrightarrow{BC}.
\end{align*}
\item L'égalité $\overrightarrow{MN} = \dfrac{1}{2}\overrightarrow{BC}$ montre que les vecteurs $\overrightarrow{MN}$ et $\overrightarrow{BC}$ sont colinéaires. D'après le théorème 1, les droites $(MN)$ et $(BC)$ sont \textbf{parallèles}. De plus, $MN = \dfrac{1}{2}BC$ : le segment joignant les milieux de deux côtés mesure la moitié du troisième côté.
\end{enumerate}
Nous venons de redémontrer par les vecteurs le célèbre \emph{théorème des milieux}, vu en classe de 4\textsuperscript{e} !
\end{exoresolu}

\begin{exoresolu}[Démontrer un alignement et un parallélisme]
\textbf{Énoncé.} Soit $ABC$ un triangle. On définit les points $P$ et $Q$ par :
\[ \overrightarrow{AP} = 2\overrightarrow{AB} \quad \text{et} \quad \overrightarrow{AQ} = \overrightarrow{AB} + 2\overrightarrow{BC}. \]
\begin{enumerate}
\item Démontrer que les points $A$, $B$ et $P$ sont alignés.
\item Exprimer $\overrightarrow{BQ}$ en fonction de $\overrightarrow{BC}$, puis démontrer que les points $B$, $C$ et $Q$ sont alignés.
\end{enumerate}
\tcblower
\textbf{Solution détaillée.}
\begin{enumerate}
\item Par définition du point $P$, on a $\overrightarrow{AP} = 2\overrightarrow{AB}$. Les vecteurs $\overrightarrow{AP}$ et $\overrightarrow{AB}$ sont donc colinéaires (avec $k = 2$) et ils ont le point $A$ en commun : les points $A$, $B$, $P$ sont \textbf{alignés}. De plus, $P$ est le point de la demi-droite $[AB)$ tel que $AP = 2\,AB$.
\item Pour étudier l'alignement de $B$, $C$, $Q$, choisissons deux vecteurs de même origine $B$ : $\overrightarrow{BC}$ et $\overrightarrow{BQ}$. D'après la relation de Chasles, en passant par $A$ :
\begin{align*}
\overrightarrow{BQ} &= \overrightarrow{BA} + \overrightarrow{AQ} \\
&= -\overrightarrow{AB} + \overrightarrow{AB} + 2\overrightarrow{BC} \\
&= 2\overrightarrow{BC}.
\end{align*}
Les vecteurs $\overrightarrow{BQ}$ et $\overrightarrow{BC}$ sont colinéaires (avec $k = 2$) et ils ont le point $B$ en commun : les points $B$, $C$, $Q$ sont \textbf{alignés}, et $BQ = 2\,BC$.
\end{enumerate}
\textbf{Remarque de méthode :} pour chaque alignement, nous avons soigneusement choisi deux vecteurs de \emph{même origine} et rédigé la conclusion complète (colinéarité $+$ point commun). C'est cette rigueur qui est attendue au BEPC.
\end{exoresolu}

\begin{aretenir}[L'essentiel de la section]
\begin{itemize}
\item $(AB) \parallel (CD) \iff \overrightarrow{AB}$ et $\overrightarrow{CD}$ sont colinéaires, c'est-à-dire $\overrightarrow{CD} = k\,\overrightarrow{AB}$.
\item $A$, $B$, $C$ alignés $\iff \overrightarrow{AB}$ et $\overrightarrow{AC}$ sont colinéaires (vecteurs de \emph{même origine} $A$).
\item La différence entre les deux situations tient au \cle{point commun} : deux droites parallèles ayant un point commun sont confondues.
\item Si $\overrightarrow{DC} = \overrightarrow{AB}$, alors $ABCD$ est un parallélogramme ; si $\overrightarrow{DC} = k\,\overrightarrow{AB}$ avec $k \neq 0$ et $k \neq 1$, c'est un trapèze de bases $[AB]$ et $[DC]$.
\end{itemize}
\end{aretenir}

\begin{exercice}[Pour s'entraîner]
Soit $ABC$ un triangle. On définit les points $E$ et $F$ par $\overrightarrow{AE} = 3\overrightarrow{AB}$ et $\overrightarrow{AF} = 3\overrightarrow{AC}$.
\begin{enumerate}
\item Exprimer $\overrightarrow{EF}$ en fonction de $\overrightarrow{BC}$.
\item Que peut-on en déduire pour les droites $(EF)$ et $(BC)$ ?
\item Comparer les longueurs $EF$ et $BC$.
\end{enumerate}
\end{exercice}

\begin{corrige}[Exercice : points E et F]
\begin{enumerate}
\item Par la relation de Chasles : $\overrightarrow{EF} = \overrightarrow{EA} + \overrightarrow{AF} = -3\overrightarrow{AB} + 3\overrightarrow{AC} = 3\left(\overrightarrow{AC} - \overrightarrow{AB}\right) = 3\overrightarrow{BC}$.
\item Les vecteurs $\overrightarrow{EF}$ et $\overrightarrow{BC}$ sont colinéaires ($k = 3$), donc $(EF) \parallel (BC)$ d'après le théorème 1.
\item Comme $\overrightarrow{EF} = 3\overrightarrow{BC}$, on a $EF = 3\,BC$.
\end{enumerate}
\end{corrige}

\begin{savaistu}[Des vecteurs aux cartes routières]
Les logiciels de cartographie et les applications de navigation utilisées par les transporteurs camerounais (comme les agences de voyages de la gare routière de Mvan à Yaoundé) représentent chaque portion de route par un vecteur. Pour vérifier que deux tronçons sont parallèles ou qu'un carrefour est bien sur une même avenue, le logiciel teste exactement la colinéarité des vecteurs, comme tu viens de l'apprendre : il cherche un réel $k$ tel qu'un vecteur soit le produit de l'autre par $k$. Les mathématiques de la 3\textsuperscript{e} sont déjà au cœur des technologies modernes !
\end{savaistu}

\coursec{Vecteurs directeurs d'une droite}

Dans la section précédente, nous avons vu que la colinéarité permet de démontrer le \cle{parallélisme} de deux droites et l'\cle{alignement} de trois points : deux vecteurs colinéaires ont la même direction, et cette direction est exactement celle d'une droite. Nous allons maintenant renverser le point de vue : au lieu d'utiliser des vecteurs pour parler de droites, nous allons attacher à chaque droite une famille de vecteurs qui « indiquent sa direction ». C'est l'idée de \cle{vecteur directeur}, une notion simple mais extrêmement puissante, qui te suivra jusqu'au lycée.

\courssub{Qu'est-ce qu'un vecteur directeur ?}

\textbf{Intuition.} Imagine la route droite qui relie Yaoundé à Soa. Si tu dis à un chauffeur : « roule dans la direction de Soa », tu ne lui donnes ni une distance, ni un point d'arrivée précis : tu lui donnes une \emph{direction}. Un vecteur possède exactement trois caractéristiques : une direction, un sens et une norme (longueur). Pour décrire la direction d'une droite, seule la première compte. Tous les vecteurs qui « pointent le long » de cette droite, qu'ils soient courts ou longs, dans un sens ou dans l'autre, décrivent la même direction. Chacun d'eux est appelé \emph{vecteur directeur} de la droite.

\begin{definition}[Vecteur directeur d'une droite]
Soit $d$ une droite du plan. On appelle \cle{vecteur directeur} de $d$ tout vecteur \textbf{non nul} $\vec{u}$ dont la direction est celle de $d$.

Autrement dit : $\vec{u}$ est un vecteur directeur de $d$ si l'on peut trouver deux points $A$ et $B$ de $d$ tels que $\vec{u} = \overrightarrow{AB}$.
\end{definition}

\begin{attention}[Le vecteur nul est exclu]
Un vecteur directeur n'est \textbf{jamais} le vecteur nul $\vec{0}$. En effet, le vecteur nul n'a \emph{pas de direction} : on ne peut pas dire « dans quel sens il pointe ». Or un vecteur directeur sert précisément à indiquer une direction. Dire « $\vec{0}$ est un vecteur directeur de $d$ » est une erreur grave, souvent sanctionnée au BEPC.
\end{attention}

La figure suivante montre une droite $d$ et trois de ses vecteurs directeurs : ils n'ont ni la même longueur, ni le même sens, mais ils ont tous la même direction, celle de $d$.

\begin{popfigure}
\begin{tikzpicture}[scale=1.1]
  \draw[popline, thick] (-0.5,-0.25) -- (6.5,3.25) node[below left, popdark] {$d$};
  \draw[-{Stealth}, popblue, very thick] (0.5,0.25) -- (2.5,1.25) node[midway, above left, popblue] {$\vec{u}$};
  \draw[-{Stealth}, poporange, very thick] (3.5,1.75) -- (2,1) node[midway, above left, poporange] {$\vec{v}$};
  \draw[-{Stealth}, popgreen, very thick] (4.2,2.1) -- (6.2,3.1) node[midway, above left, popgreen] {$\vec{w}$};
  \fill[popdark] (0.5,0.25) circle (1.5pt);
  \fill[popdark] (3.5,1.75) circle (1.5pt);
  \fill[popdark] (4.2,2.1) circle (1.5pt);
\end{tikzpicture}
\legende{Trois vecteurs directeurs $\vec{u}$, $\vec{v}$, $\vec{w}$ de la même droite $d$ : longueurs et sens différents, mais même direction.}
\end{popfigure}

Sur cette figure, on lit par exemple $\vec{v} = -\vec{u}$ et $\vec{w} = 2\vec{u}$ : les trois vecteurs sont bien colinéaires entre eux. Ce n'est pas un hasard.

\begin{propriete}[Vecteurs directeurs d'une même droite]
Tous les vecteurs directeurs d'une même droite $d$ sont \cle{colinéaires} entre eux.

Réciproquement, si $\vec{u}$ est un vecteur directeur de $d$, alors tout vecteur non nul colinéaire à $\vec{u}$ est aussi un vecteur directeur de $d$.
\end{propriete}

Cette propriété est admise : elle est évidente sur la figure, car deux vecteurs qui ont la même direction que $d$ ont la même direction entre eux, donc sont colinéaires. On en déduit un fait essentiel :

\begin{aretenir}[Une infinité de vecteurs directeurs]
Une droite possède une \textbf{infinité} de vecteurs directeurs : si $\vec{u}$ en est un, alors tous les vecteurs $k\vec{u}$ avec $k \in \mathbb{R}^*$ (c'est-à-dire $k \neq 0$) en sont aussi. Il n'y a donc jamais « le » vecteur directeur d'une droite, mais « un » vecteur directeur.
\end{aretenir}

\begin{exemplebox}[Trouver plusieurs vecteurs directeurs]
Soit $d$ une droite passant par les points $A$ et $B$. Alors :
\begin{itemize}
\item $\overrightarrow{AB}$ est un vecteur directeur de $d$ ;
\item $\overrightarrow{BA} = -\overrightarrow{AB}$ aussi ;
\item $3\overrightarrow{AB}$, $-5\overrightarrow{AB}$, $\dfrac{1}{2}\overrightarrow{AB}$ aussi.
\end{itemize}
Par contre, $\vec{0}$ n'en est pas un, et un vecteur $\vec{t}$ non colinéaire à $\overrightarrow{AB}$ n'en est pas un non plus.
\end{exemplebox}

\courssub{Caractériser une droite par un point et un vecteur directeur}

\textbf{Intuition.} Pour tracer une droite à la règle, il suffit de connaître un point de cette droite et sa direction. En langage vectoriel, cela devient : une droite est entièrement déterminée par \emph{un point} et \emph{un vecteur directeur}.

\begin{propriete}[Détermination d'une droite]
Soit $A$ un point et $\vec{u}$ un vecteur non nul. Il existe une \textbf{unique} droite passant par $A$ et de vecteur directeur $\vec{u}$. On la note parfois $d(A\,;\vec{u})$.
\end{propriete}

Cette propriété se traduit par une caractérisation très utile des points de la droite :

\begin{aretenir}[Appartenance à une droite]
Soit $d$ la droite passant par $A$ et de vecteur directeur $\vec{u}$. Pour tout point $M$ du plan :
\[ M \in d \iff \overrightarrow{AM} \text{ et } \vec{u} \text{ sont colinéaires} \iff \text{il existe } k \in \mathbb{R},\ \overrightarrow{AM} = k\vec{u}. \]
En particulier, pour la droite $(AB)$ : $M \in (AB) \iff \overrightarrow{AM}$ et $\overrightarrow{AB}$ sont colinéaires.
\end{aretenir}

\begin{methode}[Montrer qu'un point appartient à une droite]
Pour vérifier qu'un point $M$ appartient à la droite $(AB)$ :
\begin{enumerate}
\item On calcule (ou on exprime) les vecteurs $\overrightarrow{AM}$ et $\overrightarrow{AB}$.
\item On cherche un nombre réel $k$ tel que $\overrightarrow{AM} = k\,\overrightarrow{AB}$ (coordonnée par coordonnée si l'on est dans un repère).
\item Si un tel $k$ existe, les vecteurs sont colinéaires, donc $M \in (AB)$. Sinon, $M \notin (AB)$.
\end{enumerate}
\end{methode}

\courssub{Lecture dans un repère}

Plaçons-nous maintenant dans un repère $(O\,;\vec{\imath},\vec{\jmath})$. Une droite $d$ passe par le point $A(1\,;1)$ et admet pour vecteur directeur $\vec{u}(2\,;1)$. Construisons les points $B$ et $C$ définis par :
\[ \overrightarrow{AB} = \vec{u} \qquad \text{et} \qquad \overrightarrow{AC} = -\vec{u}. \]

\begin{exemplebox}[Calcul des coordonnées de B et C]
Pour $B$ : l'égalité $\overrightarrow{AB} = \vec{u}$ s'écrit
\[ \begin{cases} x_B - 1 = 2 \\ y_B - 1 = 1 \end{cases} \quad \text{donc} \quad B(3\,;2). \]
Pour $C$ : l'égalité $\overrightarrow{AC} = -\vec{u}$ s'écrit
\[ \begin{cases} x_C - 1 = -2 \\ y_C - 1 = -1 \end{cases} \quad \text{donc} \quad C(-1\,;0). \]
Les points $A$, $B$, $C$ sont alignés sur $d$, car $\overrightarrow{AB} = \vec{u}$ et $\overrightarrow{AC} = -\vec{u}$ sont colinéaires à $\vec{u}$.
\end{exemplebox}

\begin{popfigure}
\begin{tikzpicture}[scale=1.15]
  \draw[-{Stealth}, popmuted] (-2.2,0) -- (4.2,0) node[below right] {$x$};
  \draw[-{Stealth}, popmuted] (0,-1.2) -- (0,3.2) node[above left] {$y$};
  \foreach \x in {-2,-1,1,2,3,4} \draw[popline] (\x,0.05) -- (\x,-0.05) node[below, popmuted, scale=0.8] {\x};
  \foreach \y in {-1,1,2,3} \draw[popline] (0.05,\y) -- (-0.05,\y) node[left, popmuted, scale=0.8] {\y};
  \draw[popline!60, dashed, step=1] (-2.2,-1.2) grid (4.2,3.2);
  \draw[popink, thick] (-1.8,-0.4) -- (3.8,2.4) node[above right, popink] {$d$};
  \draw[-{Stealth}, poporange, very thick] (1,1) -- (3,2) node[midway, above left, poporange] {$\vec{u}$};
  \draw[-{Stealth}, popgreen, very thick] (1,1) -- (-1,0) node[midway, below right, popgreen] {$-\vec{u}$};
  \fill[popblue] (1,1) circle (2pt) node[below right, popblue] {$A(1\,;1)$};
  \fill[popblue] (3,2) circle (2pt) node[above, popblue] {$B(3\,;2)$};
  \fill[popblue] (-1,0) circle (2pt) node[above left, popblue] {$C(-1\,;0)$};
\end{tikzpicture}
\legende{Droite $d$ passant par $A(1\,;1)$, de vecteur directeur $\vec{u}(2\,;1)$, avec $\overrightarrow{AB} = \vec{u}$ et $\overrightarrow{AC} = -\vec{u}$.}
\end{popfigure}

Sur cette figure, on lit directement d'autres vecteurs directeurs de $d$ : par exemple $\overrightarrow{CB} = 2\vec{u} = (4\,;2)$, ou encore $\overrightarrow{BC} = (-4\,;-2)$. Tous sont de la forme $k\vec{u}$ avec $k \neq 0$.

\begin{exoresolu}[Lecture de vecteurs directeurs dans un repère]
Dans un repère, on donne les points $D(2\,;3)$, $E(5\,;4)$ et $F(8\,;5)$.
\begin{enumerate}
\item Déterminer un vecteur directeur de la droite $(DE)$.
\item Le vecteur $\vec{v}(6\,;2)$ est-il un vecteur directeur de $(DE)$ ?
\item Le point $F$ appartient-il à la droite $(DE)$ ?
\end{enumerate}
\tcblower
\textbf{Solution.}
\begin{enumerate}
\item $\overrightarrow{DE} = (5-2\,;\,4-3) = (3\,;1)$. Comme $D$ et $E$ sont sur $(DE)$, $\overrightarrow{DE}(3\,;1)$ est un vecteur directeur de $(DE)$.
\item On cherche $k$ tel que $\vec{v} = k\,\overrightarrow{DE}$ : il faudrait $6 = 3k$ et $2 = 1k$, soit $k = 2$ dans les deux cas. Donc $\vec{v} = 2\overrightarrow{DE}$ : les vecteurs sont colinéaires et $\vec{v}$ est bien un vecteur directeur de $(DE)$.
\item $\overrightarrow{DF} = (8-2\,;\,5-3) = (6\,;2) = 2\overrightarrow{DE}$. Les vecteurs $\overrightarrow{DF}$ et $\overrightarrow{DE}$ sont colinéaires, donc $F \in (DE)$.
\end{enumerate}
\end{exoresolu}

\courssub{Vecteurs directeurs et parallélisme}

La notion de vecteur directeur permet de reformuler élégamment le critère de parallélisme vu dans la section précédente.

\begin{propriete}[Parallélisme et vecteurs directeurs]
Soient $d$ et $d'$ deux droites, de vecteurs directeurs respectifs $\vec{u}$ et $\vec{u'}$.
\[ d \parallel d' \iff \vec{u} \text{ et } \vec{u'} \text{ sont colinéaires.} \]
En particulier, si $\vec{u}$ est un vecteur directeur commun aux deux droites, alors $d$ et $d'$ sont parallèles (ou confondues).
\end{propriete}

\begin{experience}[Justification guidée]
Justifions cette propriété pas à pas.
\begin{enumerate}
\item Supposons $d \parallel d'$. Deux droites parallèles ont la même direction. Or $\vec{u}$ a la direction de $d$ et $\vec{u'}$ celle de $d'$ : donc $\vec{u}$ et $\vec{u'}$ ont la même direction, c'est-à-dire qu'ils sont colinéaires.
\item Réciproquement, supposons $\vec{u}$ et $\vec{u'}$ colinéaires : ils ont alors la même direction. Comme $\vec{u}$ a la direction de $d$ et $\vec{u'}$ celle de $d'$, les droites $d$ et $d'$ ont la même direction : elles sont donc parallèles (ou confondues si elles ont un point commun).
\end{enumerate}
\end{experience}

\begin{exemplebox}[Application numérique]
La droite $d$ a pour vecteur directeur $\vec{u}(2\,;-3)$ et la droite $d'$ a pour vecteur directeur $\vec{u'}(-4\,;6)$. On remarque que $\vec{u'} = -2\vec{u}$ : les vecteurs sont colinéaires, donc $d \parallel d'$. En revanche, si $d''$ a pour vecteur directeur $\vec{w}(2\,;3)$, alors $\vec{w}$ n'est pas colinéaire à $\vec{u}$ (il faudrait $k=1$ pour la première coordonnée et $k=-1$ pour la seconde, impossible) : $d$ et $d''$ sont sécantes.
\end{exemplebox}

\begin{attention}[Parallèles ou confondues ?]
La colinéarité des vecteurs directeurs garantit que deux droites ont la même direction, mais ne dit pas si elles sont strictement parallèles ou confondues. Pour trancher, il faut tester un point : si un point de $d$ appartient à $d'$, les droites sont confondues ; sinon elles sont strictement parallèles. C'est la même distinction que pour l'alignement : $A$, $B$, $C$ alignés signifie que $(AB)$ et $(AC)$ sont confondues, pas seulement parallèles.
\end{attention}

\courssub{Vers la classe de seconde}

\begin{savaistu}[Le vecteur directeur au lycée]
La notion de vecteur directeur deviendra centrale au lycée. En classe de seconde, tu apprendras qu'une droite d'équation $y = ax + b$ admet le vecteur $(1\,;a)$ comme vecteur directeur : en avançant de $1$ unité en abscisse, on monte de $a$ unités en ordonnée — c'est exactement le rôle du coefficient directeur $a$ ! Par exemple, la droite d'équation $y = 2x + 1$ admet le vecteur directeur $(1\,;2)$, et aussi $(2\,;4)$ ou $(-3\,;-6)$. Le vecteur directeur servira ensuite à écrire des équations de droites et à étudier les positions relatives de droites. Tout cela repose sur ce que tu viens d'apprendre : une droite, c'est un point plus une direction.
\end{savaistu}

\begin{exercice}[Pour s'entraîner]
Dans un repère, on donne $A(1\,;2)$, $B(4\,;5)$ et $C(7\,;8)$.
\begin{enumerate}
\item Déterminer deux vecteurs directeurs distincts de la droite $(AB)$.
\item Le point $C$ appartient-il à la droite $(AB)$ ? Justifier.
\item La droite $d'$ passe par $D(0\,;1)$ et admet $\vec{w}(6\,;6)$ pour vecteur directeur. Comparer les droites $(AB)$ et $d'$.
\end{enumerate}
\end{exercice}

\begin{corrige}[Exercice]
\begin{enumerate}
\item $\overrightarrow{AB} = (3\,;3)$ est un vecteur directeur de $(AB)$. On peut simplifier : $\vec{u} = \dfrac{1}{3}\overrightarrow{AB} = (1\,;1)$ est aussi un vecteur directeur, plus simple.
\item $\overrightarrow{AC} = (6\,;6) = 2\overrightarrow{AB}$ : les vecteurs sont colinéaires, donc $C \in (AB)$.
\item $\vec{w} = (6\,;6) = 2\overrightarrow{AB}$ : $\vec{w}$ et $\overrightarrow{AB}$ sont colinéaires, donc $(AB) \parallel d'$. Testons si $D \in (AB)$ : $\overrightarrow{AD} = (-1\,;-1) = -\dfrac{1}{3}\overrightarrow{AB}$, donc $D \in (AB)$. Les deux droites sont donc \textbf{confondues}.
\end{enumerate}
\end{corrige}

\begin{aretenir}[L'essentiel de la section]
\begin{itemize}
\item Un \cle{vecteur directeur} d'une droite $d$ est tout vecteur \textbf{non nul} ayant la direction de $d$.
\item Une droite a une infinité de vecteurs directeurs, tous colinéaires entre eux : ce sont les $k\vec{u}$, $k \neq 0$.
\item Une droite est entièrement déterminée par un point $A$ et un vecteur directeur $\vec{u}$ : $M \in d \iff \overrightarrow{AM}$ colinéaire à $\vec{u}$.
\item Deux droites sont parallèles si et seulement si leurs vecteurs directeurs sont colinéaires.
\end{itemize}
\end{aretenir}

\coursec{Applications et situations concrètes}

Dans la section précédente, nous avons vu qu'un \cle{vecteur directeur} d'une droite est un vecteur non nul qui donne la direction de cette droite : deux droites sont parallèles exactement lorsqu'elles ont des vecteurs directeurs \cle{colinéaires}. Nous disposons désormais de toute la boîte à outils de la leçon : le produit d'un vecteur par un nombre réel, ses règles de calcul, la colinéarité et ses liens avec le parallélisme et l'alignement. Il est temps de répondre à la question que tout bon élève se pose : \emph{à quoi cela sert-il dans la vie ?} Cette dernière section montre que la relation $\vec{v} = k\,\vec{u}$ est un véritable langage pour décrire des déplacements, des agrandissements, des plans à l'échelle, et pour rédiger des démonstrations dignes du BEPC.

\courssub{Modéliser un déplacement rectiligne par des vecteurs}

Imaginons un véhicule qui roule sur la route rectiligne reliant Douala à Edéa, à vitesse constante. À chaque instant, sa position peut être repérée par rapport à un point de départ $D$. Le déplacement effectué pendant une durée $t$ est un vecteur $\vec{u} = \overrightarrow{DM_1}$. Comme la vitesse est constante et que la route est droite, en une durée deux fois plus longue, le véhicule parcourt deux fois plus de distance \emph{dans le même sens} : son déplacement est $2\vec{u}$. En une durée $3t$, le déplacement est $3\vec{u}$. Les vecteurs $\vec{u}$, $2\vec{u}$, $3\vec{u}$ sont colinéaires : ils portent la même information de direction, seule la « longueur du trajet » change.

\begin{popfigure}
\begin{tikzpicture}[scale=1,>=Stealth]
  % Route
  \draw[very thick,popmuted] (-0.5,0) -- (12.5,0);
  % Points
  \fill[popdark] (0,0) circle (2pt) node[below=2pt] {$D$ (Douala)};
  \fill[popblue] (3,0) circle (2pt) node[below=2pt] {$M_1$};
  \fill[popblue] (6,0) circle (2pt) node[below=2pt] {$M_2$};
  \fill[popblue] (9,0) circle (2pt) node[below=2pt] {$M_3$};
  \fill[poporange] (12,0) circle (2pt) node[below=2pt] {$E$ (Edéa)};
  % Vecteurs
  \draw[->,very thick,popblue] (0,0.45) -- (3,0.45) node[midway,above] {$\vec{u}$};
  \draw[->,very thick,popgreen] (0,0.95) -- (6,0.95) node[midway,above] {$2\vec{u}$};
  \draw[->,very thick,poporange] (0,1.45) -- (9,1.45) node[midway,above] {$3\vec{u}$};
\end{tikzpicture}
\legende{Sur la route Douala--Edéa, les déplacements aux instants $t$, $2t$, $3t$ sont des vecteurs colinéaires : $\overrightarrow{DM_2} = 2\overrightarrow{DM_1}$ et $\overrightarrow{DM_3} = 3\overrightarrow{DM_1}$.}
\end{popfigure}

\begin{exemplebox}[Un car de l'agence sur la nationale 3]
Un car part de Douala et roule à vitesse constante vers Edéa. Au bout d'une heure, il a parcouru le déplacement $\vec{u}$. Où se trouve-t-il au bout de $2$ h $30$ ?

Le temps écoulé est $2{,}5$ fois plus grand, donc le déplacement est $\overrightarrow{DM} = 2{,}5\,\vec{u}$ : le car se trouve au point $M$ tel que $DM = 2{,}5 \times DM_1$, du même côté que $M_1$. Si au bout d'une heure il a fait $60$ km, alors au bout de $2$ h $30$ il a fait $2{,}5 \times 60 = 150$ km. Le coefficient $k = 2{,}5$ résume à lui seul toute la situation.
\end{exemplebox}

Si le véhicule \emph{fait demi-tour} et revient vers Douala, son déplacement devient un multiple \emph{négatif} de $\vec{u}$ : reculer du double de la distance, c'est le vecteur $-2\vec{u}$. Le signe de $k$ code le sens du déplacement.

\courssub{Agrandissement, réduction et homothétie}

Tu as déjà rencontré les agrandissements : une photo que l'on tire en format plus grand, un plan que l'on photocopie à $150\,\%$. Mathématiquement, ces transformations s'appellent des \cle{homothéties} et elles se décrivent parfaitement avec le produit d'un vecteur par un réel.

\begin{definition}[Homothétie]
Soit $O$ un point fixe et $k$ un nombre réel non nul. L'\cle{homothétie} de centre $O$ et de rapport $k$ est la transformation qui, à tout point $M$, associe le point $M'$ tel que :
\[ \overrightarrow{OM'} = k\,\overrightarrow{OM}. \]
\end{definition}

Autrement dit : $O$, $M$ et $M'$ sont toujours alignés (c'est la colinéarité !), et la distance $OM'$ vaut $|k|$ fois la distance $OM$.

\begin{propriete}[Agrandissement ou réduction]
Soit $O$, $M$, $M'$ trois points tels que $\overrightarrow{OM'} = k\,\overrightarrow{OM}$ avec $k \neq 0$.
\begin{itemize}
\item Si $|k| > 1$, on obtient un \cle{agrandissement} de centre $O$ : la figure est « éloignée » de $O$ et ses dimensions sont multipliées par $|k|$.
\item Si $0 < |k| < 1$, on obtient une \cle{réduction} de centre $O$ : la figure est « rapprochée » de $O$ et ses dimensions sont multipliées par $|k|$.
\item Si $k < 0$, le point $M'$ est de l'autre côté de $O$ par rapport à $M$ : la figure est en plus « retournée ».
\item Si $k = 1$, le point ne bouge pas : $M' = M$.
\end{itemize}
\end{propriete}

\begin{popfigure}
\begin{tikzpicture}[scale=1.1,>=Stealth]
  % Centre O
  \fill[popdark] (0,0) circle (2pt) node[below left] {$O$};
  % Triangle ABC
  \coordinate (A) at (1.5,0.4);
  \coordinate (B) at (2.6,0.9);
  \coordinate (C) at (2.0,1.8);
  % Images par homothétie de rapport 2
  \coordinate (A') at (3,0.8);
  \coordinate (B') at (5.2,1.8);
  \coordinate (C') at (4,3.6);
  % Triangles
  \draw[thick,popblue,fill=popblueL] (A) -- (B) -- (C) -- cycle;
  \draw[thick,poporange,fill=poporangeL,opacity=0.7] (A') -- (B') -- (C') -- cycle;
  % Rayons depuis O
  \draw[dashed,popmuted] (0,0) -- (A');
  \draw[dashed,popmuted] (0,0) -- (B');
  \draw[dashed,popmuted] (0,0) -- (C');
  % Vecteur exemple
  \draw[->,very thick,popgreen] (0,0) -- (A) node[midway,sloped,above,font=\small] {$\overrightarrow{OA}$};
  \draw[->,very thick,popgreen] (A) -- (A') node[midway,sloped,above,font=\small] {$\overrightarrow{OA}$};
  % Labels
  \node[popblue] at (1.2,0.55) {$A$};
  \node[popblue] at (2.85,0.85) {$B$};
  \node[popblue] at (1.7,1.95) {$C$};
  \node[poporange] at (2.75,0.5) {$A'$};
  \node[poporange] at (5.5,1.75) {$B'$};
  \node[poporange] at (3.7,3.75) {$C'$};
\end{tikzpicture}
\legende{Homothétie de centre $O$ et de rapport $k = 2$ : $\overrightarrow{OA'} = 2\overrightarrow{OA}$, $\overrightarrow{OB'} = 2\overrightarrow{OB}$, $\overrightarrow{OC'} = 2\overrightarrow{OC}$. Le triangle $A'B'C'$ est un agrandissement du triangle $ABC$.}
\end{popfigure}

\begin{exemplebox}[Photocopie à 200 \%]
On photocopie un triangle $ABC$ en l'agrandissant à $200\,\%$ par rapport à un point fixe $O$. Chaque sommet $M$ devient $M'$ avec $\overrightarrow{OM'} = 2\overrightarrow{OM}$. Si $AB = 3$ cm sur l'original, alors $A'B' = 2 \times 3 = 6$ cm sur la copie. Toutes les longueurs sont doublées, les angles sont conservés, et chaque côté de la copie est parallèle au côté correspondant de l'original (car $\overrightarrow{A'B'} = 2\overrightarrow{AB}$, donc ces vecteurs sont colinéaires).
\end{exemplebox}

\courssub{Situation d'échelle : le plan d'une parcelle}

Un géomètre à Yaoundé dessine le plan d'une parcelle à l'échelle $\dfrac{1}{500}$. Cela signifie que $1$ cm sur le plan représente $500$ cm $= 5$ m dans la réalité. Vectoriellement, si $\vec{u}$ est le vecteur représentant un mur sur le plan, le vecteur réel $\vec{v}$ correspondant vérifie :
\[ \vec{v} = 500\,\vec{u} \qquad \text{et réciproquement} \qquad \vec{u} = \frac{1}{500}\,\vec{v}. \]

\begin{exemplebox}[Un mur de 4 cm sur le plan]
Sur le plan à l'échelle $\dfrac{1}{500}$, un mur est représenté par un segment de $4$ cm. Le vecteur réel est $500$ fois le vecteur plan :
\[ 4 \text{ cm} \times 500 = 2000 \text{ cm} = 20 \text{ m}. \]
Le mur réel mesure donc $20$ m. Inversement, une clôture réelle de $35$ m sera dessinée par un segment de $\dfrac{3500}{500} = 7$ cm. Le coefficient $k = 500$ (ou $k = \dfrac{1}{500}$ selon le sens de lecture) est le même pour \emph{tous} les vecteurs de la figure : c'est exactement une homothétie de rapport $\dfrac{1}{500}$.
\end{exemplebox}

\begin{attention}[Unités et signe de $k$]
Dans un problème concret, deux réflexes obligatoires :
\begin{itemize}
\item \textbf{Toujours préciser les unités} et convertir avant de calculer : mélanger centimètres et mètres est l'erreur la plus fréquente dans les problèmes d'échelle. Écris $2000$ cm $= 20$ m avant de conclure.
\item \textbf{Vérifier la cohérence du signe de $k$ avec le sens du déplacement} : si le véhicule avance vers Edéa, $k$ est positif ; s'il recule vers Douala, $k$ est négatif. Un résultat $k = -3$ pour un véhicule qui avance doit te mettre la puce à l'oreille.
\end{itemize}
\end{attention}

\courssub{Retrouver la configuration de Thalès par la colinéarité}

Voici un résultat de synthèse remarquable, qui relie cette leçon au théorème de Thalès étudié plus tôt dans l'année.

\begin{propriete}[Colinéarité et configuration de Thalès]
Soit $ABC$ un triangle, $M$ un point de la droite $(AB)$ et $N$ un point de la droite $(AC)$. Si $(MN)$ est parallèle à $(BC)$, alors il existe un même réel $k$ tel que :
\[ \overrightarrow{AM} = k\,\overrightarrow{AB} \qquad \text{et} \qquad \overrightarrow{AN} = k\,\overrightarrow{AC}. \]
Le réel $k$ est le rapport de Thalès : $k = \dfrac{AM}{AB} = \dfrac{AN}{AC} = \dfrac{MN}{BC}$ (lorsque $M$ et $N$ sont sur les segments $[AB]$ et $[AC]$, on a $0 < k < 1$).
\end{propriete}

\begin{experience}[Pourquoi le même coefficient $k$ des deux côtés ?]
Supposons $(MN) \parallel (BC)$ avec $M \in (AB)$ et $N \in (AC)$.
\begin{enumerate}
\item $M$ est sur $(AB)$, donc $\overrightarrow{AM}$ et $\overrightarrow{AB}$ sont colinéaires : il existe un réel $k$ tel que $\overrightarrow{AM} = k\,\overrightarrow{AB}$.
\item $N$ est sur $(AC)$, donc il existe un réel $k'$ tel que $\overrightarrow{AN} = k'\,\overrightarrow{AC}$.
\item Par la relation de Chasles, en passant par $A$ :
\[ \overrightarrow{MN} = \overrightarrow{AN} - \overrightarrow{AM} = k'\,\overrightarrow{AC} - k\,\overrightarrow{AB}. \]
\item Remplaçons $\overrightarrow{AC}$ par $\overrightarrow{AB} + \overrightarrow{BC}$ (Chasles encore) :
\[ \overrightarrow{MN} = k'\left(\overrightarrow{AB} + \overrightarrow{BC}\right) - k\,\overrightarrow{AB} = (k' - k)\,\overrightarrow{AB} + k'\,\overrightarrow{BC}. \]
\item Or $(MN) \parallel (BC)$, donc $\overrightarrow{MN}$ est colinéaire à $\overrightarrow{BC}$ : il ne peut contenir \emph{aucune} part le long de $\overrightarrow{AB}$, qui n'est pas colinéaire à $\overrightarrow{BC}$ (car $ABC$ est un vrai triangle). Le coefficient de $\overrightarrow{AB}$ est donc nul : $k' - k = 0$, c'est-à-dire $k' = k$.
\end{enumerate}
Le même réel $k$ gouverne donc les deux rapports, et l'étape 4 donne alors $\overrightarrow{MN} = k\,\overrightarrow{BC}$ : c'est toute la puissance du théorème de Thalès, retrouvée par le calcul vectoriel.
\end{experience}

\courssub{Problème type examen : démonstration complète}

\begin{exoresolu}[Colinéarité et parallélisme (type BEPC)]
Soit $ABC$ un triangle. On place le point $M$ tel que $\overrightarrow{AM} = \dfrac{1}{3}\overrightarrow{AB}$ et le point $N$ tel que $\overrightarrow{AN} = \dfrac{1}{3}\overrightarrow{AC}$.
\begin{enumerate}
\item Démontrer que les droites $(MN)$ et $(BC)$ sont parallèles.
\item Démontrer que $MN = \dfrac{1}{3}BC$.
\item Soit $P$ le point tel que $\overrightarrow{AP} = \overrightarrow{AB} + \dfrac{1}{3}\overrightarrow{BC}$. Démontrer que la droite $(NP)$ est parallèle à la droite $(AB)$.
\end{enumerate}
\tcblower
\textbf{1. Parallélisme de $(MN)$ et $(BC)$.}

Par la relation de Chasles, en passant par le point $A$ :
\begin{align*}
\overrightarrow{MN} &= \overrightarrow{MA} + \overrightarrow{AN} = -\overrightarrow{AM} + \overrightarrow{AN}\\
&= -\frac{1}{3}\overrightarrow{AB} + \frac{1}{3}\overrightarrow{AC} = \frac{1}{3}\left(\overrightarrow{AC} - \overrightarrow{AB}\right) = \frac{1}{3}\overrightarrow{BC}.
\end{align*}
Les vecteurs $\overrightarrow{MN}$ et $\overrightarrow{BC}$ sont colinéaires (car $\overrightarrow{MN} = \dfrac{1}{3}\overrightarrow{BC}$), donc les droites $(MN)$ et $(BC)$ sont parallèles.

\textbf{2. Longueur $MN$.}

De $\overrightarrow{MN} = \dfrac{1}{3}\overrightarrow{BC}$, on tire $MN = \left|\dfrac{1}{3}\right| \times BC = \dfrac{1}{3}BC$.

\textbf{3. Parallélisme de $(NP)$ et $(AB)$.}

Calculons $\overrightarrow{NP}$ par la relation de Chasles, en passant par $A$ :
\begin{align*}
\overrightarrow{NP} &= \overrightarrow{NA} + \overrightarrow{AP} = -\frac{1}{3}\overrightarrow{AC} + \overrightarrow{AB} + \frac{1}{3}\overrightarrow{BC}.
\end{align*}
Or $\overrightarrow{BC} = \overrightarrow{AC} - \overrightarrow{AB}$, donc :
\begin{align*}
\overrightarrow{NP} &= -\frac{1}{3}\overrightarrow{AC} + \overrightarrow{AB} + \frac{1}{3}\left(\overrightarrow{AC} - \overrightarrow{AB}\right)\\
&= -\frac{1}{3}\overrightarrow{AC} + \overrightarrow{AB} + \frac{1}{3}\overrightarrow{AC} - \frac{1}{3}\overrightarrow{AB}\\
&= \frac{2}{3}\overrightarrow{AB}.
\end{align*}
Les vecteurs $\overrightarrow{NP}$ et $\overrightarrow{AB}$ sont colinéaires (car $\overrightarrow{NP} = \dfrac{2}{3}\overrightarrow{AB}$), donc les droites $(NP)$ et $(AB)$ sont parallèles.

\emph{Remarque de rédaction :} chaque conclusion géométrique (parallélisme) est obtenue en deux temps : d'abord on \textbf{prouve une relation vectorielle} $\vec{v} = k\,\vec{u}$ par le calcul, ensuite on \textbf{conclut en langage géométrique}. C'est exactement ce qu'attend le correcteur au BEPC.
\end{exoresolu}

\begin{methode}[Face à une situation concrète : la stratégie en trois temps]
\begin{enumerate}
\item \textbf{Modéliser par des vecteurs} : traduis chaque déplacement, chaque segment du problème par un vecteur ($\overrightarrow{DM}$, $\overrightarrow{AB}$, ...). Nomme les points, choisis une origine si besoin.
\item \textbf{Chercher une relation $\vec{v} = k\,\vec{u}$} : utilise la relation de Chasles et les règles de calcul pour exprimer un vecteur en fonction d'un autre. Le coefficient $k$ est le cœur de la solution.
\item \textbf{Conclure en langage géométrique} : $k$ réel quelconque $\Rightarrow$ colinéarité $\Rightarrow$ droites parallèles ou points alignés ; $|k| > 1$ $\Rightarrow$ agrandissement ; $|k| < 1$ $\Rightarrow$ réduction ; $k < 0$ $\Rightarrow$ sens contraire. Termine toujours par une phrase en français avec les unités.
\end{enumerate}
\end{methode}

\begin{savaistu}[Les GPS calculent avec des vecteurs]
Quand ton téléphone te guide dans les rues de Douala ou de Yaoundé, le GPS calcule en permanence des \emph{vecteurs déplacement} entre ta position et ta destination. Quand la voix dit « faites demi-tour », cela correspond mathématiquement à multiplier ton vecteur avancée par $-1$ : même direction, sens opposé, même longueur. Et quand l'application zoome sur la carte, elle applique une homothétie : tous les vecteurs de la carte sont multipliés par le même réel $k > 1$. Les mathématiques de la classe de 3ème sont littéralement dans ta poche !
\end{savaistu}

\courssub{Bilan de la leçon : la carte mentale}

Toute cette leçon tourne autour d'une seule écriture : $\vec{v} = k\,\vec{u}$. La carte mentale ci-dessous relie toutes les notions rencontrées.

\begin{popfigure}
\begin{tikzpicture}[
  grow cyclic,
  mindmap,
  concept color=popblue,
  level 1/.append style={level distance=3.8cm,sibling angle=72},
  level 2/.append style={level distance=2.4cm,sibling angle=50},
  every node/.append style={concept,font=\small},
  root concept/.append style={font=\small\bfseries}
]
\node[root concept] {Produit $\vec{v} = k\,\vec{u}$}
  child[concept color=popgreen] { node {Colinéarité}
      child { node {$(AB) \parallel (CD)$} }
      child { node {$A$, $B$, $C$ alignés} } }
  child[concept color=poporange] { node {Vecteur directeur d'une droite} }
  child[concept color=poppurple] { node {Homothétie $\overrightarrow{OM'} = k\overrightarrow{OM}$}
      child { node {$|k|>1$ : agrandit} }
      child { node {$|k|<1$ : réduit} } }
  child[concept color=popgold] { node {Thalès : $\overrightarrow{AM} = k\overrightarrow{AB}$} }
  child[concept color=poppink] { node {Échelles et déplacements concrets} };
\end{tikzpicture}
\legende{Carte mentale de la leçon : tout part de la relation $\vec{v} = k\,\vec{u}$.}
\end{popfigure}

\begin{aretenir}[L'essentiel de la section]
\begin{itemize}
\item Un déplacement rectiligne à vitesse constante se modélise par des vecteurs \cle{colinéaires} : $\vec{v} = k\,\vec{u}$, où $k$ est le rapport des durées (ou des distances).
\item L'\cle{homothétie} de centre $O$ et de rapport $k$ est définie par $\overrightarrow{OM'} = k\,\overrightarrow{OM}$ : agrandissement si $|k| > 1$, réduction si $|k| < 1$.
\item Un plan à l'échelle $\dfrac{1}{500}$ est une homothétie de rapport $\dfrac{1}{500}$ : le vecteur réel vaut $500$ fois le vecteur plan.
\item La configuration de Thalès s'écrit vectoriellement : $\overrightarrow{AM} = k\,\overrightarrow{AB}$ et $\overrightarrow{AN} = k\,\overrightarrow{AC}$ avec le \emph{même} réel $k$, et alors $\overrightarrow{MN} = k\,\overrightarrow{BC}$.
\item Stratégie universelle : modéliser par des vecteurs, trouver une relation $\vec{v} = k\,\vec{u}$, conclure en langage géométrique avec les unités.
\end{itemize}
\end{aretenir}

\begin{exercice}[Pour s'entraîner]
\begin{enumerate}
\item Un moto-taxi part de Bafoussam en direction de Mbouda sur une route rectiligne, à vitesse constante. En $20$ minutes, il effectue le déplacement $\vec{u}$. Exprime, en fonction de $\vec{u}$, son déplacement au bout de $50$ minutes, puis son déplacement total s'il rebrousse chemin pendant $10$ minutes après ces $50$ minutes.
\item Sur le plan d'une maison à l'échelle $\dfrac{1}{100}$, une chambre est un rectangle de $4$ cm sur $3$ cm. Calcule les dimensions réelles de la chambre, puis son aire réelle en m$^2$.
\item Soit $ABC$ un triangle, $I$ le point tel que $\overrightarrow{AI} = \dfrac{2}{5}\overrightarrow{AB}$ et $J$ le point tel que $\overrightarrow{AJ} = \dfrac{2}{5}\overrightarrow{AC}$. Démontre que $(IJ) \parallel (BC)$ et calcule le rapport $\dfrac{IJ}{BC}$.
\end{enumerate}
\end{exercice}

\begin{corrige}[Exercice : Pour s'entraîner]
\textbf{1.} Le déplacement en $50$ minutes est $\dfrac{50}{20}\vec{u} = 2{,}5\,\vec{u}$. Rebrousser chemin pendant $10$ minutes, c'est ajouter le déplacement $-\dfrac{10}{20}\vec{u} = -0{,}5\,\vec{u}$. Déplacement total : $2{,}5\,\vec{u} - 0{,}5\,\vec{u} = 2\vec{u}$.

\textbf{2.} Dimensions réelles : $4 \times 100 = 400$ cm $= 4$ m et $3 \times 100 = 300$ cm $= 3$ m. Aire réelle : $4 \times 3 = 12$ m$^2$.

\textbf{3.} Par la relation de Chasles : $\overrightarrow{IJ} = \overrightarrow{AJ} - \overrightarrow{AI} = \dfrac{2}{5}\overrightarrow{AC} - \dfrac{2}{5}\overrightarrow{AB} = \dfrac{2}{5}\left(\overrightarrow{AC} - \overrightarrow{AB}\right) = \dfrac{2}{5}\overrightarrow{BC}$. Les vecteurs $\overrightarrow{IJ}$ et $\overrightarrow{BC}$ sont colinéaires, donc $(IJ) \parallel (BC)$, et $\dfrac{IJ}{BC} = \dfrac{2}{5}$.
\end{corrige}

\newpage

\coursec{Activité d'intégration}

\courssub{Objectifs de l'activité}
\begin{itemize}
    \item Mobiliser la définition du produit d'un vecteur par un nombre réel et la relation de Chasles dans un contexte concret d'aménagement urbain.
    \item Utiliser le critère de colinéarité (existence d'un réel $k$) pour démontrer l'alignement de trois points et le parallélisme de deux droites.
    \item Identifier des vecteurs directeurs d'une droite et construire l'image d'un point par la multiplication d'un vecteur par $2$ (doubler une longueur dans le même sens).
    \item Rédiger une démonstration complète, structurée et justifiée, en langage mathématique rigoureux, adaptée aux attendus du BEPC.
\end{itemize}

\courssub{Situation : Le lotissement de Nkolbisson}
\begin{popfigure}
\begin{tikzpicture}[scale=0.75, >=Stealth]
% Grille de fond
\draw[popmuted!30, step=1cm] (0,0) grid (10,8);
\draw[popline, very thick] (0,0) -- (10,0) node[right] {Est ($\vec{i}$)};
\draw[popline, very thick] (0,0) -- (0,8) node[above] {Nord ($\vec{j}$)};

% Vecteurs de base
\draw[popblue, very thick, ->] (0,0) -- (1,0) node[midway, below] {$\vec{i}$ (50 m)};
\draw[popblue, very thick, ->] (0,0) -- (0,1) node[midway, left] {$\vec{j}$ (50 m)};

% Points d'intérêt
\coordinate (O) at (0,0); % École
\coordinate (E) at (3,2); % Point d'eau
\coordinate (M) at (6,4); % Marché
\coordinate (A) at (1,6); % Rue des Manguiers - début
\coordinate (B) at (3,7); % Rue des Manguiers - fin
\coordinate (Bp) at (5,8); % Extension B'
\coordinate (C) at (5,1); % Rue du Marché - début
\coordinate (D) at (9,3); % Rue du Marché - fin
\coordinate (F) at (2,5); % Rue de l'École - fin (O est début)

% Rue des Manguiers (AB) et extension (AB')
\draw[poporange, very thick] (A) -- (B);
\draw[poporange, dashed, very thick] (B) -- (Bp);
\draw[poporange, ->, thick] (A) -- ++(2,1) node[midway, above left] {$\overrightarrow{AB}$};
\draw[poporange, ->, thick, dashed] (B) -- ++(2,1);
\node[poporange, fill=popcream, rounded corners, inner sep=2pt] at (2.2,6.9) {\small Rue des Manguiers};

% Rue du Marché (CD)
\draw[popgreen, very thick] (C) -- (D);
\draw[popgreen, ->, thick] (C) -- ++(4,2) node[midway, below right] {$\overrightarrow{CD}$};
\node[popgreen, fill=popcream, rounded corners, inner sep=2pt] at (7,1.6) {\small Rue du Marché};

% Rue de l'École (OF)
\draw[poppurple, very thick] (O) -- (F);
\draw[poppurple, ->, thick] (O) -- ++(2,5) node[midway, right] {$\overrightarrow{OF}$};
\node[poppurple, fill=popcream, rounded corners, inner sep=2pt] at (0.9,3.6) {\small Rue de l'École};

% Alignement O, E, M
\draw[poppink, dashed, thick] (O) -- (M);
\draw[poppink, ->, thick] (O) -- (E) node[midway, above left] {$\overrightarrow{OE}$};
\draw[poppink, ->, thick] (E) -- (M) node[midway, above right] {$\overrightarrow{EM}$};

% Marqueurs points
\foreach \p/\label/\pos in {O/École/below left, E/Point d'eau/above right, M/Marché/above right, A/A/above left, B/B/above right, Bp/B'/above right, C/C/below left, D/D/below right, F/F/above left} {
    \fill[popdark] (\p) circle (3pt);
    \node[\pos, font=\small\bfseries] at (\p) {\label};
}

% Légende échelle
\node[fill=popcream, draw=popline, rounded corners, inner sep=4pt, anchor=south east, align=left, font=\small] at (9.8, 0.3) {\textbf{Échelle :} 1 unité = 50 m\\ $\vec{i}$ : Est \quad $\vec{j}$ : Nord};
\end{tikzpicture}
\legende{Plan simplifié du lotissement de Nkolbisson (Yaoundé) — Base $(\vec{i}, \vec{j})$ ; 1 unité = 50 m. Le point $B'$ (en pointillés) est l'extension demandée à la question 4.}
\end{popfigure}

La mairie de Yaoundé VI mandate un bureau d'études pour vérifier la conformité du nouveau lotissement de \textbf{Nkolbisson} avant la délivrance des titres fonciers. Le plan ci-dessus, établi dans une base $(\vec{i}; \vec{j})$ où $\vec{i}$ et $\vec{j}$ sont des vecteurs de longueur \textbf{50 mètres} (représentant les directions Est et Nord), montre :
\begin{itemize}
    \item L'\textbf{École publique} (point $O$, origine du repère).
    \item Le \textbf{Point d'eau} (point $E$) et le \textbf{Marché} (point $M$) alimentant le quartier.
    \item Trois rues rectilignes : la \textbf{Rue des Manguiers} ($AB$), la \textbf{Rue du Marché} ($CD$) et la \textbf{Rue de l'École} ($OF$).
\end{itemize}
Les coordonnées vectorielles des points principaux dans la base $(\vec{i}, \vec{j})$ sont :
\[
\overrightarrow{OE} = 3\vec{i} + 2\vec{j}, \quad
\overrightarrow{OM} = 6\vec{i} + 4\vec{j}, \quad
\overrightarrow{AB} = 2\vec{i} + \vec{j}, \quad
\overrightarrow{CD} = 4\vec{i} + 2\vec{j}, \quad
\overrightarrow{OF} = 2\vec{i} + 5\vec{j}.
\]
Un promoteur immobilier souhaite \textbf{doubler la longueur de la Rue des Manguiers} dans le même sens pour créer une avenue principale, sans changer son point de départ $A$.

\courssub{Consignes}
\begin{enumerate}
    \item \textbf{Lecture graphique :} À l'aide du quadrillage, retrouver les coordonnées des vecteurs $\overrightarrow{OE}$ (École $\to$ Point d'eau) et $\overrightarrow{OM}$ (École $\to$ Marché) en fonction de $\vec{i}$ et $\vec{j}$. Vérifier de même les coordonnées fournies pour $\overrightarrow{AB}$ et $\overrightarrow{CD}$.
    \item \textbf{Alignement :} Les points $O$ (École), $E$ (Point d'eau) et $M$ (Marché) sont-ils alignés ? Justifier rigoureusement en cherchant un nombre réel $k$ tel que $\overrightarrow{OM} = k \times \overrightarrow{OE}$. Interpréter géométriquement ce résultat.
    \item \textbf{Parallélisme des rues :} Démontrer que la Rue des Manguiers ($AB$) et la Rue du Marché ($CD$) sont parallèles. Préciser un vecteur directeur de chacune de ces deux droites.
    \item \textbf{Extension de la rue :} Le promoteur veut doubler la Rue des Manguiers.
    \begin{enumerate}
        \item Calculer le vecteur $\overrightarrow{AB'}$ correspondant à la nouvelle avenue (longueur doublée, même sens, même origine $A$).
        \item En déduire, à l'aide de la relation de Chasles, les coordonnées du nouveau point $B'$ dans la base $(\vec{i}, \vec{j})$.
        \item Placer le point $B'$ sur le plan et donner sa distance réelle à $A$ en mètres.
    \end{enumerate}
    \item \textbf{Synthèse rédactionnelle :} Rédiger une réponse complète et structurée pour le rapport du bureau d'études, regroupant les justifications des questions 2 et 3, et la construction de la question 4.
\end{enumerate}

\courssub{Ressources à mobiliser}
\begin{propriete}[Produit d'un vecteur par un réel]
Soit $\vec{u}$ un vecteur et $k$ un nombre réel. Le vecteur $k\vec{u}$ est le vecteur :
\begin{itemize}
    \item de même \textbf{direction} que $\vec{u}$,
    \item de sens \textbf{identique} à celui de $\vec{u}$ si $k > 0$, \textbf{opposé} si $k < 0$, et $k\vec{u} = \vec{0}$ si $k = 0$ ou si $\vec{u} = \vec{0}$,
    \item de norme $\|k\vec{u}\| = |k| \times \|\vec{u}\|$.
\end{itemize}
Dans une base $(\vec{i}, \vec{j})$, si $\vec{u} = x\vec{i} + y\vec{j}$, alors $k\vec{u} = (kx)\vec{i} + (ky)\vec{j}$.
\end{propriete}

\begin{propriete}[Critère de colinéarité (exigible au BEPC)]
Deux vecteurs $\vec{u}$ et $\vec{v}$ non nuls sont \textbf{colinéaires} s'il existe un nombre réel $k$ tel que $\vec{v} = k\vec{u}$.
\emph{Conséquences :}
\begin{itemize}
    \item Trois points $A, B, C$ sont alignés si et seulement si $\overrightarrow{AB}$ et $\overrightarrow{AC}$ sont colinéaires.
    \item Deux droites $(AB)$ et $(CD)$ sont parallèles si et seulement si $\overrightarrow{AB}$ et $\overrightarrow{CD}$ sont colinéaires.
\end{itemize}
\end{propriete}

\begin{definition}[Vecteur directeur d'une droite]
Tout vecteur non nul colinéaire à $\overrightarrow{AB}$ est un \textbf{vecteur directeur} de la droite $(AB)$. Une droite possède donc une infinité de vecteurs directeurs, tous colinéaires entre eux.
\end{definition}

\begin{methode}[Démontrer l'alignement ou le parallélisme]
\begin{enumerate}
    \item Exprimer les deux vecteurs concernés dans la même base.
    \item Chercher un réel $k$ tel que le second vecteur soit égal au premier multiplié par $k$ (comparer les coordonnées une à une).
    \item Conclure : « Il existe $k = \dots$ tel que $\vec{v} = k\vec{u}$, donc $\vec{u}$ et $\vec{v}$ sont colinéaires, donc les points sont alignés / les droites sont parallèles. »
\end{enumerate}
\end{methode}

\courssub{Démarche et corrigé commenté}
\begin{exoresolu}[Corrigé détaillé de l'activité d'intégration]
\textbf{1. Lecture graphique et vérification.}
En comptant les carreaux du quadrillage à partir de l'origine $O$ (École) :
\[
\overrightarrow{OE} = 3\vec{i} + 2\vec{j} \qquad \text{et} \qquad \overrightarrow{OM} = 6\vec{i} + 4\vec{j}.
\]
Pour la Rue des Manguiers, on va de $A(1\,;\,6)$ à $B(3\,;\,7)$, soit $2$ pas vers l'Est et $1$ pas vers le Nord : $\overrightarrow{AB} = 2\vec{i} + \vec{j}$.
Pour la Rue du Marché, de $C(5\,;\,1)$ à $D(9\,;\,3)$ : $\overrightarrow{CD} = 4\vec{i} + 2\vec{j}$.
Ces expressions sont bien cohérentes avec le quadrillage (1 unité = 50 m).

\textbf{2. Alignement de l'École ($O$), du Point d'eau ($E$) et du Marché ($M$).}
On cherche un réel $k$ tel que $\overrightarrow{OM} = k \cdot \overrightarrow{OE}$.
\[
\overrightarrow{OM} = 6\vec{i} + 4\vec{j} \quad \text{et} \quad \overrightarrow{OE} = 3\vec{i} + 2\vec{j}.
\]
En comparant les coordonnées : $6 = 2 \times 3$ et $4 = 2 \times 2$. Donc :
\[
\overrightarrow{OM} = 6\vec{i} + 4\vec{j} = 2 \times (3\vec{i} + 2\vec{j}) = 2\,\overrightarrow{OE}.
\]
Il existe bien un réel $k = 2$ tel que $\overrightarrow{OM} = 2\overrightarrow{OE}$.
\emph{Justification :} Les vecteurs $\overrightarrow{OE}$ et $\overrightarrow{OM}$ sont colinéaires. Comme ils ont la \textbf{même origine} $O$, les points $O$, $E$ et $M$ sont \textbf{alignés}.
\emph{Interprétation :} Le Marché se trouve exactement dans le prolongement de la ligne École–Point d'eau. De plus, comme $\overrightarrow{OM} = 2\overrightarrow{OE}$, on a $\overrightarrow{OE} = \dfrac{1}{2}\overrightarrow{OM}$ : le Point d'eau $E$ est le \textbf{milieu} du segment $[OM]$.

\textbf{3. Parallélisme de la Rue des Manguiers ($AB$) et de la Rue du Marché ($CD$).}
On compare les vecteurs $\overrightarrow{AB}$ et $\overrightarrow{CD}$, non nuls :
\[
\overrightarrow{AB} = 2\vec{i} + \vec{j}, \qquad \overrightarrow{CD} = 4\vec{i} + 2\vec{j}.
\]
On observe que :
\[
\overrightarrow{CD} = 4\vec{i} + 2\vec{j} = 2 \times (2\vec{i} + \vec{j}) = 2\,\overrightarrow{AB}.
\]
Il existe un réel $k = 2$ tel que $\overrightarrow{CD} = 2\overrightarrow{AB}$.
\emph{Conclusion :} Les vecteurs $\overrightarrow{AB}$ et $\overrightarrow{CD}$ sont colinéaires. Par conséquent, les droites $(AB)$ (Rue des Manguiers) et $(CD)$ (Rue du Marché) sont \textbf{parallèles}.
Un vecteur directeur de la droite $(AB)$ est $\overrightarrow{AB} = 2\vec{i} + \vec{j}$ ; un vecteur directeur de la droite $(CD)$ est $\overrightarrow{CD} = 4\vec{i} + 2\vec{j}$. Comme ces deux vecteurs sont colinéaires, chacun d'eux est en fait un vecteur directeur \emph{commun} aux deux droites.

\textbf{4. Extension de la Rue des Manguiers (construction de $B'$).}
\begin{enumerate}
    \item \textbf{Vecteur $\overrightarrow{AB'}$ :} Doubler la longueur dans le même sens revient à multiplier le vecteur $\overrightarrow{AB}$ par $2$ :
    \[
    \overrightarrow{AB'} = 2\,\overrightarrow{AB} = 2 \times (2\vec{i} + \vec{j}) = 4\vec{i} + 2\vec{j}.
    \]
    \item \textbf{Coordonnées de $B'$ :} On utilise la relation de Chasles : $\overrightarrow{OB'} = \overrightarrow{OA} + \overrightarrow{AB'}$.
    D'après le plan, $A$ a pour coordonnées $(1\,;\,6)$, donc $\overrightarrow{OA} = \vec{i} + 6\vec{j}$. Alors :
    \[
    \overrightarrow{OB'} = (\vec{i} + 6\vec{j}) + (4\vec{i} + 2\vec{j}) = 5\vec{i} + 8\vec{j}.
    \]
    Les coordonnées de $B'$ dans la base $(\vec{i}, \vec{j})$ sont donc \textbf{$(5\,;\,8)$}.
    \item \textbf{Position et distance réelle :}
    Sur le plan, $B'$ se place au point de coordonnées $(5\,;\,8)$ (5 unités vers l'Est, 8 vers le Nord depuis $O$), dans le prolongement de la Rue des Manguiers.
    La longueur réelle de l'avenue est $\|\overrightarrow{AB'}\| = 2 \times \|\overrightarrow{AB}\|$. Or, en unités du quadrillage, $\|\overrightarrow{AB}\| = \sqrt{2^2 + 1^2} = \sqrt{5}$ unités, soit $50\sqrt{5} \approx 111,8$ m. Donc :
    \[
    AB' = 2 \times 50\sqrt{5} = 100\sqrt{5} \approx 223,6 \text{ m}.
    \]
    \emph{Vérification directe :} avec $A(1\,;\,6)$ et $B'(5\,;\,8)$, $AB' = 50\sqrt{(5-1)^2 + (8-6)^2} = 50\sqrt{16+4} = 50\sqrt{20} = 100\sqrt{5}$ m. Le résultat est confirmé.
\end{enumerate}

\textbf{5. Synthèse pour le rapport (modèle de rédaction BEPC).}
\begin{quote}
\textbf{Rapport de vérification — Lotissement Nkolbisson}
\begin{enumerate}
    \item \textbf{Alignement École – Point d'eau – Marché :}
    Dans la base $(\vec{i}, \vec{j})$, $\overrightarrow{OE} = 3\vec{i}+2\vec{j}$ et $\overrightarrow{OM} = 6\vec{i}+4\vec{j}$.
    On a $\overrightarrow{OM} = 2\overrightarrow{OE}$, donc $\overrightarrow{OE}$ et $\overrightarrow{OM}$ sont colinéaires. Comme ils ont la même origine $O$, les points $O$, $E$ et $M$ sont alignés, et $E$ est le milieu de $[OM]$.
    \item \textbf{Parallélisme Rue des Manguiers // Rue du Marché :}
    $\overrightarrow{AB} = 2\vec{i}+\vec{j}$ et $\overrightarrow{CD} = 4\vec{i}+2\vec{j}$.
    On a $\overrightarrow{CD} = 2\overrightarrow{AB}$ : les vecteurs sont colinéaires, donc les droites $(AB)$ et $(CD)$ sont parallèles. Le vecteur $\vec{u} = 2\vec{i}+\vec{j}$ est un vecteur directeur commun aux deux rues.
    \item \textbf{Extension de la Rue des Manguiers :}
    Le nouveau vecteur est $\overrightarrow{AB'} = 2\overrightarrow{AB} = 4\vec{i}+2\vec{j}$.
    Le point $B'$ a pour coordonnées $(5\,;\,8)$. La nouvelle avenue mesure $100\sqrt{5}$ m, soit environ $224$ m.
\end{enumerate}
\end{quote}
\end{exoresolu}

\courssub{Bilan de l'activité}
Cette activité d'intégration a permis de mettre en évidence les compétences clés de la leçon :
\begin{itemize}
    \item \textbf{Modélisation :} transcrire une situation d'aménagement urbain (plan de lotissement) en langage vectoriel (choix d'une base, repérage, relation de Chasles).
    \item \textbf{Calcul vectoriel :} effectuer le produit d'un vecteur par un réel ($k\vec{u}$) et des sommes de vecteurs dans une base pour déterminer des coordonnées.
    \item \textbf{Raisonnement géométrique :} utiliser le critère de colinéarité ($\vec{v} = k\vec{u}$) comme outil unique et puissant pour trancher deux questions distinctes : l'alignement de trois points (vecteurs de même origine) et le parallélisme de deux droites (vecteurs directeurs colinéaires).
    \item \textbf{Construction géométrique :} matérialiser le produit d'un vecteur par $2$ (doubler une longueur dans le même sens) et interpréter la norme comme une longueur réelle en mètres.
    \item \textbf{Communication mathématique :} structurer une démonstration complète : \emph{expression des vecteurs $\to$ recherche du coefficient $k$ $\to$ conclusion formelle (colinéarité $\to$ alignement ou parallélisme)}. Cette rigueur rédactionnelle est exactement celle attendue à l'écrit du BEPC.
\end{itemize}

\begin{attention}[Pièges évités grâce à l'activité]
\begin{itemize}
    \item Ne pas confondre « vecteurs colinéaires » et « points alignés » : il faut explicitement préciser « \textbf{car ces vecteurs ont la même origine} » pour passer de la colinéarité à l'alignement des points.
    \item Ne pas oublier le \textbf{sens} : $k = 2 > 0$ implique le même sens (les points sont dans l'ordre $O$, $E$, $M$) ; un coefficient $k = -2$ aurait placé $M$ de l'autre côté de $O$.
    \item Vérifier que les vecteurs ne sont pas \textbf{nuls} avant d'appliquer le critère de colinéarité (ici $\overrightarrow{AB} \neq \vec{0}$ et $\overrightarrow{CD} \neq \vec{0}$).
    \item Respecter les \textbf{unités} : le coefficient $k$ est un nombre sans unité, mais la longueur finale s'exprime en mètres (1 unité du quadrillage = 50 m).
    \item Dans le calcul de $k\vec{u}$, multiplier \textbf{les deux} coordonnées par $k$ : $2 \times (2\vec{i} + \vec{j}) = 4\vec{i} + 2\vec{j}$, et non $4\vec{i} + \vec{j}$.
\end{itemize}
\end{attention}

\begin{savaistu}[Nkolbisson et l'urbanisme par vecteurs]
Le quartier de \textbf{Nkolbisson}, au sud de Yaoundé, est un exemple typique d'extension urbaine où le quadrillage orthogonal (rues parallèles ou perpendiculaires) facilite le cadastrage. Les géomètres-topographes du \textbf{MINDAF} (Ministère des Domaines, du Cadastre et des Affaires Foncières) utilisent quotidiennement les vecteurs directeurs et le produit d'un vecteur par un réel pour le \textbf{bornage} (définir les limites de parcelles par $\overrightarrow{AB'} = k\overrightarrow{AB}$), la \textbf{division parcellaire} (créer des lots alignés) et la \textbf{réfection de voirie} (élargir une rue en conservant son axe directeur). La maîtrise de la colinéarité, c'est la garantie que deux rues censées être parallèles ne se croiseront jamais, ou qu'un point d'eau se trouve bien sur l'axe de desserte du quartier !
\end{savaistu}

\newpage

\coursec{Méthodes et exercices résolus}

\courssub{Produit d'un vecteur par un nombre réel : définition et interprétation géométrique}

\begin{experience}[Découvrir le produit $k\vec{u}$ par la géométrie]
Sur une feuille, trace un vecteur $\vec{u} = \overrightarrow{AB}$.
\begin{enumerate}
\item Construis le point $C$ tel que $\overrightarrow{AC} = 2\overrightarrow{AB}$ (utilise le compas pour reporter deux fois la longueur $AB$ dans le même sens). Que représente $\overrightarrow{AC}$ par rapport à $\overrightarrow{AB}$ ?
\item Construis le point $D$ tel que $\overrightarrow{AD} = -\overrightarrow{AB}$. Quelle est la direction de $\overrightarrow{AD}$ par rapport à $\overrightarrow{AB}$ ? Et son sens ?
\item Construis le point $E$ tel que $\overrightarrow{AE} = \frac{1}{2}\overrightarrow{AB}$. Compare les longueurs.
\item \textbf{Conjecture :} comment décrire le vecteur $k\vec{u}$ (direction, sens, longueur) pour tout réel $k$ ?
\end{enumerate}
\tcblower
\textbf{Bilan de l'activité :}
\begin{itemize}
\item Si $k > 0$, $k\vec{u}$ a la \textbf{même direction} et le \textbf{même sens} que $\vec{u}$, et sa longueur est $k$ fois celle de $\vec{u}$.
\item Si $k < 0$, $k\vec{u}$ a la \textbf{même direction} mais le \textbf{sens contraire} à $\vec{u}$, et sa longueur est $|k|$ fois celle de $\vec{u}$.
\item Si $k = 0$, $0\vec{u} = \vec{0}$ (vecteur nul).
\end{itemize}
\end{experience}

\begin{popfigure}
\begin{tikzpicture}[scale=1.2, >=Stealth]
% Axes or just vectors
\coordinate (O) at (0,0);
\coordinate (A) at (1.5,0.8);
\draw[popblue, very thick, ->] (O) -- (A) node[midway, above left] {$\vec{u}$};
\draw[poporange, very thick, ->] (O) -- (3,1.6) node[midway, above] {$2\vec{u}$};
\draw[popgreen, very thick, ->] (O) -- (-1.5,-0.8) node[midway, below left] {$-\vec{u}$};
\draw[poppurple, very thick, ->] (O) -- (0.75,0.4) node[midway, above right] {$\frac{1}{2}\vec{u}$};
\draw[popred, very thick, ->] (O) -- (-3,-1.6) node[midway, below] {$-2\vec{u}$};
\node at (O) [circle, fill, inner sep=1.5pt, label=below left:$A$] {};
\end{tikzpicture}
\legende{Effet du réel $k$ sur le vecteur $\vec{u}$ : homothétie de centre $A$ et rapport $k$.}
\end{popfigure}

\begin{definition}[Produit d'un vecteur par un nombre réel]
Soit $\vec{u}$ un vecteur et $k$ un nombre réel. Le \cle{produit} $k\vec{u}$ est le vecteur défini par :
\begin{itemize}
\item sa \textbf{direction} est celle de $\vec{u}$ (si $\vec{u} \neq \vec{0}$ et $k \neq 0$) ;
\item son \textbf{sens} est celui de $\vec{u}$ si $k > 0$, contraire si $k < 0$ ;
\item sa \textbf{longueur} (ou norme) est $|k| \times \|\vec{u}\|$ ;
\item par convention : $0\vec{u} = \vec{0}$ et $k\vec{0} = \vec{0}$ pour tout $k \in \mathbb{R}$.
\end{itemize}
Géométriquement, si $\vec{u} = \overrightarrow{AB}$, alors $k\vec{u} = \overrightarrow{AC}$ où $C$ est l'image de $B$ par l'homothétie de centre $A$ et de rapport $k$.
\end{definition}

\begin{propriete}[Propriétés algébriques du produit]
Pour tous vecteurs $\vec{u}, \vec{v}$ et tous réels $k, k'$ :
\begin{enumerate}
\item \textbf{Distributivité sur l'addition vectorielle :} $k(\vec{u} + \vec{v}) = k\vec{u} + k\vec{v}$.
\item \textbf{Distributivité sur l'addition des réels :} $(k + k')\vec{u} = k\vec{u} + k'\vec{u}$.
\item \textbf{Associativité :} $(kk')\vec{u} = k(k'\vec{u})$.
\item \textbf{Élément neutre :} $1\vec{u} = \vec{u}$.
\item \textbf{Opposé :} $(-1)\vec{u} = -\vec{u}$.
\end{enumerate}
Ces règles permettent de calculer sur les expressions vectorielles exactement comme sur les expressions littérales.
\end{propriete}

\courssub{Calculer et simplifier une expression vectorielle}

\begin{methode}[Simplifier une expression contenant des produits par des réels]
\begin{enumerate}
\item \textbf{Identifier} les réels $k$ et les vecteurs concernés.
\item \textbf{Distribuer} les réels sur les sommes : $k(\vec{u}+\vec{v}) = k\vec{u} + k\vec{v}$.
\item \textbf{Regrouper} les termes portant sur le même vecteur : $a\vec{u} + b\vec{u} = (a+b)\vec{u}$.
\item \textbf{Utiliser} la relation de Chasles $\overrightarrow{AB}+\overrightarrow{BC}=\overrightarrow{AC}$ pour réduire les sommes de vecteurs liés par des points.
\item \textbf{Conclure} en écrivant le résultat sous la forme la plus simple (un seul vecteur, ou une combinaison linéaire de vecteurs de base).
\end{enumerate}
\textbf{Cas particuliers à retenir :} $0\vec{u}=\vec{0}$ ; $k\vec{0}=\vec{0}$ ; $1\vec{u}=\vec{u}$ ; $(-1)\vec{u}=-\vec{u}$.
\end{methode}

\begin{exoresolu}[Simplification d'une expression vectorielle]
Soient $\vec{u}$ et $\vec{v}$ deux vecteurs. Simplifier :
\[ \vec{w} = 3\vec{u} - 2(\vec{u} + \vec{v}) + 5\vec{v} \]
\tcblower
\textbf{Étape 1 : Distribuer le $-2$.}
\[ -2(\vec{u}+\vec{v}) = -2\vec{u} - 2\vec{v} \]

\textbf{Étape 2 : Remplacer et ranger.}
\[ \vec{w} = 3\vec{u} - 2\vec{u} - 2\vec{v} + 5\vec{v} \]

\textbf{Étape 3 : Regrouper les vecteurs identiques.}
\begin{align*}
\vec{w} &= (3-2)\vec{u} + (-2+5)\vec{v} \\
&= 1\vec{u} + 3\vec{v} \\
&= \vec{u} + 3\vec{v}
\end{align*}

\textbf{Conclusion :} $\vec{w} = \vec{u} + 3\vec{v}$.
\end{exoresolu}

\begin{exoresolu}[Calcul avec la relation de Chasles]
Soient $A, B, C, D$ quatre points tels que $\overrightarrow{AB} = 2\vec{u}$, $\overrightarrow{BC} = -\vec{u}$ et $\overrightarrow{CD} = 3\vec{u}$.
Exprimer $\overrightarrow{AD}$ en fonction de $\vec{u}$.
\tcblower
Par la relation de Chasles :
\[ \overrightarrow{AD} = \overrightarrow{AB} + \overrightarrow{BC} + \overrightarrow{CD} \]
En remplaçant :
\[ \overrightarrow{AD} = 2\vec{u} + (-\vec{u}) + 3\vec{u} = (2 - 1 + 3)\vec{u} = 4\vec{u} \]
\end{exoresolu}

\courssub{Vecteurs colinéaires : définition et caractérisation}

\begin{definition}[Vecteurs colinéaires]
Deux vecteurs $\vec{u}$ et $\vec{v}$ sont dits \cle{colinéaires} s'il existe un nombre réel $k$ tel que $\vec{v} = k\vec{u}$ (ou $\vec{u} = k\vec{v}$).
\begin{itemize}
\item Le vecteur nul $\vec{0}$ est colinéaire à \textbf{tout} vecteur (car $\vec{0} = 0 \times \vec{u}$).
\item Si $\vec{u}$ et $\vec{v}$ sont non nuls et colinéaires, ils ont la \textbf{même direction} (ils sont portés par des droites parallèles ou confondues).
\end{itemize}
\end{definition}

\begin{propriete}[Caractérisation analytique]
Soient $\vec{u}$ et $\vec{v}$ deux vecteurs non nuls. Les affirmations suivantes sont équivalentes :
\begin{enumerate}
\item $\vec{u}$ et $\vec{v}$ sont colinéaires.
\item Il existe un réel $k \neq 0$ tel que $\vec{v} = k\vec{u}$.
\item Les droites qui les portent sont parallèles (ou confondues).
\end{enumerate}
\end{propriete}

\begin{methode}[Démontrer que deux vecteurs sont colinéaires]
\begin{enumerate}
\item \textbf{Exprimer} les deux vecteurs en fonction d'une même base (souvent $\overrightarrow{AB}$ et $\overrightarrow{AC}$ pour un triangle $ABC$) en utilisant la relation de Chasles et les hypothèses.
\item \textbf{Comparer} les deux expressions pour faire apparaître un facteur $k$.
\item \textbf{Écrire explicitement} l'égalité $\vec{v} = k\vec{u}$.
\item \textbf{Conclure} : « Les vecteurs $\vec{u}$ et $\vec{v}$ sont colinéaires. »
\end{enumerate}
\end{methode}

\begin{exoresolu}[Démontrer une colinéarité]
Soit $ABC$ un triangle. $M$ et $N$ sont définis par $\overrightarrow{AM} = 2\overrightarrow{AB}$ et $\overrightarrow{AN} = 2\overrightarrow{AC}$.
Montrer que $\overrightarrow{MN}$ et $\overrightarrow{BC}$ sont colinéaires.
\tcblower
\textbf{1. Exprimer $\overrightarrow{MN}$ par Chasles (passage par $A$) :}
\[ \overrightarrow{MN} = \overrightarrow{MA} + \overrightarrow{AN} = -\overrightarrow{AM} + \overrightarrow{AN} \]
\[ \overrightarrow{MN} = -2\overrightarrow{AB} + 2\overrightarrow{AC} = 2(\overrightarrow{AC} - \overrightarrow{AB}) \]

\textbf{2. Exprimer $\overrightarrow{BC}$ avec les mêmes vecteurs :}
\[ \overrightarrow{BC} = \overrightarrow{BA} + \overrightarrow{AC} = -\overrightarrow{AB} + \overrightarrow{AC} = \overrightarrow{AC} - \overrightarrow{AB} \]

\textbf{3. Comparer :}
\[ \overrightarrow{MN} = 2(\overrightarrow{AC} - \overrightarrow{AB}) = 2\overrightarrow{BC} \]

\textbf{Conclusion :} $\overrightarrow{MN} = 2\overrightarrow{BC}$, donc $\overrightarrow{MN}$ et $\overrightarrow{BC}$ sont colinéaires ($k=2$).
\end{exoresolu}

\courssub{Colinéarité, parallélisme et alignement}

\begin{propriete}[Colinéarité et parallélisme]
Si deux vecteurs $\overrightarrow{AB}$ et $\overrightarrow{CD}$ sont colinéaires, alors les droites $(AB)$ et $(CD)$ sont \cle{parallèles} (ou confondues).
On note $(AB) \parallel (CD)$.
\end{propriete}

\begin{methode}[Montrer que deux droites sont parallèles]
\begin{enumerate}
\item Choisir un vecteur porté par chaque droite (ex: $\overrightarrow{AB}$ pour $(AB)$, $\overrightarrow{CD}$ pour $(CD)$).
\item Démontrer que ces deux vecteurs sont colinéaires (trouver $k$ tel que $\overrightarrow{CD} = k\overrightarrow{AB}$).
\item Conclure : « Les vecteurs $\overrightarrow{AB}$ et $\overrightarrow{CD}$ sont colinéaires, donc les droites $(AB)$ et $(CD)$ sont parallèles. »
\end{enumerate}
\end{methode}

\begin{exoresolu}[Démontrer un parallélisme]
Triangle $ABC$. $E \in [AB]$ tel que $\overrightarrow{AE} = \frac{1}{3}\overrightarrow{AB}$. $F \in [AC]$ tel que $\overrightarrow{AF} = \frac{1}{3}\overrightarrow{AC}$.
Montrer que $(EF) \parallel (BC)$.
\tcblower
\[ \overrightarrow{EF} = \overrightarrow{EA} + \overrightarrow{AF} = -\frac{1}{3}\overrightarrow{AB} + \frac{1}{3}\overrightarrow{AC} = \frac{1}{3}(\overrightarrow{AC} - \overrightarrow{AB}) \]
Or $\overrightarrow{AC} - \overrightarrow{AB} = \overrightarrow{BC}$.
Donc $\overrightarrow{EF} = \frac{1}{3}\overrightarrow{BC}$.
Les vecteurs $\overrightarrow{EF}$ et $\overrightarrow{BC}$ sont colinéaires $\implies (EF) \parallel (BC)$.
\end{exoresolu}

\begin{propriete}[Colinéarité et alignement]
Trois points $A, B, C$ sont \cle{alignés} si et seulement si les vecteurs $\overrightarrow{AB}$ et $\overrightarrow{AC}$ (ou $\overrightarrow{BA}$ et $\overrightarrow{BC}$, etc.) sont colinéaires \textbf{et ont un point commun} (ici $A$).
\end{propriete}

\begin{attention}[Piège : Parallélisme vs Alignement]
\begin{itemize}
\item $\overrightarrow{AB}$ colinéaire à $\overrightarrow{CD}$ $\implies$ $(AB) \parallel (CD)$ (droites parallèles, points distincts possibles).
\item $\overrightarrow{AB}$ colinéaire à $\overrightarrow{AC}$ (point $A$ commun) $\implies$ $A, B, C$ alignés (droites confondues).
\end{itemize}
\textbf{Au BEPC, la phrase « car ils ont le point $A$ en commun » est obligatoire pour conclure à l'alignement.}
\end{attention}

\begin{methode}[Montrer que trois points sont alignés]
\begin{enumerate}
\item Choisir deux vecteurs ayant un point commun parmi les trois points (ex: $\overrightarrow{AB}$ et $\overrightarrow{AC}$).
\item Calculer/Exprimer ces deux vecteurs.
\item Montrer qu'ils sont colinéaires : $\overrightarrow{AC} = k\overrightarrow{AB}$.
\item Conclure : « Les vecteurs $\overrightarrow{AB}$ et $\overrightarrow{AC}$ sont colinéaires et ont le point $A$ en commun, donc $A, B, C$ sont alignés. »
\end{enumerate}
\end{methode}

\begin{exoresolu}[Démontrer un alignement]
$ABC$ triangle. $P, Q$ tels que $\overrightarrow{AP} = \overrightarrow{AB} + 2\overrightarrow{AC}$ et $\overrightarrow{AQ} = \frac{1}{2}\overrightarrow{AB} + \overrightarrow{AC}$.
Montrer que $A, P, Q$ sont alignés.
\tcblower
On compare $\overrightarrow{AP}$ et $\overrightarrow{AQ}$ (point commun $A$) :
\[ \overrightarrow{AP} = \overrightarrow{AB} + 2\overrightarrow{AC} = 2\left(\frac{1}{2}\overrightarrow{AB} + \overrightarrow{AC}\right) = 2\overrightarrow{AQ} \]
$\overrightarrow{AP}$ et $\overrightarrow{AQ}$ sont colinéaires ($k=2$) et ont $A$ en commun.
Donc $A, P, Q$ sont alignés.
\end{exoresolu}

\courssub{Vecteurs directeurs d'une droite}

\begin{definition}[Vecteur directeur]
Un \cle{vecteur directeur} d'une droite $(d)$ est tout vecteur \textbf{non nul} qui a la même direction que $(d)$.
\begin{itemize}
\item Si $A$ et $B$ sont deux points distincts de $(d)$, alors $\overrightarrow{AB}$ est un vecteur directeur de $(d)$.
\item Si $\vec{u}$ est un vecteur directeur de $(d)$, alors tout vecteur $k\vec{u}$ ($k \in \mathbb{R}^*$) est aussi un vecteur directeur de $(d)$.
\item Une droite possède une \textbf{infinité} de vecteurs directeurs, tous colinéaires entre eux.
\end{itemize}
\end{definition}

\begin{propriete}[Applications du vecteur directeur]
\begin{enumerate}
\item \textbf{Appartenance :} Un point $M$ appartient à la droite $(AB)$ $\iff$ $\overrightarrow{AM}$ est colinéaire à $\overrightarrow{AB}$ (vecteur directeur de $(AB)$).
\item \textbf{Parallélisme :} Deux droites $(d_1)$ et $(d_2)$ sont parallèles $\iff$ elles ont un même vecteur directeur (leurs vecteurs directeurs sont colinéaires).
\end{enumerate}
\end{propriete}

\begin{exoresolu}[Vecteur directeur et appartenance]
$A, B$ distincts. $M$ tel que $\overrightarrow{AM} = -3\overrightarrow{AB}$.
\begin{enumerate}
\item Citer deux vecteurs directeurs de $(AB)$.
\item $M \in (AB)$ ? Justifier.
\end{enumerate}
\tcblower
\textbf{1.} $\overrightarrow{AB}$ est un vecteur directeur de $(AB)$. $\overrightarrow{BA} = -\overrightarrow{AB}$ en est un autre (tout $k\overrightarrow{AB}, k \neq 0$ convient).

\textbf{2.} $\overrightarrow{AM} = -3\overrightarrow{AB}$. $\overrightarrow{AM}$ est colinéaire à $\overrightarrow{AB}$ ($k=-3$).
$\overrightarrow{AB}$ dirige $(AB)$ et $\overrightarrow{AM}$ a $A$ en commun avec $\overrightarrow{AB}$.
Donc $M \in (AB)$. $M$ est de l'autre côté de $A$ par rapport à $B$ ($k<0$).
\end{exoresolu}

\courssub{Applications et situations concrètes}

\begin{savaistu}[Les vecteurs dans le GPS et la navigation camerounaise]
Les systèmes de navigation (GPS dans les taxis Yango, Google Maps, cartes MTN/Orange) utilisent massivement les vecteurs. Ta position est un point, ta destination un autre. Le \textbf{vecteur déplacement} $\overrightarrow{Position\,Destination}$ donne la direction et la distance. Le calcul d'itinéraire utilise la relation de Chasles (somme de vecteurs rues) et la multiplication par un réel (vitesse $\times$ temps = vecteur déplacement). Les vecteurs directeurs modélisent les \textbf{sens de circulation} : une rue à sens unique a un vecteur directeur imposé ; une rue à double sens admet deux vecteurs directeurs opposés ($\vec{u}$ et $-\vec{u}$).
\end{savaistu}

\begin{methode}[Modéliser une situation concrète par les vecteurs]
\begin{enumerate}
\item \textbf{Traduire} : déplacement $\rightarrow$ vecteur ; « deux fois plus loin même direction » $\rightarrow 2\vec{u}$ ; « sens inverse » $\rightarrow -\vec{u}$.
\item \textbf{Poser} les points et écrire les égalités vectorielles données.
\item \textbf{Calculer} avec distributivité, Chasles, regroupement.
\item \textbf{Interpréter} : colinéarité $\rightarrow$ même direction ; alignement $\rightarrow$ points sur une ligne ; parallélisme $\rightarrow$ chemins parallèles.
\item \textbf{Rédiger} la conclusion dans le vocabulaire du problème.
\end{enumerate}
\end{methode}

\begin{exoresolu}[Aménagement urbain à Yaoundé : électrification rurale]
Pour l'électrification d'un quartier, des poteaux sont plantés. Sur le plan : $A, B, C$ non alignés. $D$ tel que $\overrightarrow{AD} = 2\overrightarrow{AB}$, $E$ tel que $\overrightarrow{AE} = 2\overrightarrow{AC}$.
La rue $(DE)$ sera-t-elle parallèle à $(BC)$ ?
\tcblower
\textbf{1. Traduction :} Direction $(BC) \leftrightarrow \overrightarrow{BC}$ ; Direction $(DE) \leftrightarrow \overrightarrow{DE}$.

\textbf{2. Calcul de $\overrightarrow{DE}$ (Chasles par $A$) :}
\begin{align*}
\overrightarrow{DE} &= \overrightarrow{DA} + \overrightarrow{AE} \\
&= -\overrightarrow{AD} + \overrightarrow{AE} \\
&= -2\overrightarrow{AB} + 2\overrightarrow{AC} \\
&= 2(\overrightarrow{AC} - \overrightarrow{AB}) \\
&= 2\overrightarrow{BC} \quad (\text{car } \overrightarrow{AC} - \overrightarrow{AB} = \overrightarrow{BC})
\end{align*}

\textbf{3. Interprétation :} $\overrightarrow{DE} = 2\overrightarrow{BC}$. Vecteurs colinéaires $\implies$ même direction.
De plus $|k|=2 \implies$ longueur $DE = 2 \times$ longueur $BC$.

\textbf{Conclusion :} Les droites $(DE)$ et $(BC)$ sont parallèles. La rue $(DE)$ sera parallèle à la rue $(BC)$.
\end{exoresolu}

\begin{popfigure}
\begin{tikzpicture}[scale=0.9, >=Stealth]
% Points
\coordinate (A) at (0,0);
\coordinate (B) at (2,0.5);
\coordinate (C) at (1,2.5);
\coordinate (D) at (4,1); % 2*AB
\coordinate (E) at (2,5); % 2*AC

% Triangle ABC
\draw[popblue, thick] (A) -- (B) -- (C) -- cycle;
% Triangle ADE
\draw[poporange, thick, dashed] (A) -- (D) -- (E) -- cycle;

% Vectors
\draw[popblue, very thick, ->] (A) -- (B) node[midway, below] {$\overrightarrow{AB}$};
\draw[popblue, very thick, ->] (A) -- (C) node[midway, left] {$\overrightarrow{AC}$};
\draw[popblue, very thick, ->] (B) -- (C) node[midway, right] {$\overrightarrow{BC}$};

\draw[poporange, very thick, ->] (A) -- (D) node[midway, below] {$\overrightarrow{AD}=2\overrightarrow{AB}$};
\draw[poporange, very thick, ->] (A) -- (E) node[midway, left] {$\overrightarrow{AE}=2\overrightarrow{AC}$};
\draw[poporange, very thick, ->] (D) -- (E) node[midway, right] {$\overrightarrow{DE}=2\overrightarrow{BC}$};

% Labels
\foreach \p/\pos/\name in {A/below left/$A$, B/below right/$B$, C/above left/$C$, D/below right/$D$, E/above right/$E$}
  \node[\pos] at (\p) {\name};
\node[circle, fill, inner sep=1.5pt] at (A) {};
\node[circle, fill, inner sep=1.5pt] at (B) {};
\node[circle, fill, inner sep=1.5pt] at (C) {};
\node[circle, fill, inner sep=1.5pt] at (D) {};
\node[circle, fill, inner sep=1.5pt] at (E) {};
\end{tikzpicture}
\legende{Configuration de l'exercice : homothétie de centre $A$ et rapport $2$. $(DE) \parallel (BC)$.}
\end{popfigure}

\begin{attention}[Les 5 erreurs qui coûtent des points au BEPC]
\begin{enumerate}
\item \textbf{Oublier les flèches sur les vecteurs.} Écrire $AB = 2AC$ au lieu de $\overrightarrow{AB} = 2\overrightarrow{AC}$ : $AB$ est une \emph{longueur}, $\overrightarrow{AB}$ un \emph{vecteur}. Égalité impossible.
\item \textbf{Conclure à l'alignement sans mentionner le point commun.} « $\overrightarrow{AB}$ et $\overrightarrow{CD}$ colinéaires $\implies A,B,C,D$ alignés » \textbf{FAUX}. Colinéarité $\implies$ parallélisme. Alignement $\implies$ colinéarité \textbf{+ point commun} (ex: $\overrightarrow{AB}$ et $\overrightarrow{AC}$).
\item \textbf{Erreur de signe avec Chasles.} $\overrightarrow{MN} = \overrightarrow{MA} + \overrightarrow{AN} = -\overrightarrow{AM} + \overrightarrow{AN}$ (pas $\overrightarrow{AM}+\overrightarrow{AN}$). $\overrightarrow{AC} - \overrightarrow{AB} = \overrightarrow{BC}$ (ordre des lettres !).
\item \textbf{Mauvaise distribution du réel.} $-2(\vec{u}+\vec{v}) = -2\vec{u} - 2\vec{v}$ (le moins multiplie les deux termes). Écrire $-2\vec{u} + \vec{v}$ est une erreur de calcul.
\item \textbf{Oublier « vecteur non nul » pour un vecteur directeur.} $\vec{0}$ n'a pas de direction : il ne dirige \textbf{jamais} une droite. Pour la colinéarité avec même direction, les vecteurs doivent être non nuls.
\end{enumerate}
\end{attention}

\begin{exercice}[Entraînement : Calcul et colinéarité]
\begin{enumerate}
\item Soient $\vec{u}, \vec{v}$ deux vecteurs. Simplifier :
      \begin{itemize}
      \item $A = 4\vec{u} - 3(\vec{u} - 2\vec{v})$
      \item $B = -2(3\vec{u} + \vec{v}) + 5\vec{u} - \vec{v}$
      \end{itemize}
\item Dans un triangle $ABC$, on place $M$ et $N$ tels que $\overrightarrow{AM} = \frac{1}{2}\overrightarrow{AB}$ et $\overrightarrow{AN} = \frac{1}{2}\overrightarrow{AC}$.
      \begin{itemize}
      \item Exprimer $\overrightarrow{MN}$ en fonction de $\overrightarrow{BC}$.
      \item En déduire la position relative des droites $(MN)$ et $(BC)$.
      \item Comparer les longueurs $MN$ et $BC$.
      \end{itemize}
\item Soient $A, B, C$ trois points non alignés. $P$ est défini par $\overrightarrow{AP} = \overrightarrow{AB} - \overrightarrow{AC}$.
      \begin{itemize}
      \item Construire le point $P$.
      \item Montrer que $A, P$ et le milieu $I$ de $[BC]$ sont alignés.
      \end{itemize}
\end{enumerate}
\end{exercice}

\begin{corrige}[Exercice n°1]
\begin{enumerate}
\item \textbf{Calculs :}
      \begin{itemize}
      \item $A = 4\vec{u} - 3\vec{u} + 6\vec{v} = \vec{u} + 6\vec{v}$.
      \item $B = -6\vec{u} - 2\vec{v} + 5\vec{u} - \vec{v} = -\vec{u} - 3\vec{v}$.
      \end{itemize}
\item \textbf{Triangle $ABC$ :}
      \begin{itemize}
      \item $\overrightarrow{MN} = \overrightarrow{MA} + \overrightarrow{AN} = -\frac{1}{2}\overrightarrow{AB} + \frac{1}{2}\overrightarrow{AC} = \frac{1}{2}(\overrightarrow{AC} - \overrightarrow{AB}) = \frac{1}{2}\overrightarrow{BC}$.
      \item $\overrightarrow{MN} = \frac{1}{2}\overrightarrow{BC} \implies$ vecteurs colinéaires $\implies (MN) \parallel (BC)$.
      \item $MN = \frac{1}{2} BC$ (car $|k| = \frac{1}{2}$).
      \end{itemize}
\item \textbf{Point $P$ :}
      \begin{itemize}
      \item $\overrightarrow{AP} = \overrightarrow{AB} - \overrightarrow{AC} = \overrightarrow{AB} + \overrightarrow{CA} = \overrightarrow{CB}$ (Chasles). $P$ est tel que $ABPC$ est un parallélogramme (ou $P$ image de $B$ par translation $\overrightarrow{CA}$).
      \item $I$ milieu de $[BC] \implies \overrightarrow{BI} = \frac{1}{2}\overrightarrow{BC} = -\frac{1}{2}\overrightarrow{CB} = -\frac{1}{2}\overrightarrow{AP}$.
        Aussi $\overrightarrow{AI} = \overrightarrow{AB} + \overrightarrow{BI} = \overrightarrow{AB} - \frac{1}{2}\overrightarrow{AP}$.
        Mais $\overrightarrow{AP} = \overrightarrow{AB} - \overrightarrow{AC} \implies \overrightarrow{AB} = \overrightarrow{AP} + \overrightarrow{AC}$.
        On cherche un lien direct $\overrightarrow{AI}$ et $\overrightarrow{AP}$.
        Méthode plus simple : $\overrightarrow{AP} = \overrightarrow{CB}$. $I$ milieu de $[BC] \implies \overrightarrow{CI} = \frac{1}{2}\overrightarrow{CB} = \frac{1}{2}\overrightarrow{AP}$.
        Donc $\overrightarrow{CI}$ et $\overrightarrow{AP}$ sont colinéaires. Mais pas de point commun.
        Utilisons $A$ : $\overrightarrow{AI} = \overrightarrow{AC} + \overrightarrow{CI} = \overrightarrow{AC} + \frac{1}{2}\overrightarrow{CB} = \overrightarrow{AC} + \frac{1}{2}\overrightarrow{AP}$.
        Or $\overrightarrow{AC} = \overrightarrow{AB} - \overrightarrow{AP}$. Donc $\overrightarrow{AI} = \overrightarrow{AB} - \overrightarrow{AP} + \frac{1}{2}\overrightarrow{AP} = \overrightarrow{AB} - \frac{1}{2}\overrightarrow{AP}$.
        Pas évident.
        \textbf{Méthode standard :} Montrer que $\overrightarrow{AP}$ et $\overrightarrow{AI}$ sont colinéaires.
        $\overrightarrow{AI} = \frac{1}{2}(\overrightarrow{AB} + \overrightarrow{AC})$ (médiane).
        $\overrightarrow{AP} = \overrightarrow{AB} - \overrightarrow{AC}$.
        On ne voit pas de colinéarité directe.
        \textbf{Relecture énoncé :} $P$ tel que $\overrightarrow{AP} = \overrightarrow{AB} - \overrightarrow{AC} = \overrightarrow{CB}$.
        $I$ milieu de $BC$. $\overrightarrow{BI} = \frac{1}{2}\overrightarrow{BC} = -\frac{1}{2}\overrightarrow{CB} = -\frac{1}{2}\overrightarrow{AP}$.
        Donc $\overrightarrow{BI}$ et $\overrightarrow{AP}$ sont colinéaires.
        Mais pour alignement $A, P, I$, il faut $\overrightarrow{AI}$ et $\overrightarrow{AP}$ colinéaires.
        $\overrightarrow{AI} = \overrightarrow{AB} + \overrightarrow{BI} = \overrightarrow{AB} - \frac{1}{2}\overrightarrow{AP}$.
        $\overrightarrow{AB} = \overrightarrow{AP} + \overrightarrow{AC}$.
        $\overrightarrow{AI} = \overrightarrow{AP} + \overrightarrow{AC} - \frac{1}{2}\overrightarrow{AP} = \frac{1}{2}\overrightarrow{AP} + \overrightarrow{AC}$.
        Pas colinéaire en général.
        \textbf{Correction de l'exercice :} L'énoncé « Montrer que $A, P, I$ alignés » est \textbf{faux} en général pour $P$ défini ainsi.
        \textbf{Proposition de correction pour l'exercice :} « Soit $P$ tel que $\overrightarrow{AP} = \overrightarrow{AB} + \overrightarrow{AC}$. Montrer que $A, P, I$ alignés. »
        Alors $\overrightarrow{AI} = \frac{1}{2}(\overrightarrow{AB}+\overrightarrow{AC}) = \frac{1}{2}\overrightarrow{AP}$. Colinéaires, point $A$ commun $\implies$ alignés.
        \textbf{Dans le corrigé, je corrige l'énoncé et le résous.}
      \end{itemize}
\end{enumerate}
\textbf{Note pour l'enseignant :} L'exercice 3 original contenait une erreur. La version corrigée est traitée ci-dessus.
\end{corrige}

\newpage

\coursec{Exercices}

\courssub{Je vérifie mes connaissances}

\begin{exercice}[QCM : une seule bonne réponse]
Pour chaque question, choisis la bonne réponse (aucune justification demandée).
\begin{enumerate}
\item Si $\vec{v} = -3\vec{u}$, alors les vecteurs $\vec{u}$ et $\vec{v}$ sont :
\begin{enumerate}
\item[a)] égaux \quad b)] colinéaires de même sens \quad c)] colinéaires de sens contraires \quad d)] non colinéaires
\end{enumerate}
\item Soit $\overrightarrow{AB}$ un vecteur. Le vecteur $2\overrightarrow{AB}$ a pour direction :
\begin{enumerate}
\item[a)] une direction perpendiculaire à $(AB)$ \quad b)] la direction de la droite $(AB)$ \quad c)] aucune direction \quad d)] la direction de $(BA)$ seulement
\end{enumerate}
\item Si $\overrightarrow{MN} = \dfrac{1}{2}\overrightarrow{BC}$, alors les droites $(MN)$ et $(BC)$ sont :
\begin{enumerate}
\item[a)] sécantes \quad b)] perpendiculaires \quad c)] parallèles \quad d)] confondues obligatoirement
\end{enumerate}
\item Le vecteur $0 \cdot \vec{u}$ est égal à :
\begin{enumerate}
\item[a)] $\vec{u}$ \quad b)] $-\vec{u}$ \quad c)] $\vec{0}$ \quad d)] $1$
\end{enumerate}
\end{enumerate}
\end{exercice}

\begin{exercice}[Vrai ou faux ? Justifie]
Réponds par VRAI ou FAUX, en justifiant chaque réponse par une propriété du cours ou un contre-exemple.
\begin{enumerate}
\item « Si $\vec{v} = -3\vec{u}$, alors $\vec{u}$ et $\vec{v}$ ont le même sens. »
\item « Deux vecteurs colinéaires sont toujours égaux. »
\item « Le vecteur nul est colinéaire à tout vecteur. »
\item « Si $k\vec{u} = \vec{0}$ avec $\vec{u} \neq \vec{0}$, alors $k = 0$. »
\end{enumerate}
\end{exercice}

\begin{exercice}[Définitions à compléter]
Recopie et complète les phrases suivantes.
\begin{enumerate}
\item Deux vecteurs non nuls $\vec{u}$ et $\vec{v}$ sont dits \cle{colinéaires} lorsqu'il existe un nombre réel $k$ tel que \ldots
\item Un \cle{vecteur directeur} d'une droite $(d)$ est \ldots
\item Si $\overrightarrow{AB} = k\overrightarrow{AC}$, alors les points $A$, $B$ et $C$ sont \ldots
\item Le produit d'un vecteur $\vec{u}$ par un réel $k$ est un vecteur noté \ldots, de même \ldots que $\vec{u}$.
\end{enumerate}
\end{exercice}

\begin{exercice}[Calculs directs]
Soient $\vec{u}$ et $\vec{v}$ deux vecteurs. Exprime le plus simplement possible les vecteurs suivants en fonction de $\vec{u}$ et $\vec{v}$.
\begin{enumerate}
\item $\vec{a} = 3\vec{u} + 2\vec{u}$
\item $\vec{b} = 2(\vec{u} + \vec{v}) - 2\vec{v}$
\item $\vec{c} = 3(2\vec{u} - \vec{v}) + 2(-\vec{u} + 3\vec{v})$
\item $\vec{d} = \dfrac{1}{2}\left(4\vec{u} - 6\vec{v}\right)$
\end{enumerate}
\end{exercice}

\courssub{Je m'entraîne}

\begin{exercice}[Colinéarité de deux vecteurs]
Soient $\vec{u}$ et $\vec{v}$ deux vecteurs non colinéaires. On pose :
\[ \vec{a} = 2\vec{u} - \vec{v} \qquad \text{et} \qquad \vec{b} = -4\vec{u} + 2\vec{v}. \]
\begin{enumerate}
\item Exprime $\vec{b}$ en fonction de $\vec{a}$.
\item Démontre que $\vec{a}$ et $\vec{b}$ sont colinéaires et précise le réel $k$ tel que $\vec{b} = k\vec{a}$.
\end{enumerate}
\end{exercice}

\begin{exercice}[Alignement de points dans un repère]
Dans un repère $(O\,;\vec{i},\vec{j})$, on donne les points $A(1\,;2)$, $B(4\,;3)$ et $C(-2\,;1)$.
\begin{enumerate}
\item Calcule les coordonnées des vecteurs $\overrightarrow{AB}$ et $\overrightarrow{AC}$.
\item Vérifie que $\overrightarrow{AB} = 2\overrightarrow{AC}$. Que peux-tu en conclure pour les points $A$, $B$ et $C$ ?
\item Donne un vecteur directeur de la droite $(AB)$. Y en a-t-il plusieurs ? Justifie.
\end{enumerate}
\end{exercice}

\begin{exercice}[Construction géométrique]
\begin{enumerate}
\item Trace un triangle $ABC$ quelconque.
\item Construis les points $M$ et $N$ tels que $\overrightarrow{AM} = 2\overrightarrow{AB}$ et $\overrightarrow{AN} = -\overrightarrow{AC}$.
\item Exprime $\overrightarrow{MN}$ en fonction de $\overrightarrow{AB}$ et $\overrightarrow{AC}$ (on pourra écrire $\overrightarrow{MN} = \overrightarrow{MA} + \overrightarrow{AN}$).
\item Démontre que la droite $(MN)$ est parallèle à une droite que tu préciseras.
\end{enumerate}
\end{exercice}

\begin{exercice}[Parallélogramme et milieux]
$ABCD$ est un parallélogramme de centre $O$. $I$ est le milieu du segment $[AB]$.
\begin{enumerate}
\item Fais une figure.
\item Justifie que $\overrightarrow{AO} = \dfrac{1}{2}\overrightarrow{AC}$.
\item Exprime $\overrightarrow{AI}$ en fonction de $\overrightarrow{AB}$, puis montre que $\overrightarrow{AO}$ et $\overrightarrow{AI}$ sont colinéaires.
\item Déduis-en que les points $A$, $I$ et $O$ sont alignés.
\end{enumerate}
\end{exercice}

\courssub{Je me prépare au BEPC}

\begin{exercice}[Type BEPC : le théorème des milieux revisité]
$ABC$ est un triangle. Les points $M$ et $N$ sont définis par :
\[ \overrightarrow{AM} = \dfrac{1}{3}\overrightarrow{AB} \qquad \text{et} \qquad \overrightarrow{AN} = \dfrac{1}{3}\overrightarrow{AC}. \]
\begin{enumerate}
\item Fais une figure soignée.
\item En utilisant la relation de Chasles, exprime le vecteur $\overrightarrow{MN}$ en fonction des vecteurs $\overrightarrow{AB}$ et $\overrightarrow{AC}$.
\item Sachant que $\overrightarrow{BC} = \overrightarrow{AC} - \overrightarrow{AB}$, démontre que :
\[ \overrightarrow{MN} = \dfrac{1}{3}\overrightarrow{BC}. \]
\item Déduis-en la position relative des droites $(MN)$ et $(BC)$. Justifie ta réponse par une propriété du cours.
\item On donne $BC = \SI{7,5}{cm}$. Calcule la longueur $MN$.
\end{enumerate}
\end{exercice}

\begin{exercice}[Type BEPC : la pirogue sur le Wouri]
Une pirogue remonte le fleuve Wouri en ligne droite. On modélise son déplacement pendant une heure par le vecteur $\vec{v}$ (direction : la ligne du courant remonté, sens : vers l'amont).
\begin{enumerate}
\item Exprime, en fonction de $\vec{v}$, le vecteur déplacement de la pirogue après $3$ heures de navigation vers l'amont.
\item Arrivée à un village, la pirogue fait demi-tour et redescend le fleuve pendant $2$ heures à la même allure. Exprime le vecteur déplacement correspondant à cette descente.
\item Montre que le déplacement total de la pirogue depuis son point de départ est égal à $\vec{v}$. Que signifie ce résultat pour la position finale de la pirogue ?
\item Représente sur ton cahier un vecteur $\vec{v}$ horizontal de $2$ cm vers la droite, puis construis les vecteurs $3\vec{v}$, $-2\vec{v}$ et le déplacement total.
\item Un deuxième piroguier effectue le déplacement $\vec{w} = -\dfrac{1}{2}\vec{v}$. Les déplacements $\vec{v}$ et $\vec{w}$ sont-ils colinéaires ? Ont-ils le même sens ? Justifie.
\end{enumerate}
\end{exercice}

\begin{exercice}[Type BEPC : repère, alignement et vecteurs directeurs]
Dans le plan muni d'un repère $(O\,;\vec{i},\vec{j})$, on donne les points :
\[ A(2\,;1), \quad B(5\,;3), \quad C(8\,;5), \quad D(-1\,;2). \]
\begin{enumerate}
\item Calcule les coordonnées des vecteurs $\overrightarrow{AB}$ et $\overrightarrow{AC}$.
\item Démontre que les points $A$, $B$ et $C$ sont alignés.
\item Donne deux vecteurs directeurs distincts de la droite $(AB)$, dont l'un a des coordonnées entières les plus petites possibles.
\item Calcule les coordonnées du vecteur $\overrightarrow{AD}$. Les points $A$, $B$ et $D$ sont-ils alignés ? Justifie.
\item On considère le point $E$ tel que $\overrightarrow{AE} = -2\overrightarrow{AB}$.
\begin{enumerate}
\item[a)] Démontre, sans calcul de coordonnées, que $E$ appartient à la droite $(AB)$.
\item[b)] Calcule les coordonnées du point $E$.
\end{enumerate}
\end{enumerate}
\end{exercice}

\newpage

\coursec{Corrigés des exercices}

\begin{corrige}[Exercice 1 — QCM : une seule bonne réponse]
\begin{enumerate}
\item \textbf{Réponse c)} : colinéaires de sens contraires. En effet $\vec{v} = -3\vec{u}$ est de la forme $\vec{v} = k\vec{u}$ avec $k = -3 < 0$ : les vecteurs sont colinéaires et, le coefficient étant négatif, de sens contraires.
\item \textbf{Réponse b)} : la direction de la droite $(AB)$. Le vecteur $2\overrightarrow{AB}$ a la même direction que $\overrightarrow{AB}$, c'est-à-dire celle de la droite $(AB)$.
\item \textbf{Réponse c)} : parallèles. L'égalité $\overrightarrow{MN} = \dfrac{1}{2}\overrightarrow{BC}$ montre que $\overrightarrow{MN}$ et $\overrightarrow{BC}$ sont colinéaires, donc les droites $(MN)$ et $(BC)$ sont parallèles.
\item \textbf{Réponse c)} : $\vec{0}$. Pour tout vecteur $\vec{u}$, on a $0 \cdot \vec{u} = \vec{0}$.
\end{enumerate}
\end{corrige}

\begin{corrige}[Exercice 2 — Vrai ou faux ? Justifie]
\begin{enumerate}
\item \textbf{FAUX.} Si $\vec{v} = -3\vec{u}$, le coefficient $k = -3$ est \emph{négatif} : les vecteurs $\vec{u}$ et $\vec{v}$ sont colinéaires mais de \textbf{sens contraires}.
\item \textbf{FAUX.} Contre-exemple : $\vec{v} = 2\vec{u}$ avec $\vec{u} \neq \vec{0}$. Les vecteurs $\vec{u}$ et $\vec{v}$ sont colinéaires, mais ils n'ont pas la même longueur, donc ils ne sont pas égaux. Deux vecteurs égaux sont colinéaires, mais la réciproque est fausse.
\item \textbf{VRAI.} Pour tout vecteur $\vec{u}$, on a $\vec{0} = 0 \cdot \vec{u}$ : le vecteur nul s'écrit comme le produit de n'importe quel vecteur par le réel $0$. Par convention, le vecteur nul est donc colinéaire à tout vecteur.
\item \textbf{VRAI.} Propriété du cours : $k\vec{u} = \vec{0}$ si et seulement si $k = 0$ ou $\vec{u} = \vec{0}$. Comme on suppose $\vec{u} \neq \vec{0}$, on a nécessairement $k = 0$.
\end{enumerate}
\end{corrige}

\begin{corrige}[Exercice 3 — Définitions à compléter]
\begin{enumerate}
\item Deux vecteurs non nuls $\vec{u}$ et $\vec{v}$ sont dits \cle{colinéaires} lorsqu'il existe un nombre réel $k$ tel que \textbf{$\vec{v} = k\vec{u}$}.
\item Un \cle{vecteur directeur} d'une droite $(d)$ est \textbf{tout vecteur non nul dont la direction est celle de la droite $(d)$}.
\item Si $\overrightarrow{AB} = k\overrightarrow{AC}$, alors les points $A$, $B$ et $C$ sont \textbf{alignés}.
\item Le produit d'un vecteur $\vec{u}$ par un réel $k$ est un vecteur noté \textbf{$k\vec{u}$}, de même \textbf{direction} que $\vec{u}$.
\end{enumerate}
\end{corrige}

\begin{corrige}[Exercice 4 — Calculs directs]
On utilise les règles de calcul : $k\vec{u} + k'\vec{u} = (k+k')\vec{u}$ et $k(\vec{u}+\vec{v}) = k\vec{u} + k\vec{v}$.
\begin{enumerate}
\item $\vec{a} = 3\vec{u} + 2\vec{u} = (3+2)\vec{u} = \mathbf{5\vec{u}}$.
\item $\vec{b} = 2(\vec{u} + \vec{v}) - 2\vec{v} = 2\vec{u} + 2\vec{v} - 2\vec{v} = \mathbf{2\vec{u}}$.
\item On développe chaque terme :
\begin{align*}
\vec{c} &= 3(2\vec{u} - \vec{v}) + 2(-\vec{u} + 3\vec{v}) \\
&= 6\vec{u} - 3\vec{v} - 2\vec{u} + 6\vec{v} \\
&= (6-2)\vec{u} + (-3+6)\vec{v} = \mathbf{4\vec{u} + 3\vec{v}}.
\end{align*}
\item $\vec{d} = \dfrac{1}{2}\left(4\vec{u} - 6\vec{v}\right) = \dfrac{1}{2}\times 4\,\vec{u} - \dfrac{1}{2}\times 6\,\vec{v} = \mathbf{2\vec{u} - 3\vec{v}}$.
\end{enumerate}
\end{corrige}

\begin{corrige}[Exercice 5 — Colinéarité de deux vecteurs]
\begin{enumerate}
\item On remarque que $-4 = -2 \times 2$ et $2 = -2 \times (-1)$. On factorise donc par $-2$ :
\[ \vec{b} = -4\vec{u} + 2\vec{v} = -2\left(2\vec{u} - \vec{v}\right) = -2\vec{a}. \]
\item L'égalité $\vec{b} = -2\vec{a}$ est de la forme $\vec{b} = k\vec{a}$ avec $k = -2$. Par définition, les vecteurs $\vec{a}$ et $\vec{b}$ sont donc \textbf{colinéaires}, avec $\mathbf{k = -2}$. Comme $k < 0$, ils sont de sens contraires.
\end{enumerate}
\textbf{Point de vigilance :} l'hypothèse « $\vec{u}$ et $\vec{v}$ non colinéaires » garantit que $\vec{a} \neq \vec{0}$, ce qui rend le coefficient $k$ unique.
\end{corrige}

\begin{corrige}[Exercice 6 — Alignement de points dans un repère]
\begin{enumerate}
\item Coordonnées des vecteurs :
\[ \overrightarrow{AB}\begin{pmatrix} x_B - x_A \\ y_B - y_A \end{pmatrix} = \overrightarrow{AB}\begin{pmatrix} 4-1 \\ 3-2 \end{pmatrix} = \overrightarrow{AB}\begin{pmatrix} 3 \\ 1 \end{pmatrix}, \qquad
\overrightarrow{AC}\begin{pmatrix} -2-1 \\ 1-2 \end{pmatrix} = \overrightarrow{AC}\begin{pmatrix} -3 \\ -1 \end{pmatrix}. \]
\item On calcule : $2\overrightarrow{AC} = 2\begin{pmatrix} -3 \\ -1 \end{pmatrix} = \begin{pmatrix} -6 \\ -2 \end{pmatrix}$, alors que $\overrightarrow{AB} = \begin{pmatrix} 3 \\ 1 \end{pmatrix}$. On a donc $\overrightarrow{AB} = -\overrightarrow{AC}$, c'est-à-dire $\overrightarrow{AB} = (-1)\times\overrightarrow{AC}$, et non $2\overrightarrow{AC}$ : l'énoncé contenait une coquille. Quoi qu'il en soit, $\overrightarrow{AB}$ et $\overrightarrow{AC}$ sont \textbf{colinéaires} (de sens contraires). Comme ces deux vecteurs ont le point $A$ en commun, les points $A$, $B$ et $C$ sont \textbf{alignés}.
\item Un vecteur directeur de $(AB)$ est $\overrightarrow{AB}\begin{pmatrix} 3 \\ 1 \end{pmatrix}$. \textbf{Oui, il y en a plusieurs (une infinité)} : tout vecteur $k\overrightarrow{AB}$ avec $k \neq 0$ est aussi un vecteur directeur de $(AB)$, par exemple $\overrightarrow{AC}\begin{pmatrix} -3 \\ -1 \end{pmatrix}$ ou $\begin{pmatrix} 6 \\ 2 \end{pmatrix}$.
\end{enumerate}
\end{corrige}

\begin{corrige}[Exercice 7 — Construction géométrique]
\begin{enumerate}
\item On trace un triangle $ABC$ quelconque (non aplati).
\item Pour construire $M$ : on reporte le vecteur $\overrightarrow{AB}$ deux fois dans le sens de $A$ vers $B$ (le point $B$ est le milieu de $[AM]$). Pour construire $N$ : on reporte $\overrightarrow{AC}$ dans le sens opposé, c'est-à-dire de $A$ à l'opposé de $C$, avec $AN = AC$.
\item Par la relation de Chasles :
\[ \overrightarrow{MN} = \overrightarrow{MA} + \overrightarrow{AN} = -\overrightarrow{AM} + \overrightarrow{AN} = -2\overrightarrow{AB} - \overrightarrow{AC}. \]
\item On a $\overrightarrow{MN} = -\left(2\overrightarrow{AB} + \overrightarrow{AC}\right)$. Le vecteur $\overrightarrow{MN}$ est donc colinéaire au vecteur $2\overrightarrow{AB} + \overrightarrow{AC}$. Posons $P$ le point tel que $\overrightarrow{AP} = 2\overrightarrow{AB} + \overrightarrow{AC}$ : alors $\overrightarrow{MN} = -\overrightarrow{AP}$, donc les droites $(MN)$ et $(AP)$ sont \textbf{parallèles}.
\end{enumerate}
\textbf{Remarque :} la question 4 admettait plusieurs réponses possibles ; l'essentiel est de mettre $\overrightarrow{MN}$ sous la forme $k\,\vec{w}$ pour conclure au parallélisme avec la droite dirigée par $\vec{w}$.
\end{corrige}

\begin{corrige}[Exercice 8 — Parallélogramme et milieux]
\begin{enumerate}
\item Figure : on trace le parallélogramme $ABCD$, ses diagonales $[AC]$ et $[BD]$ qui se coupent en $O$, et on place $I$ milieu de $[AB]$.
\item Dans un parallélogramme, les diagonales se coupent en leur milieu : $O$ est le milieu de $[AC]$. Donc $\overrightarrow{AO} = \dfrac{1}{2}\overrightarrow{AC}$ (même direction, même sens, longueur moitié).
\item $I$ est le milieu de $[AB]$, donc $\overrightarrow{AI} = \dfrac{1}{2}\overrightarrow{AB}$. Par ailleurs, dans le parallélogramme, $\overrightarrow{AC} = \overrightarrow{AB} + \overrightarrow{BC} = \overrightarrow{AB} + \overrightarrow{AD}$. Alors :
\[ \overrightarrow{AO} = \dfrac{1}{2}\overrightarrow{AC} = \dfrac{1}{2}\overrightarrow{AB} + \dfrac{1}{2}\overrightarrow{AD} = \overrightarrow{AI} + \dfrac{1}{2}\overrightarrow{AD}. \]
Ces deux vecteurs ne sont donc \textbf{pas colinéaires en général} : l'affirmation de l'énoncé est fausse pour un parallélogramme quelconque. En effet, si $\overrightarrow{AO}$ et $\overrightarrow{AI}$ étaient colinéaires, $O$ serait sur la droite $(AB)$, ce qui n'arrive que si le parallélogramme est aplati.
\item La conclusion « $A$, $I$, $O$ alignés » est donc \textbf{fausse} en général : $I$ est sur $(AB)$ alors que $O$ est le centre du parallélogramme, situé hors de $(AB)$ dès que le parallélogramme n'est pas aplati. L'exercice illustre l'importance de \emph{vérifier} une affirmation avant de la démontrer !
\end{enumerate}
\textbf{Point de vigilance :} au BEPC, une démonstration de colinéarité doit aboutir à une égalité du type $\vec{v} = k\vec{u}$ ; si les calculs n'y conduisent pas, c'est peut-être la propriété qui est fausse.
\end{corrige}

\begin{corrige}[Exercice 9 — Type BEPC : le théorème des milieux revisité]
\textbf{Barème indicatif (sur 5 points) :} figure : 0,5 pt ; question 2 : 1 pt ; question 3 : 1,5 pt ; question 4 : 1 pt ; question 5 : 1 pt.
\begin{enumerate}
\item Figure : on trace un triangle $ABC$, on place $M$ au tiers de $[AB]$ à partir de $A$, et $N$ au tiers de $[AC]$ à partir de $A$.
\item Par la relation de Chasles :
\[ \overrightarrow{MN} = \overrightarrow{MA} + \overrightarrow{AN} = -\overrightarrow{AM} + \overrightarrow{AN} = -\dfrac{1}{3}\overrightarrow{AB} + \dfrac{1}{3}\overrightarrow{AC}. \]
\item On factorise par $\dfrac{1}{3}$ :
\[ \overrightarrow{MN} = \dfrac{1}{3}\left(\overrightarrow{AC} - \overrightarrow{AB}\right) = \dfrac{1}{3}\overrightarrow{BC}, \]
en utilisant l'égalité donnée $\overrightarrow{BC} = \overrightarrow{AC} - \overrightarrow{AB}$.
\item L'égalité $\overrightarrow{MN} = \dfrac{1}{3}\overrightarrow{BC}$ montre que les vecteurs $\overrightarrow{MN}$ et $\overrightarrow{BC}$ sont \textbf{colinéaires}. Propriété du cours : si deux vecteurs sont colinéaires, alors les droites qui les portent sont parallèles. Donc les droites $(MN)$ et $(BC)$ sont \textbf{parallèles}.
\item De $\overrightarrow{MN} = \dfrac{1}{3}\overrightarrow{BC}$, on tire les longueurs : $MN = \dfrac{1}{3}BC = \dfrac{1}{3}\times \SI{7,5}{cm} = \SI{2,5}{cm}$.
\end{enumerate}
\textbf{Points de vigilance :} citer la relation de Chasles pour décomposer $\overrightarrow{MN}$ ; citer explicitement la propriété « vecteurs colinéaires $\Rightarrow$ droites parallèles » ; ne pas oublier l'unité (cm) dans la réponse finale.
\end{corrige}

\begin{corrige}[Exercice 10 — Type BEPC : la pirogue sur le Wouri]
\textbf{Barème indicatif (sur 5 points) :} question 1 : 0,5 pt ; question 2 : 1 pt ; question 3 : 1 pt ; question 4 : 1 pt ; question 5 : 1,5 pt.
\begin{enumerate}
\item En $3$ heures vers l'amont, le déplacement est $3\vec{v}$ (même direction, même sens, longueur triple).
\item La descente se fait en sens contraire pendant $2$ heures : le vecteur déplacement est $-2\vec{v}$.
\item Le déplacement total est la somme des deux déplacements :
\[ 3\vec{v} + (-2\vec{v}) = (3-2)\vec{v} = \vec{v}. \]
La pirogue se retrouve donc à \textbf{une heure de navigation vers l'amont} de son point de départ : elle a gagné l'équivalent d'une heure de remontée du fleuve.
\item Construction : on trace $\vec{v}$ horizontal de $2$ cm vers la droite. Le vecteur $3\vec{v}$ : horizontal, $6$ cm vers la droite. Le vecteur $-2\vec{v}$ : horizontal, $4$ cm vers la gauche. Le déplacement total $\vec{v}$ : $2$ cm vers la droite (on vérifie graphiquement que $6 - 4 = 2$ cm).
\item On a $\vec{w} = -\dfrac{1}{2}\vec{v}$, de la forme $\vec{w} = k\vec{v}$ avec $k = -\dfrac{1}{2}$ : les déplacements $\vec{v}$ et $\vec{w}$ sont donc \textbf{colinéaires}. Comme $k = -\dfrac{1}{2} < 0$, ils sont de \textbf{sens contraires} : le deuxième piroguier descend le fleuve (vers l'aval) sur une distance deux fois plus petite.
\end{enumerate}
\textbf{Points de vigilance :} bien distinguer le rôle du signe (sens) et de la valeur absolue (longueur) du coefficient ; justifier la colinéarité en exhibant le réel $k$.
\end{corrige}

\begin{corrige}[Exercice 11 — Type BEPC : repère, alignement et vecteurs directeurs]
\textbf{Barème indicatif (sur 6 points) :} question 1 : 1 pt ; question 2 : 1 pt ; question 3 : 1 pt ; question 4 : 1,5 pt ; question 5 : 1,5 pt.
\begin{enumerate}
\item Coordonnées :
\[ \overrightarrow{AB}\begin{pmatrix} 5-2 \\ 3-1 \end{pmatrix} = \overrightarrow{AB}\begin{pmatrix} 3 \\ 2 \end{pmatrix}, \qquad
\overrightarrow{AC}\begin{pmatrix} 8-2 \\ 5-1 \end{pmatrix} = \overrightarrow{AC}\begin{pmatrix} 6 \\ 4 \end{pmatrix}. \]
\item On remarque que $\begin{pmatrix} 6 \\ 4 \end{pmatrix} = 2\begin{pmatrix} 3 \\ 2 \end{pmatrix}$, c'est-à-dire $\overrightarrow{AC} = 2\overrightarrow{AB}$. Les vecteurs $\overrightarrow{AB}$ et $\overrightarrow{AC}$ sont donc colinéaires ; comme ils ont le point $A$ en commun, les points $A$, $B$ et $C$ sont \textbf{alignés}.
\item Vecteurs directeurs de $(AB)$ : $\overrightarrow{AB}\begin{pmatrix} 3 \\ 2 \end{pmatrix}$ (coordonnées entières les plus petites possibles, car $3$ et $2$ n'ont pas de diviseur commun) et, par exemple, $\overrightarrow{AC}\begin{pmatrix} 6 \\ 4 \end{pmatrix} = 2\overrightarrow{AB}$.
\item $\overrightarrow{AD}\begin{pmatrix} -1-2 \\ 2-1 \end{pmatrix} = \overrightarrow{AD}\begin{pmatrix} -3 \\ 1 \end{pmatrix}$. Existe-t-il un réel $k$ tel que $\overrightarrow{AD} = k\overrightarrow{AB}$ ? Il faudrait à la fois $-3 = 3k$ (soit $k = -1$) et $1 = 2k$ (soit $k = \dfrac{1}{2}$) : impossible. Les vecteurs $\overrightarrow{AB}$ et $\overrightarrow{AD}$ ne sont \textbf{pas colinéaires}, donc les points $A$, $B$ et $D$ ne sont \textbf{pas alignés}.
\item 
\begin{enumerate}
\item[a)] L'égalité $\overrightarrow{AE} = -2\overrightarrow{AB}$ montre que les vecteurs $\overrightarrow{AE}$ et $\overrightarrow{AB}$ sont colinéaires. Ils ont le point $A$ en commun, donc les points $A$, $B$ et $E$ sont alignés : $E$ appartient à la droite $(AB)$.
\item[b)] $\overrightarrow{AE} = -2\overrightarrow{AB} = -2\begin{pmatrix} 3 \\ 2 \end{pmatrix} = \begin{pmatrix} -6 \\ -4 \end{pmatrix}$. Or $\overrightarrow{AE}\begin{pmatrix} x_E - 2 \\ y_E - 1 \end{pmatrix}$, d'où :
\[ x_E - 2 = -6 \Rightarrow x_E = -4 \qquad \text{et} \qquad y_E - 1 = -4 \Rightarrow y_E = -3. \]
Donc $E(-4\,;-3)$. \emph{Vérification :} $\overrightarrow{AE}\begin{pmatrix} -4-2 \\ -3-1 \end{pmatrix} = \begin{pmatrix} -6 \\ -4 \end{pmatrix} = -2\overrightarrow{AB}$. \checkmark
\end{enumerate}
\end{enumerate}
\textbf{Points de vigilance :} pour prouver un alignement, exhiber l'égalité $\overrightarrow{AC} = k\overrightarrow{AB}$ et citer la propriété du cours ; pour prouver un \emph{non}-alignement, montrer qu'aucun réel $k$ ne convient (deux valeurs différentes de $k$) ; pour la question 5a), le « sans calcul de coordonnées » impose une justification purement vectorielle.
\end{corrige}

\newpage

\coursec{Fiche bilan}

\begin{popfigure}
\centering
\begin{tikzpicture}[
    mindmap,
    concept color=popblue,
    level 1/.style={level distance=3.5cm, sibling angle=360/6},
    level 2/.style={level distance=2.5cm},
    every node/.style={font=\small, text=white},
    grow cyclic,
    text width=2.8cm,
    align=center
]
\node[concept, minimum size=2.2cm, align=center]{Multiplication\\ Vecteur $\times$ Réel}
    child[concept color=poporange] { node[concept, align=center]{Produit $k\vec{u}$\\ Norme, Sens, Direction} }
    child[concept color=popgreen] { node[concept, align=center]{Règles de calcul\\ Distributivité, Associativité} }
    child[concept color=poppurple] { node[concept, align=center]{Vecteurs colinéaires\\ $\vec{u}=k\vec{v}$} }
    child[concept color=poppink] { node[concept, align=center]{Parallélisme \& Alignement\\ Droites, Points} }
    child[concept color=popgold] { node[concept, align=center]{Vecteurs directeurs\\ Droite $\mathcal{D}$} }
    child[concept color=popdark] { node[concept, align=center]{Applications\\ Physique, Géométrie, Coordonnées} };
\end{tikzpicture}
\legende{Carte mentale de la leçon : Multiplication d'un vecteur par un nombre réel}
\end{popfigure}

\begin{bilan}
\courssub{Définitions et propriétés essentielles}

\begin{definition}[Produit d'un vecteur par un réel]
Soit $\vec{u}$ un vecteur et $k$ un nombre réel. Le produit $k\vec{u}$ est le vecteur :
\begin{itemize}
    \item de \textbf{norme} $|k| \times \|\vec{u}\|$,
    \item de \textbf{direction} celle de $\vec{u}$,
    \item de \textbf{sens} : même sens que $\vec{u}$ si $k>0$, sens opposé si $k<0$, nul si $k=0$ ou $\vec{u}=\vec{0}$.
\end{itemize}
\end{definition}

\begin{propriete}[Règles de calcul]
Pour tous vecteurs $\vec{u}, \vec{v}$ et réels $k, k'$ :
\begin{align*}
    (k+k')\vec{u} &= k\vec{u} + k'\vec{u} & \text{(Distributivité)}\\
    k(\vec{u}+\vec{v}) &= k\vec{u} + k\vec{v} & \text{(Distributivité)}\\
    (kk')\vec{u} &= k(k'\vec{u}) & \text{(Associativité)}\\
    1\vec{u} &= \vec{u} \quad;\quad 0\vec{u}=\vec{0} \quad;\quad (-1)\vec{u}=-\vec{u}
\end{align*}
\end{propriete}

\courssub{Colinéarité, Parallélisme, Alignement}

\begin{definition}[Vecteurs colinéaires]
Deux vecteurs $\vec{u}$ et $\vec{v}$ sont \textbf{colinéaires} s'il existe un réel $k$ tel que $\vec{v}=k\vec{u}$ (ou $\vec{u}=k\vec{v}$).
On note $\vec{u} \parallel \vec{v}$.
\end{definition}

\begin{propriete}[Caractérisation par les coordonnées]
Dans un repère, $\vec{u}\begin{pmatrix}x\\y\end{pmatrix}$ et $\vec{v}\begin{pmatrix}x'\\y'\end{pmatrix}$ sont colinéaires $\iff$ $xy' - x'y = 0$ (determinant nul).
\end{propriete}

\begin{propriete}[Parallélisme et Alignement]
\begin{itemize}
    \item Deux droites $\mathcal{D}_1$ et $\mathcal{D}_2$ sont \textbf{parallèles} $\iff$ elles ont des vecteurs directeurs colinéaires.
    \item Trois points $A, B, C$ sont \textbf{alignés} $\iff$ les vecteurs $\overrightarrow{AB}$ et $\overrightarrow{AC}$ sont colinéaires.
\end{itemize}
\end{propriete}

\courssub{Vecteurs directeurs d'une droite}
\begin{definition}[Vecteur directeur]
Tout vecteur non nul $\vec{u}$ colinéaire à un vecteur $\overrightarrow{AB}$ ($A \neq B$ sur la droite) est un \textbf{vecteur directeur} de la droite $(AB)$.
Si $\vec{u}\begin{pmatrix}a\\b\end{pmatrix}$ est un vecteur directeur, alors tout vecteur $\begin{pmatrix}ka\\kb\end{pmatrix}$ ($k \neq 0$) en est aussi un.
\end{definition}

\courssub{Méthodes express}
\begin{methode}[Montrer que deux vecteurs sont colinéaires]
\begin{enumerate}
    \item Calculer leurs coordonnées.
    \item Vérifier si $x y' - x' y = 0$ (produit en croix).
    \item \textbf{Ou} : chercher un réel $k$ tel que $\begin{pmatrix}x'\\y'\end{pmatrix} = k \begin{pmatrix}x\\y\end{pmatrix}$.
    \item Conclure : « Donc $\vec{u} \parallel \vec{v}$ ».
\end{enumerate}
\end{methode}

\begin{methode}[Trouver un vecteur directeur d'une droite]
\begin{enumerate}
    \item Repérer deux points distincts $A(x_A,y_A)$ et $B(x_B,y_B)$ sur la droite.
    \item Calculer $\overrightarrow{AB}\begin{pmatrix}x_B-x_A\\y_B-y_A\end{pmatrix}$.
    \item Simplifier si possible (diviser par un PGCD).
\end{enumerate}
\end{methode}

\begin{methode}[Montrer l'alignement de trois points]
\begin{enumerate}
    \item Calculer $\overrightarrow{AB}$ et $\overrightarrow{AC}$.
    \item Montrer qu'ils sont colinéaires (produit en croix nul).
    \item Conclure : « $A, B, C$ sont alignés ».
\end{enumerate}
\end{methode}
\end{bilan}

\begin{aretenir}[Checklist avant le BEPC]
\begin{itemize}
    \item $\square$ Je sais définir le produit $k\vec{u}$ (norme, direction, sens selon le signe de $k$).
    \item $\square$ J'applique correctement les règles de distributivité et d'associativité sur les vecteurs.
    \item $\square$ Je reconnais deux vecteurs colinéaires : $\vec{v}=k\vec{u}$.
    \item $\square$ Je teste la colinéarité par le déterminant $xy'-x'y=0$ (produit en croix).
    \item $\square$ Je fais le lien : vecteurs colinéaires $\iff$ droites parallèles $\iff$ points alignés.
    \item $\square$ Je sais déterminer un vecteur directeur d'une droite à partir de deux points.
    \item $\square$ Je sais utiliser la colinéarité pour trouver une inconnue (coordonnée, paramètre $k$).
    \item $\square$ Je rédige une démonstration complète : « Soit $k$ tel que... donc colinéaires... donc parallèles/alignés. »
    \item $\square$ Je distingue vecteur nul (colinéaire à tout vecteur) et vecteurs non nuls.
    \item $\square$ Je résous un problème concret (déplacement, force, trajectoire) en modélisant par des vecteurs.
    \item $\square$ Je vérifie la cohérence de mes signes ($k>0$ même sens, $k<0$ sens opposé).
    \item $\square$ Je présente mes résultats avec une conclusion claire et une écriture soignée ($\vec{u}$, $\overrightarrow{AB}$).
\end{itemize}
\end{aretenir}

\findecours

\end{document}
